\documentclass[10pt,a4paper]{article}

% ----------------- Paquetes -------------------%
\usepackage[T1]{fontenc}
\usepackage[utf8]{inputenc}
\usepackage[a4paper,margin=2.5cm]{geometry}
\usepackage{mathptmx}             
\usepackage{changepage} 
\usepackage{setspace}
\usepackage[spanish]{babel}
\usepackage[labelfont=bf,labelsep=space]{caption}
\usepackage{amssymb}
\usepackage{amsmath}
\usepackage{ebgaramond}
\usepackage{placeins}
\usepackage{etoolbox}
\usepackage{setspace}
\usepackage{parskip} 
\usepackage{tcolorbox}
\usepackage{needspace}
\usepackage{pdflscape}% páginas apaisadas (\imagenfithorizontal)
\usepackage{xparse}
\usepackage{ocgx2}
\usepackage{hyperref}
\usepackage{hypcap}
\usepackage{soul}\sethlcolor{yellow}% resaltado amarillo: \hl{texto} (Panel TCEE, ⌃H)

%%%%%%%%%%%%%%%%%%%%%%%%%%%%%%%%%%%%%%%%%%%%%%%%%%%%%%%%%%%%%%%%
% --- Fija primero tus valores globales "buenos" ---
\setstretch{1.2}                 % Interlineado global
\setlength{\parskip}{0.7\baselineskip}
\setlength{\parindent}{0pt}
\newlength{\GlobalParskip}
\setlength{\GlobalParskip}{\parskip}

\newcounter{eqnum}[subsection]
\renewcommand{\theeqnum}{\arabic{eqnum}}
\newcounter{alphaitem}
\newcounter{numitem}

%% ------------------- Recuadros----------------------%%
\newcommand{\puntosclave}[1]{%
  \begin{tcolorbox}[
    colframe=black,
    title={Puntos clave},
    before upper={\setlength{\parindent}{0.5cm}\setlength{\parskip}{0.5\baselineskip}}
  ]
  #1
  \end{tcolorbox}
}

\newcommand{\modificaciones}[1]{
  \begin{tcolorbox}[
    colback=white,
    colframe=red,
    title={Modificaciones a realizar},
    before upper={\setlength{\parindent}{0.5cm}\setlength{\parskip}{0.5\baselineskip}}
  ]
  #1
  \end{tcolorbox}
}

%% ------------------- Preguntas TEST----------------------%%
% Opciones: radio (single-choice)
\NewDocumentCommand{\OptRadio}{m m m}{%
  \noindent\CheckBox[radio,name=#1,value=#2,width=1.2ex,height=1.2ex]{}%
  \;\textbf{#2.}~#3\par
}

% Opciones: checkbox (multi-choice)
\NewDocumentCommand{\OptCheck}{m m}{%
  \noindent\CheckBox[name=#1,width=1.2ex,height=1.2ex,bordercolor={0 0 0}]{}%
  \;#2\par
}

% Pregunta single-choice (radio)
% Args: 1=id, 2=enunciado, 3=opciones, 4=solución+justificación, 5=imagen (o {}), 6=origen
\NewDocumentCommand{\PreguntaTestSingle}{m m m m m m}{%
  \par\medskip
  \noindent\textbf{#1.}~#2\par\smallskip

    % Imagen (si no es {})
    \IfBlankTF{#5}{}{%
      \begin{center}
        \includegraphics[width=0.75\linewidth]{#5}
      \end{center}
    }

    \begin{Form}
      #3
    \end{Form}

    \par\smallskip
    \switchocg{sol-#1}{\fbox{Mostrar / ocultar solución}}%
    \hfill{\footnotesize #6}\par\smallskip

    \begin{ocg}{Solución}{sol-#1}{0}
      \noindent #4
    \end{ocg}
  \par\medskip
}

% Pregunta multi-choice (checkboxes)
\NewDocumentCommand{\PreguntaTestMulti}{m m m m m m}{%
  \par\medskip
  \noindent\textbf{#1.}~#2\par\smallskip

    % Imagen (si no es {})
    \IfBlankTF{#5}{}{%
      \begin{center}
        \includegraphics[width=0.75\linewidth]{#5}
      \end{center}
    }
    \begin{Form}
      #3
    \end{Form}

    \par\smallskip
    \switchocg{sol-#1}{\fbox{Mostrar / ocultar solución}}%
    \hfill{\footnotesize #6}\par\smallskip

    \begin{ocg}{Solución}{sol-#1}{0}
      \noindent #4
    \end{ocg}
  \par\medskip
}


%% ------------------- CITA ----------------------%%
% \begin{cita}[Autor][Año][Obra] texto \end{cita}: cita sangrada y, debajo a la derecha, «AUTOR (Año), Obra». Los tres datos son opcionales.
% En el Panel TCEE: escribir \cita e Intro (el cursor va al texto; Tab pasa a Autor, Año y Obra).
\NewDocumentEnvironment{cita}{ O{} O{} O{} +b }{%
  \begin{adjustwidth}{2cm}{0pt}
  \raggedright
  #4\par
  \raggedleft
  \if\relax\detokenize{#1#2#3}\relax
    % sin firma
  \else
    #1%
    \if\relax\detokenize{#2}\relax\else\if\relax\detokenize{#1}\relax\else\space\fi(#2)\fi
    \if\relax\detokenize{#3}\relax\else\if\relax\detokenize{#1#2}\relax\else, \fi\textit{#3}\fi
  \fi
  \end{adjustwidth}
}{}



%% ------------------- Listas----------------------%%
    % --- Lista alfabética ---
    \newenvironment{la}
      {\begin{list}{\alph{alphaitem})}{%
          \usecounter{alphaitem}%
          \setlength{\leftmargin}{1cm}%
          \setlength{\parskip}{3pt}% 
          \setlength{\itemsep}{0pt}% ← espacio entre ítems
          \setlength{\parsep}{0pt}%  ← párrafos dentro de un ítem
      }}
      {\end{list}}
      
    % --- Lista numérica ---
    \newenvironment{lnum}
      {\begin{list}{\arabic{numitem}.}{%
          \usecounter{numitem}%
          \setlength{\leftmargin}{1cm}%
          \setlength{\parskip}{3pt}% 
          \setlength{\itemsep}{0pt}% ← espacio entre ítems
          \setlength{\parsep}{0pt}%  ← párrafos dentro de un ítem
      }}
      {\end{list}}
      
% ------------------- Imágenes----------------------%
    \graphicspath{{Img/}}% Rutas por defecto 
    % --- Imagen adaptativa ---
    % \imagenfit[alto]{fichero}{título}{fuente}: la imagen queda fija entre los párrafos donde se escribe (no flota), con su
    % título en el índice de figuras y la fuente debajo. Sin [alto]: las figuras PDF (TikZ) salen al tamaño con que se
    % dibujaron y solo se reducen si no caben en el ancho o en el alto de la página; las PNG/JPG ocupan el ancho del texto
    % hasta \imgaltoras del alto de la página. Con [alto] (fracción del alto de la página; [1] = página entera) se fija el
    % tamaño de esa imagen. Si no cabe en lo que queda de página, pasa entera a la siguiente.
    % \imagenfithorizontal{fichero}{título}{fuente}: la imagen en una página girada (apaisada), en ese mismo punto del texto.
    % Reglas: main/GUIA_TEMAS.md, §4 (Figuras).
    \newdimen\imgcap
    \newdimen\imgh
    \newdimen\imgw
    \newdimen\imgneed
    \newsavebox\imgcaja
    \newcommand{\imgaltoras}{0.30}% alto por defecto de las PNG/JPG (fracción del alto de la página)
    \makeatletter
    \newcommand{\imagenfit}[4][]{%
      \addvspace{10pt}%
      \par\begingroup
        \imgcap=3.5\baselineskip% alto aproximado de título + fuente
        \imgw=\linewidth
        \filename@parse{#2}\def\img@pdf{pdf}%
        \if\relax\detokenize{#1}\relax
          \ifx\filename@ext\img@pdf
            \sbox\imgcaja{\includegraphics{\detokenize{#2}}}% tamaño natural
            \ifdim\wd\imgcaja<\imgw \imgw=\wd\imgcaja \fi
            \imgh=\ht\imgcaja
          \else
            \imgh=\imgaltoras\textheight
          \fi
        \else
          \imgh=#1\textheight
        \fi
        \dimen@=\textheight \advance\dimen@ by -\imgcap
        \ifdim\imgh>\dimen@ \imgh=\dimen@ \fi
        \sbox\imgcaja{\includegraphics[width=\imgw,height=\imgh,keepaspectratio]{\detokenize{#2}}}%
        \imgneed=\ht\imgcaja \advance\imgneed by \imgcap
        \Needspace{\imgneed}%
        \noindent\begin{minipage}{\textwidth}
          \centering
          \captionof{figure}{\textemdash\ #3}\par\vspace{0.3em}% título arriba
          \usebox\imgcaja\par
          \vspace{0.3em}\small\textit{Fuente: }#4% fuente abajo
        \end{minipage}%
      \endgroup
      \par\addvspace{10pt}\RequestRestore
    }
    \newcommand{\imagenfithorizontal}[3]{%
      \begin{landscape}
        \centering
        \captionof{figure}{\textemdash\ #2}\par\vspace{0.3em}%
        \includegraphics[width=\linewidth,height=\dimexpr\textheight-4\baselineskip\relax,keepaspectratio]{\detokenize{#1}}\par
        \vspace{0.3em}{\small\textit{Fuente: }#3}%
      \end{landscape}%
    }
    \makeatother

%------------ Bloque de RECUADRO ESPECIAL -------------%
    \newenvironment{specialblock}[1]{%
      \begin{quote} % crea sangría a izquierda y derecha
      \small % texto más pequeño
      \noindent\textbf{#1}\\[-0.3em] % título en negrita
      \rule{\linewidth}{0.4pt}\\[0.6em]}{
      \end{quote}}
      
%-------------------------- Portada -----------------------%
\newcommand{\oldtitle}[1]{\def\OldTitle{#1}}
\newcommand{\changerationale}[1]{\def\ChangeWhy{#1}}

\newcommand{\changebox}{%
\begin{tcolorbox}[title={Historial de cambios del tema}]
\textbf{Título anterior:} \OldTitle\par
\textbf{Motivo del cambio:} \ChangeWhy
\end{tcolorbox}
}

%-------------------------- Author Font -----------------------%
    \newcommand{\authorfont}[1]{{\fontfamily{EBGaramond-TLF}\selectfont #1}}

%-------------------------- Anexos -----------------------%
    % \anexo{Título}: un anexo equivale a un \section (su título sale del mismo tamaño). Dentro se usan \subsection,
    % \subsubsection y \paragraph, nunca \section: se numeran I.1, I.1.1, I.1.1.1 y las páginas del anexo, I-1, I-2…
    % En el índice de contenido inicial solo figura el nombre del anexo. Sus apartados forman el índice del anexo, que sale
    % solo si el anexo los tiene, justo después de la introducción del anexo (antes del primer \subsection).
    % \anexo*{Título}: el mismo anexo, sin ese índice. Reglas: main/GUIA_TEMAS.md, §6.
    \makeatletter
    \newcounter{anexo}
    \renewcommand{\theanexo}{\Roman{anexo}}
    \newif\ifEnAnexo \EnAnexofalse
    \newif\ifanx@pendiente \anx@pendientefalse
    \newcommand{\anexo}{\@ifstar{\anx@empezar\anx@pendientefalse}{\anx@empezar\anx@pendientetrue}}
    \newcommand{\anx@empezar}[2]{%
      \clearpage
      \refstepcounter{anexo}%
      \global\EnAnexotrue
      \setcounter{page}{1}\gdef\thepage{\theanexo-\arabic{page}}%
      \setcounter{section}{\numexpr\value{anexo}-1\relax}\stepcounter{section}% el anexo hace de \section
      \gdef\thesection{\Roman{section}}\gdef\theHsection{anexo.\Roman{section}}%
      \phantomsection
      \protected@xdef\leq@current@section{Anexo \theanexo.- \unexpanded{#2}}\gdef\leq@current@subsection{}% índice de ecuaciones
      \addtocontents{stc}{\protect\contentsline{section}{Anexo \theanexo.- #2}{}{\@currentHref}}% índice inicial: solo el nombre, sin página
      \addcontentsline{toc}{section}{Anexo \theanexo.- #2}% marcadores del PDF
      {\noindent\normalfont\Large\bfseries Anexo \theanexo.- #2\par}\nobreak\vspace{1.5ex \@plus .2ex}%
      #1\@ifundefined{anx@lista@\theanexo}{\global\anx@pendientefalse}{}%
    }
    % Índice del anexo: cada apartado escribe una línea en el fichero .atc (anexo, nivel, texto, página, enlace); al
    % empezar el documento se leen las de la compilación anterior y se guardan por anexo (dos pasadas, como cualquier índice).
    % En ese índice los apartados suben un nivel, porque el anexo hace de \section.
    \def\anx@entrada#1#2#3#4#5{\@ifundefined{anx@lista@#1}{\global\@namedef{anx@lista@#1}{}}{}%
      \expandafter\g@addto@macro\csname anx@lista@#1\endcsname{\contentsline{#2}{#3}{#4}{#5}}}
    \tcbuselibrary{breakable}
    \def\anx@indice{\ifanx@pendiente\global\anx@pendientefalse
        \begin{tcolorbox}[breakable, colback=white, colframe=black!35, title={Contenido del anexo}, fonttitle=\small\bfseries,
            before upper={\parindent=0pt\parskip=2pt}]
          \small
          % números más anchos que en el índice inicial (I.1, I.1.1, I.1.1.1; páginas I-1)
          \def\@pnumwidth{2.6em}\def\@tocrmarg{3.4em}%
          \def\l@section##1##2{\addpenalty\@secpenalty\addvspace{1.0em \@plus\p@}\setlength\@tempdima{2.6em}%
            \begingroup\parindent\z@\rightskip\@pnumwidth\parfillskip-\@pnumwidth\leavevmode\bfseries
            \advance\leftskip\@tempdima\hskip-\leftskip##1\nobreak\hfil\nobreak\hb@xt@\@pnumwidth{\hss##2\kern-\p@\kern\p@}\par\endgroup}%
          \def\l@subsection{\@dottedtocline{2}{2.6em}{3.2em}}\def\l@subsubsection{\@dottedtocline{3}{5.8em}{4.1em}}%
          \csname anx@lista@\theanexo\endcsname
        \end{tcolorbox}\par\medskip
      \fi}
    \def\anx@escribir#1#2#3{\ifEnAnexo
        \addtocontents{atc}{\protect\anx@entrada{\theanexo}{#2}{\protect\numberline{\csname the#1\endcsname}#3}{\thepage}{\@currentHref}}%
      \else\addcontentsline{stc}{#1}{\protect\numberline{\csname the#1\endcsname}#3}\fi}
    % Encabezados: dentro de un anexo escriben en su índice (y el primero saca antes ese índice); fuera, en el índice de contenido.
    % Se activan al empezar el documento, cuando ya están definidos los del índice propio (\MyTOC, más abajo).
    \def\anx@subsection[#1]#2{\ifEnAnexo\anx@indice\fi\MyTOC@orig@subsection[{#1}]{#2}\anx@escribir{subsection}{section}{#2}}
    \def\anx@subsubsection[#1]#2{\ifEnAnexo\anx@indice\fi\MyTOC@orig@subsubsection[{#1}]{#2}\anx@escribir{subsubsection}{subsection}{#2}}
    \def\anx@paragraph[#1]#2{\ifEnAnexo\anx@indice\fi\MyTOC@orig@paragraph[{#1}]{#2}\anx@escribir{paragraph}{subsubsection}{#2}}
    \AtBeginDocument{\@starttoc{atc}%
      \let\MyTOC@subsection@nostar\anx@subsection
      \let\MyTOC@subsubsection@nostar\anx@subsubsection
      \let\MyTOC@paragraph@nostar\anx@paragraph}
    \makeatother

    %-------------------------- Cabecera -----------------------%
    % En las páginas numeradas del tema (de la Introducción a la Conclusión), arriba, el título completo del tema tal como está en
    % \title, sin «Tema X.Y.N --»: título y subtítulo seguidos. En letra pequeña; más pequeña aún si no cabe en dos líneas. Va en el
    % margen superior, así que no cambia las páginas del tema. Se enciende con \pagenumbering{arabic} y se apaga con
    % \pagenumbering{gobble} (bibliografía): índices, bibliografía y anexos quedan sin cabecera. Reglas: main/GUIA_TEMAS.md, §4.
    \usepackage{fancyhdr}
    \geometry{headheight=28pt}% hasta tres líneas de letra pequeña
    \ExplSyntaxOn
    \tl_new:N \g_cab_titulo_tl
    \cs_new_protected:Npn \cab_guardar:n #1
      {
        \tl_gset:Nn \g_cab_titulo_tl {#1}
        \regex_replace_once:nnN { \A \s* Tema \s+ \S+ \s* \-\-? \s* } { } \g_cab_titulo_tl % sin «Tema X.Y.N --»
        \regex_replace_once:nnN { \.? \s* \c{\\} (\[ [^\]]* \])? \s* (\c{large} \s*)? } { .\  } \g_cab_titulo_tl % título. subtítulo
        \regex_replace_once:nnN { [\.\s]* \Z } { } \g_cab_titulo_tl % sin subtítulo: nada detrás
      }
    \cs_new_eq:NN \cabGuardarTitulo \cab_guardar:n
    \cs_new:Npn \cabTitulo { \tl_use:N \g_cab_titulo_tl }
    \ExplSyntaxOff
    \makeatletter
    \let\cab@title\title
    \renewcommand{\title}[1]{\cab@title{#1}\cabGuardarTitulo{#1}}
    \newcommand{\cab@letra}{\footnotesize}
    \AtBeginDocument{% si en letra pequeña ocupa más de dos líneas, una talla menos
      \settowidth{\dimen@}{\footnotesize\itshape\cabTitulo}%
      \ifdim\dimen@>1.9\textwidth \renewcommand{\cab@letra}{\scriptsize}\fi}
    \fancypagestyle{tema}{%
      \fancyhf{}%
      \fancyhead[L]{\parbox[b]{\textwidth}{\raggedright\cab@letra\itshape\color{black!65}\cabTitulo}}%
      \fancyfoot[C]{\thepage}%
      \renewcommand{\headrulewidth}{0.4pt}}
    \apptocmd{\pagenumbering}{\ifnum\pdfstrcmp{#1}{arabic}=0 \pagestyle{tema}\else\pagestyle{plain}\fi}{}{}
    \makeatother

    %-------------------------- Lista de figuras (personal) -----------------------%
\makeatletter
\newcommand{\MyListOfFigures}{%
  \clearpage
  \phantomsection
  {\centering\bfseries\Large Índice de figuras\par}%
  \medskip
  % Entrada en nuestro índice de contenido (stc)
  \addcontentsline{stc}{section}{Índice de figuras}%
  % Recuperación del listado (sin cabecera propia)
  \begingroup
    \parindent=0pt\parskip=2pt
    \@starttoc{lof}%
  \endgroup
}
\makeatother

%-------------------------- Bibliografía -----------------------%
\makeatletter
% Contador para las referencias
\newcounter{myref}
% Cabecera de la bibliografía, entra en el índice stc
\newcommand{\MyBibliography}{%
  \clearpage
  \phantomsection
  {\centering\bfseries\Large Bibliografía y referencias\par}%
  \medskip
  \addcontentsline{stc}{section}{Bibliografía y referencias}%
  \setcounter{myref}{0}%
}
\newcommand{\MyBibItem}[4]{%
  \refstepcounter{myref}%
  \noindent[\themyref]\ #1.\ \emph{#2}. #3. #4\par
}
\makeatother

%--------- Establecer cuatro niveles de índice --------%
    \setcounter{tocdepth}{4}   
    \setcounter{secnumdepth}{4}% numera hasta \paragraph

%%%%%%%%%%%%%%%%%%%%%%%%%%%%%%%%

\makeatletter

    %--------- Interlineado Global --------%
    \edef\GlobalBaseLineStretch{\baselinestretch}
    
    %--------- Espaciado de seccinones --------%
    \renewcommand\section{\@startsection{section}
      {1}{\z@}
      {2.5ex \@plus 1ex \@minus .2ex} % espacio antes
      {1.5ex \@plus .2ex}             % espacio después
      {\normalfont\Large\bfseries}    % formato del título
    }
    \renewcommand\subsection{\@startsection{subsection}{2}{\z@}
      {3ex \@plus 0.8ex \@minus .2ex}
      {1.5ex \@plus .2ex}
      {\normalfont\large\bfseries}
    }
    \renewcommand\subsubsection{\@startsection{subsubsection}{3}{\z@}
      {3ex \@plus 0.5ex \@minus .2ex}
      {1ex \@plus .2ex}
      {\normalfont\normalsize\bfseries}
    }
    \renewcommand\paragraph{\@startsection{paragraph}{4}{\z@}%
    {-2ex \@plus -0.5ex \@minus -.2ex}% espacio antes
    {0.8ex \@plus .2ex}% espacio después
    {\normalfont\normalsize\itshape}}% ← cursiva en lugar de negrita

    %--------- Ecuaciones --------%   
    % Variables globales para almacenar títulos y notas
    \gdef\leq@current@section{}%
    \gdef\leq@current@subsection{}%
    \gdef\leq@last@printed@section{}%
    \gdef\leq@last@printed@subsection{}%
    
    % Comando para añadir nota (invisible en el cuerpo)
    \newcommand{\eqnote}[1]{%
      \addcontentsline{leq}{leqnote}{#1}%
    }
    
    % Comando para ecuaciones
    \newcommand{\eqblock}[2]{%
      % Verificar si necesitamos imprimir el título de sección
      \ifx\leq@current@section\leq@last@printed@section\else
        \addcontentsline{leq}{leqsection}{\leq@current@section}%
        \global\let\leq@last@printed@section\leq@current@section
        \gdef\leq@last@printed@subsection{}% Reset subsección al cambiar sección
      \fi
      % Verificar si necesitamos imprimir el título de subsección
      \ifx\leq@current@subsection\leq@last@printed@subsection\else
        \ifx\leq@current@subsection\@empty\else
          \addcontentsline{leq}{leqsubsection}{\leq@current@subsection}%
          \global\let\leq@last@printed@subsection\leq@current@subsection
        \fi
      \fi
      % Ecuación
      \refstepcounter{eqnum}%
      \begin{equation*}
        #1\tag{E.\theeqnum}
      \end{equation*}%
      \addcontentsline{leq}{leqt}{[E.\theeqnum] -- #2}%
      \addcontentsline{leq}{leqm}{#1}%
    }
    
    %%% ----- Índice de ecuaciones ----------%%%
    % Tamaño para []
    \newcommand{\leqIndexFont}{\normalsize}
    \newcommand{\leqTitleFont}{\tiny}
    \newcommand{\leqNoteFont}{\tiny}
    % Cabecera de sección 
    \newcommand*\l@leqsection[2]{%
      \addvspace{25pt}% Mayor separación antes de nueva sección
      \noindent\leqIndexFont
      \makebox[\linewidth]{\rule{.38\linewidth}{0.4pt}\ \ \textbf{  \MakeUppercase{#1}}\ \ \rule{.38\linewidth}{0.4pt}}%
      \par\vspace{-10pt}
    }
    % Cabecera de subsección 
    \newcommand*\l@leqsubsection[2]{%
      \par\addvspace{15pt}%
      \noindent\leqIndexFont #1
      \par\vspace{-7pt}
    }
    % Título de ecuación 
    \newcommand*\l@leqt[2]{%
    \par\addvspace{10pt}%
      \par\noindent\leqTitleFont #1\par
    }
    % Miniatura de ecuación
    \newcommand*\l@leqm[2]{%
      \vspace{0.3\baselineskip}%
      \par\scriptsize\noindent\hspace*{1.5em}\(#1\)\par%
    }
    % Nota de ecuación 
    \newcommand*\l@leqnote[2]{%
      \vspace{0.2\baselineskip}%
      \par\noindent\leqNoteFont\hspace*{1.5em}\textbf{Nota.--} #1\par\vspace{0.5\baselineskip}%
    }
    % Generación del índice
    \newcommand{\mylistofequations}{%
      \clearpage
      \phantomsection
      % Título visual de la lista (ajústalo a tu gusto):
      {\centering\bfseries\Large Índice de ecuaciones\par}
      % Entrada en nuestro índice stc:
      \addcontentsline{stc}{section}{Índice de ecuaciones}%
      % Llamada a tu lista real (usa el nombre exacto de tu macro):
      \@starttoc{leq}%
    }
    
    %--------- Captura de títulos de sección --------%
    \let\leq@original@section\section
    \let\leq@original@subsection\subsection
    % Flag para saber si estamos en una sección asterisco
    \newif\ifLeqStarSection
    \LeqStarSectionfalse
    
    % Redefinir \section CON CONTROL DE ASTERISCO
    \renewcommand{\section}{%
      \@ifstar{\leq@section@star}{\@dblarg\leq@section@nostar}%
    }
    \def\leq@section@star#1{%
      \global\LeqStarSectiontrue
      \gdef\leq@current@section{#1}%
      \gdef\leq@current@subsection{}%
      \leq@original@section*{#1}%
      \global\LeqStarSectionfalse
    }
    \def\leq@section@nostar[#1]#2{%
      \global\LeqStarSectionfalse
      \gdef\leq@current@section{#2}%
      \gdef\leq@current@subsection{}%
      \leq@original@section[#1]{#2}%
    }
    % Redefinir \subsection
    \renewcommand{\subsection}{%
      \@ifstar{\leq@subsection@star}{\@dblarg\leq@subsection@nostar}%
    }
    \def\leq@subsection@star#1{%
      \global\LeqStarSectiontrue
      \gdef\leq@current@subsection{#1}%
      \leq@original@subsection*{#1}%
      \global\LeqStarSectionfalse
    }
    \def\leq@subsection@nostar[#1]#2{%
      \global\LeqStarSectionfalse
      \gdef\leq@current@subsection{#2}%
      \leq@original@subsection[#1]{#2}%
    }

    % ----- Restauración de valores Globales de estilo ----------%
    % Flags
    \newif\ifInFootnote
    \InFootnotefalse
    \pretocmd{\@footnotetext}{\InFootnotetrue}{}{}
    \apptocmd{\@footnotetext}{\InFootnotefalse}{}{}
    % Profundidad de listas (soporta anidadas)
    \newcount\MyListDepth \MyListDepth=0
    \AtBeginEnvironment{la}{\advance\MyListDepth by 1}
    \AfterEndEnvironment{la}{\advance\MyListDepth by -1}
    \AtBeginEnvironment{lnum}{\advance\MyListDepth by 1}
    \AfterEndEnvironment{lnum}{\advance\MyListDepth by -1}
    \AtBeginEnvironment{itemize}{\advance\MyListDepth by 1}
    \AfterEndEnvironment{itemize}{\advance\MyListDepth by -1}
    \AtBeginEnvironment{enumerate}{\advance\MyListDepth by 1}
    \AfterEndEnvironment{enumerate}{\advance\MyListDepth by -1}
    % Restauración diferida: hazlo al INICIO del próximo párrafo real
    \newif\if@pendingRestore
    \@pendingRestorefalse
    \newcommand{\RequestRestore}{\global\@pendingRestoretrue}
    \everypar{%
      \if@pendingRestore
        \global\@pendingRestorefalse
        \ifInFootnote\else
          \ifnum\MyListDepth=0
            \setlength{\parskip}{\GlobalParskip}%
            \setstretch{\GlobalBaseLineStretch}%
          \fi
        \fi
      \fi
    }
    % Al final de bloques:
    \AfterEndEnvironment{equation}{\RequestRestore}
    \AfterEndEnvironment{equation*}{\RequestRestore}
    \AfterEndEnvironment{la}{\RequestRestore}
    \AfterEndEnvironment{lnum}{\RequestRestore}
    % Al final de macros:
    \apptocmd{\eqblock}{\RequestRestore}{}{}
    
\makeatother

% ===================== ToC propio con 4 niveles (único bloque) =====================
\makeatletter

% --- Escritores: añadimos una línea al .stc en cada cabecera numerada ---
% Guardamos las versiones actuales para llamar al comportamiento normal
\let\MyTOC@orig@section\section
\let\MyTOC@orig@subsection\subsection
\let\MyTOC@orig@subsubsection\subsubsection
\let\MyTOC@orig@paragraph\paragraph

% SECTION
\renewcommand{\section}{%
  \@ifstar{\MyTOC@section@star}{\@dblarg\MyTOC@section@nostar}%
}
\def\MyTOC@section@star#1{%
  \MyTOC@orig@section*{#1}% sin entrada al .stc
}
\def\MyTOC@section@nostar[#1]#2{%
  \MyTOC@orig@section[{#1}]{#2}% comportamiento normal (escribe al .toc)
  % Escribimos también nuestra línea al .stc (texto con \numberline)
  \addcontentsline{stc}{section}{\protect\numberline{\thesection}#2}%
}

% SUBSECTION
\renewcommand{\subsection}{%
  \@ifstar{\MyTOC@subsection@star}{\@dblarg\MyTOC@subsection@nostar}%
}
\def\MyTOC@subsection@star#1{%
  \MyTOC@orig@subsection*{#1}%
}
\def\MyTOC@subsection@nostar[#1]#2{%
  \MyTOC@orig@subsection[{#1}]{#2}%
  \addcontentsline{stc}{subsection}{\protect\numberline{\thesubsection}#2}%
}

% SUBSUBSECTION
\renewcommand{\subsubsection}{%
  \@ifstar{\MyTOC@subsubsection@star}{\@dblarg\MyTOC@subsubsection@nostar}%
}
\def\MyTOC@subsubsection@star#1{%
  \MyTOC@orig@subsubsection*{#1}%
}
\def\MyTOC@subsubsection@nostar[#1]#2{%
  \MyTOC@orig@subsubsection[{#1}]{#2}%
  \addcontentsline{stc}{subsubsection}{\protect\numberline{\thesubsubsection}#2}%
}

% PARAGRAPH
\renewcommand{\paragraph}{%
  \@ifstar{\MyTOC@paragraph@star}{\@dblarg\MyTOC@paragraph@nostar}%
}
\def\MyTOC@paragraph@star#1{%
  \MyTOC@orig@paragraph*{#1}%
}
\def\MyTOC@paragraph@nostar[#1]#2{%
  \MyTOC@orig@paragraph[{#1}]{#2}%
  % Si no numeras los \paragraph, cambia \theparagraph por texto plano sin \numberline
  \addcontentsline{stc}{paragraph}{\protect\numberline{\theparagraph}#2}%
}

% --- Impresor del índice propio (lee el .stc) ---
\newcommand{\MyTOC}{%
  \par\bigskip
  {\centering\bfseries\Large Índice de contenido\par}%
  \medskip
  \begingroup
    \parindent=0pt\parskip=2pt
    \@starttoc{stc}%
  \endgroup
}

\makeatother
% ===================== fin ToC propio =====================


%%%%%%%%%%%%%%


% --- Datos del tema ---
\title{Tema 3.A.14 -- Teoría de Juegos.\\
\large Principales conceptos. Aplicaciones, con especial referencia a las subastas}
\author{Marcelo Ortega}
\date{Diciembre 2025}
\newcommand{\TemaCode}{3A14}

\oldtitle{Tema nuevo}
\changerationale{Se busca tratar de forma independiente la teoría de juegos en un tema separado, desvinculándolo del tratamiento dentro del tema de oligopolio que se le daba en la anterior versión del temario. Además, gracias a esta separación, se pretende que el opositor incluya entre otras cuestiones desarrollos bayesianos y más concretamente pueda introducir la teoría de subastas.}

\begin{document}


\maketitle
\changebox

\newpage

% --- Página de resumen y modificaciones ---
\puntosclave{ %% Escribir puntos clave Apuntes CECO
    
    }
\vspace{1cm}

\modificaciones{
    Estructurar el Tema de tal forma que se organice con el siguiente minutero:
    2’ Introducción; 5’ Conceptos básicos; 12’ Juegos no-cooperativos; 5’ Juegos cooperativos (OJO! Cambiar esta distribución, se ha introducido la parte de cooperativos en Análisis de mercados (IV) [Tema 3.A.19]); 5’ Diseño de Mecanismos y especial referencia a las subastas; 1’ Conclusión
    
    En Juegos cooperativos introducir concepto de Variaciones “conjunturales” = 0 (conocimiento de las ventajas de la cooperación) para ligarlo con la siguiente parte “Juegos cooperativos” (“Si levantamos el supuesto de que las variaciones ”)
    
    Introducir más ejemplos en la parte de Juegos cooperativos, que no sea tan técnica la exposición, muy dinámica y narrada; Utilizar lo característico de la Teoría del Oligopolio, por ejemplo, u otros ejemplos pero fáciles y sencillos

    \href{https://www.gametheory.net/news/news.pl?Concept=Auc&highlight=AUC}{Game Theory}
    }

\newpage
\MyTOC
\newpage

\mylistofequations
\clearpage

\MyListOfFigures
\clearpage


\pagenumbering{arabic}
\setcounter{page}{1}


\section{Introducción}\label{intro}





\subsection{Contextualización}\label{context}
 

\subsubsection{Relevancia}\label{importance}


\subsubsection{Historia}\label{history}


\subsubsection{Problemática}\label{problem} 


\subsection{Estructura de la exposición}\label{structure}












































\clearpage
\section{Principales conceptos}
\puntosclave{
    Un \textbf{juego} describe una situación de interdependencia estratégica: el resultado de cada participante depende de sus decisiones y de las de los demás. Se describe con sus jugadores (qué quieren, cómo razonan y qué saben), sus reglas (qué pueden hacer, cuándo y qué observan al decidir) y sus estrategias, que son planes completos de actuación; se representa en forma extensiva (el árbol) o en forma normal o estratégica (la matriz).

    Dos dimensiones ordenan los juegos y sus conceptos de solución: el orden de las jugadas (\textbf{juegos simultáneos o dinámicos}) y lo que cada uno sabe de los demás (\textbf{información completa o incompleta}). La información incompleta se trata considerando que el tipo de preferencias de cada jugador es la realización de una variable aleatoria; así, el juego se convierte en uno de información imperfecta (HARSANYI).

    Cada concepto de solución exige más a las expectativas de los jugadores que el anterior. En un juego simultáneo con información completa, la \textbf{dominancia} permite predecir sin anticipar al rival, pero rara vez basta; el \textbf{equilibrio de NASH} (EN) exige conjeturas correctas y existe siempre, al menos en estrategias mixtas, si el juego es finito, pero puede haber varios sin que quede claro cuál cabe esperar.

    Con información incompleta, el \textbf{equilibrio bayesiano de NASH} (EBN) es una extensión del EN a los juegos bayesianos: cada tipo de cada jugador juega su mejor respuesta dadas las probabilidades de los tipos ajenos. Es la herramienta de las subastas y del diseño de mecanismos.

    En los \textbf{juegos dinámicos} con información completa, el EN no descarta las amenazas no creíbles. El \textit{principio de racionalidad secuencial} lleva al \textbf{equilibrio de NASH perfecto en subjuegos} (ENPS), que en horizonte finito se halla por inducción hacia atrás y en horizonte infinito exige otros métodos.

    En los juegos dinámicos con información imperfecta, la perfección en subjuegos no basta: se introducen las creencias y, con ellas, el equilibrio bayesiano perfecto débil (EBPD). Al imponer requisitos adicionales a las creencias se obtienen sus refuerzos: el equilibrio bayesiano perfecto (EBP), el equilibrio secuencial y el equilibrio perfecto de mano temblorosa (en forma extensiva). Cuando el jugador informado mueve primero, su jugada es una señal: en los equilibrios separadores revela su tipo y en los agrupadores no revela nada.
}

Durante mucho tiempo, la teoría económica estudió decisiones individuales tomadas en un entorno que el decisor toma como dado: el consumidor y la empresa precio-aceptantes reaccionan a precios que nadie controla, y el monopolista se enfrenta a una curva de demanda que no responde estratégicamente. 

En cuanto unos pocos agentes se afectan mutuamente esa simplificación deja de valer: lo que conviene a cada uno depende de lo que hagan los demás, que a su vez razonan sobre lo que hará él. La \textbf{teoría de juegos} es el instrumental que formaliza esa \textbf{interdependencia estratégica}. Nace con el teorema del minimax de VON NEUMANN (1928) y con VON NEUMANN y MORGENSTERN (1944), que la concibieron para los juegos de suma cero, y adquiere su forma actual con NASH (1950a, 1951), SELTEN (1965, 1975) y HARSANYI (1967-1968)\footnote{%
  NASH, HARSANYI y SELTEN recibieron el premio Nobel de Economía en 1994. Lo han recibido después, por desarrollos de la teoría de juegos o por sus aplicaciones directas, AUMANN y SCHELLING (2005), HURWICZ, MASKIN y MYERSON (2007, diseño de mecanismos), ROTH y SHAPLEY (2012, asignaciones estables y diseño de mercados) y MILGROM y WILSON (2020, subastas).
}.

La teoría se divide en dos ramas según su unidad de análisis. La \textbf{no cooperativa} toma como unidad al jugador individual y supone que no puede firmar acuerdos vinculantes: si coopera, ha de ser porque le conviene hacerlo. La \textbf{cooperativa} toma como unidad las coaliciones y supone que los acuerdos se pueden hacer cumplir, de modo que la pregunta pasa a ser qué repartos son estables o justos ([\authorfont{Ver Tema 3.A.19}]). Este bloque se ocupa de la primera, que es la que sostiene el diseño de mecanismos y la teoría de subastas.

Describir un juego es relativamente sencillo; lo difícil es predecir cómo se jugará. La teoría no ofrece un único concepto de solución, sino una familia, y cada uno nace de un fallo del anterior. Lo que los distingue es cuánto exigen a las expectativas de los jugadores: la dominancia no exige prever lo que harán los demás; el equilibrio de Nash exige que cada uno acierte las estrategias de los otros; y sus refinamientos exigen además que esas expectativas sean coherentes en situaciones que no llegan a producirse o que se actualicen con lo que revelan las jugadas. Un concepto más exigente predice con más precisión, pero corre el riesgo de no existir o de descansar en supuestos poco verosímiles sobre la capacidad de razonar de los jugadores. Esa tensión entre existencia y precisión recorre todo el bloque.

\subsection{Reglas del juego: ¿cómo describimos un juego?}

Un \textbf{juego} es la representación formal de una situación en la que un número determinado de individuos interactúa en un entorno de interdependencia estratégica. 

Para describirla hay que responder a cuatro preguntas: quién participa (los \textit{jugadores}); cuándo actúa cada uno, qué puede hacer y qué sabe al hacerlo (las \textit{reglas}); qué resultado produce cada combinación de decisiones (los \textit{resultados}); y cómo valora cada jugador esos resultados (los \textit{pagos}) (MAS-COLELL, WHINSTON y GREEN, 1995). RASMUSEN (2007) lo resume en cuatro elementos (jugadores, acciones, pagos e información), que juntos forman las reglas del juego. 

A partir de ellos se definen las estrategias, y el conjunto se representa en forma extensiva o en forma normal. Todo ello es la descripción del juego, el dato de partida; la predicción de cómo se jugará, el concepto de solución, viene después.

Un ejemplo muy sencillo nos acompañará en todo el bloque. Una empresa \textbf{entrante} (J1) estudia si entrar en un mercado que hoy abastece en solitario una empresa \textbf{incumbente} (J2). Primero decide J1 si entra o no entra; después, J2, que ha visto lo que ha hecho J1, fija un precio alto o bajo.

\imagenfit{JE_ejemplo_arbol.pdf}{Un juego de entrada}{Elaboración propia.}

Cada uno de los cuatro desenlaces posibles se designa con una letra, y en él cada jugador obtiene un pago, que lleva su número como subíndice: $(a_1,a_2)$ si J1 entra y J2 mantiene el precio alto; $(b_1,b_2)$ si entra y J2 lo baja; $(c_1,c_2)$ si no entra y el precio es alto; y $(d_1,d_2)$ si no entra y el precio es bajo.\footnote{%
  Usamos letras y no cifras porque lo que decide cómo se juega es el orden entre los pagos: cambiándolo obtendremos, con el mismo juego, cada uno de los conceptos del bloque. 
  
  La lectura más natural del ejemplo sugiere algunas relaciones, que mantendremos salvo que se diga otra cosa. A J1 le interesa entrar si J2 se acomoda, pero no si desata una guerra de precios, y si no entra gana lo mismo fije J2 el precio que fije ($a_1>c_1=d_1>b_1$). A J2, que sin competidor es monopolista, le conviene en ese caso el precio alto ($c_2>d_2$), y prefiere seguir solo a compartir el mercado ($c_2>a_2$). Lo que el ejemplo deja abierto, y explotaremos, es qué prefiere J2 si J1 entra: acomodarse ($a_2>b_2$) o la guerra de precios ($b_2>a_2$).
}

\subsubsection{Jugadores: ¿quiénes son y qué quieren?}

Los jugadores son los agentes que toman decisiones: consumidores, empresas, sindicatos, gobiernos o bancos centrales. Su número puede ser pequeño, como en un duopolio, o muy grande, como en los modelos de poblaciones. A menudo se añade un jugador ficticio, la \textbf{naturaleza} (el jugador $0$), que no tiene preferencias y elige al azar con probabilidades que todos conocen: es la forma de introducir en el juego la incertidumbre exógena, como el estado de la demanda o las características de un rival.

En el ejemplo, los jugadores son las dos empresas, J1 y J2, y la naturaleza no interviene, porque nada se decide al azar.

Los jugadores son \textbf{racionales}: cada uno elige lo que maximiza su pago esperado, dadas sus conjeturas sobre lo que harán los demás. La teoría de juegos añade un supuesto que la teoría de la decisión individual no necesita: que la racionalidad y la propia estructura del juego son de \textbf{conocimiento común}. No basta con que todos sean racionales y conozcan el juego; hace falta que todos sepan que todos lo son, que todos sepan que todos lo saben, y así indefinidamente. El concepto procede de LEWIS (1969) y lo formalizó AUMANN (1976). Esto es crucial, pues permite a cada jugador ponerse en el lugar de los demás y anticipar sus decisiones; los distintos conceptos de solución se apoyan en él en grados muy distintos, y es también el supuesto más discutido de la teoría.\footnote{%
    \textbf{NOTA.-} La evidencia experimental sugiere que pocos jugadores encadenan más de dos o tres niveles de razonamiento sobre lo que piensan los demás. 
    
    En el juego de adivinar una fracción de la media de las cifras elegidas por todos (el concurso de belleza de KEYNES, 1936), NAGEL (1995) observó elecciones muy alejadas del único equilibrio, al que se llega eliminando iterativamente estrategias dominadas. Los modelos de niveles de razonamiento (STAHL y WILSON, 1995) y de jerarquía cognitiva (CAMERER, HO y CHONG, 2004) explican esas elecciones mejor que el supuesto de conocimiento común.
} En el ejemplo, para decidir si entra, J1 ha de prever qué precio fijará J2, y solo puede hacerlo si sabe que J2 es racional y conoce sus pagos.

\paragraph{Información: ¿cómo se interrelaciona la información?}
Lo que cada jugador sabe de los demás define la primera gran distinción. Un juego es de \textbf{información completa} si la estructura del juego, incluidos los pagos de todos, es de conocimiento común: cada uno sabe qué quieren los demás, aunque no sepa qué van a hacer. Es de \textbf{información incompleta} si algún jugador desconoce las características relevantes de otro, normalmente sus pagos: el entrante potencial no sabe si los costes del incumbente son altos o bajos, y quien puja en una subasta no sabe cuánto valoran el objeto los demás. 

Lo habitual es que esa información sea \textit{privada}: cada uno conoce sus propias características y desconoce las ajenas, de modo que la información es \textbf{asimétrica}, y quien no la tiene solo puede valorar la situación en términos probabilísticos. Si nadie conociera el dato y todos compartieran la misma probabilidad sobre él, la incertidumbre sería simétrica y bastaría con trabajar con pagos esperados; lo que da a estos juegos su interés es la asimetría.\footnote{%
    \textbf{OJO.-} Información incompleta no equivale a que intervenga la naturaleza. Si la naturaleza mueve y todos observan su jugada (una tirada de dados a la vista de todos), el juego sigue siendo de información completa. Lo que define la información incompleta es que algún jugador no conoce los pagos o las características de otro; la naturaleza es solo el recurso con que HARSANYI (1967-1968) la modeliza, como vemos a continuación.
} 

HARSANYI (1967-1968) resumió la información privada de cada jugador en su \textbf{tipo}($\theta$) : \textbf{todo lo que sabe él y no saben los demás}. Los demás no conocen su tipo, pero sí la probabilidad con que puede ser de cada clase, que es de conocimiento común.\footnote{%
  \textbf{NOTA.-} Los juegos de información incompleta plantean un problema de representación propio: quien desconoce los pagos de otro no sabe a qué juego está jugando. La solución de HARSANYI (1967-1968), la \textbf{transformación de HARSANYI}, consiste en añadir al principio del árbol una jugada de la naturaleza que elige el tipo de cada jugador con probabilidades que todos conocen. Cada jugador observa su propio tipo, pero no el de los demás.
} El juego de información incompleta se convierte así en un \textbf{juego bayesiano}: un juego en el que algún jugador no ha visto una jugada anterior, la de la naturaleza, es decir, un juego de información imperfecta, concepto que precisaremos al tratar las reglas. En él, la estrategia de un jugador es un plan para cada uno de sus tipos (qué haría si fuera de cada clase), y quien no conoce el tipo de otro forma sus conjeturas ponderando lo que haría cada tipo por su probabilidad.

\imagenfit{JE_ejemplo_informacion.pdf}{Información completa, información incompleta y transformación de HARSANYI}{Elaboración propia a partir de HARSANYI (1967-1968).}

El ejemplo lo muestra con claridad. En el panel (a) de la figura, la información es completa: cada empresa conoce los pagos de la otra en cada desenlace. Deja de serlo si J1 no sabe qué prefiere J2 cuando entra. Supongamos que el coste de J2 puede ser alto o bajo y que solo J2 lo sabe. Con un coste alto, a J2 le conviene acomodarse; con uno bajo, la guerra de precios le sale a cuenta. Ese coste es su tipo: si es de tipo \textit{acomodaticio} ($\theta_{ac}$), sus pagos son los de siempre; si es de tipo \textit{agresivo} ($\theta_{ag}$), son otros, que distinguimos con un acento circunflejo, con $\hat{b}_2>\hat{a}_2$. J1, en cambio, tiene que decidir si entra sin saber a qué incumbente se enfrenta. El panel (b) recoge la transformación de HARSANYI. La naturaleza, el jugador $0$, elige primero el tipo de J2: acomodaticio con probabilidad $p$ y agresivo con probabilidad $1-p$. J2 observa el resultado; J1 no, de modo que sus dos nodos de decisión, unidos por una línea discontinua, son indistinguibles para él. Resolveremos este juego al tratar el equilibrio bayesiano de NASH.

\paragraph{Función de pagos:}

Cada jugador tiene preferencias sobre los resultados posibles, representadas por una \textbf{función de pagos}: es usual llamar función de pagos a su función de utilidad y pago a su nivel de utilidad. 

El pago recoge todo lo que el jugador valora del resultado, no solo el dinero; si a un jugador le importa el bienestar de otro, eso ya está en su pago, de modo que la teoría no supone egoísmo. Cuando los resultados son aleatorios, porque mueve la naturaleza o porque algún jugador elige al azar, se supone que cada posibilidad (\textit{prospect}) aleatoria se evalúa por su utilidad esperada, en el sentido de VON NEUMANN y MORGENSTERN (1944). Por eso los pagos de un juego no expresan solo un orden: su escala importa, porque con ellos se calculan medias ponderadas por probabilidades.

\imagenfit{JE_ejemplo_pagos.pdf}{Pagos ciertos y pagos esperados}{Elaboración propia.}

En el ejemplo, los pagos son los beneficios de cada empresa y figuran en los nodos finales del árbol, primero el de J1 y después el de J2. Que su escala importa se ve en cuanto aparece la incertidumbre. En el panel (a) de la figura, J1 sabe que, si entra, J2 mantendrá el precio alto: entrar le da $a_1$ con seguridad y no entrar, $c_1$, y le basta el orden de los pagos para decidir. En el panel (b), J1 cree que, si entra, J2 bajará el precio con probabilidad $q$: entrar le da en promedio $(1-q)\,a_1+q\,b_1$, y le conviene si ese pago esperado supera a $c_1$, es decir, si $q<(a_1-c_1)/(a_1-b_1)$. El umbral no depende solo del orden de los pagos, sino de sus distancias: si la guerra de precios fuese más dañina para J1 (un $b_1$ menor), el orden sería el mismo, pero J1 solo entraría con una probabilidad de guerra menor.

\subsubsection{Reglas: ¿qué son las acciones y el orden de juego?}

Las reglas fijan, en primer lugar, las \textbf{acciones}: la decisión o movimiento concreto que un jugador puede realizar en un punto determinado del juego. Fijan también el \textbf{orden} en que se decide. En un \textbf{juego simultáneo} (o estático) cada jugador decide sin conocer lo que han elegido los demás; en un \textbf{juego secuencial} (o dinámico) algunos deciden después de observar lo que otros han hecho.

En el ejemplo, J1 elige entre entrar y no entrar, y J2, entre un precio alto y uno bajo; como J1 decide primero y J2 después, el juego es secuencial.

Debemos tener siempre muy presente que la \textbf{simultaneidad es una cuestión de información}, no de calendario: en una subasta a sobre cerrado los sobres pueden entregarse en días distintos, pero el juego es simultáneo porque nadie conoce las pujas ajenas al escribir la suya. Del mismo modo, el juego de entrada pasaría a ser simultáneo si J2 tuviera que fijar su precio sin saber aún si J1 ha entrado, aunque lo fijara después.

\paragraph{Secuencialidad: ¿cómo influye el orden sobre la información?}
Lo que cada jugador sabe al mover es el elemento más delicado de las reglas, y se describe con los \textbf{conjuntos de información} (\textit{information sets}). Cada vez que un jugador ha de decidir, su conjunto de información reúne todas las situaciones (los nodos del árbol del juego) que, dado lo que ha observado, podrían haberse producido hasta ese momento\footnote{%
    Técnicamente, se dice que los conjuntos de información forman una «partición» de los nodos de decisión de cada jugador.
}. Si el conjunto tiene un único nodo (en inglés, \textit{singleton}), el jugador conoce toda la historia anterior, porque un nodo resume completamente lo sucedido hasta él. Si tiene varios, el jugador \textit{sabe qué ha podido pasar}, pero \textit{no lo que realmente ha pasado}. La noción procede de VON NEUMANN, a quien el farol del póquer llevó a modelizar jugadores que no ven las cartas del rival, y KUHN (1953) le dio su formulación general sobre el árbol del juego.

Con ellos se define la segunda gran distinción. Un juego es de \textbf{información perfecta} si todos los conjuntos de información tienen un único nodo: cada jugador, al mover, conoce todas las jugadas anteriores, como en el ajedrez o en el duopolio de STACKELBERG (1934). Si algún conjunto tiene varios nodos, el juego es de \textbf{información imperfecta}: todo juego simultáneo lo es, porque quien mueve no ha visto la jugada de los otros.

\imagenfit{JE_ejemplo_conjuntos.pdf}{Conjuntos de información: información perfecta e imperfecta}{Elaboración propia.}

La figura lo muestra con el ejemplo, en el que J2 decide en dos nodos, tras la entrada y tras la no entrada. En el panel (a), J2 ve lo que ha hecho J1, de modo que cada uno de sus nodos es un conjunto de información distinto: el juego es de información perfecta. En el panel (b), J2 fija el precio sin saber si J1 ha entrado: sus dos nodos forman un único conjunto de información, que en el árbol se dibuja uniéndolos con una línea discontinua, y el juego es de información imperfecta. A todos los efectos, el juego del panel (b) es simultáneo.

Por tanto, información perfecta e información completa son cosas distintas. Por un lado, la perfección se refiere a lo que se \textit{observa de las jugadas} y es una regla del juego. Por otro lado, la completitud, que se refiere a lo que se \textit{sabe de los pagos} ajenos, es un rasgo de los jugadores.\footnote{%
  El duopolio de COURNOT (1838) es un juego de información completa e imperfecta; el de STACKELBERG, de información completa y perfecta; una subasta a sobre cerrado, de información incompleta, que la transformación de HARSANYI convierte, como vimos, en imperfecta.
} La distinción condiciona el concepto de solución, como se verá al tratar la credibilidad y las creencias.

\subsection{Representación del juego: ¿cómo podemos representar el juego?}
Como habíamos anticipado, jugadores, pagos y reglas son las definiciones necesarias para poder hablar de un juego: con ellas sabemos quién decide, qué puede hacer, cuándo, sabiendo qué y con qué consecuencias. Falta el objeto sobre el que razona cada jugador cuando se pregunta cómo jugar, que no es una acción suelta, sino un plan para todo el juego.

Ese plan, la estrategia, es también el concepto común del que se desprenden las dos formas de representar un juego, aunque de manera distinta: la forma extensiva dibuja las reglas, y en ella se leen las estrategias; la forma normal se construye directamente con ellas.

\subsubsection{Perfil de estrategias: ¿qué es y qué determina una estrategia?}

A partir de los jugadores y las reglas se define la pieza más importante, y conceptualmente más difícil, de la teoría: la \textbf{estrategia}. 

Una estrategia es un \textbf{{plan de acción completo}}: especifica qué hará el jugador en cada una de las situaciones en que puede tener que mover, es decir, en cada uno de sus conjuntos de información, porque cada conjunto de información representa una circunstancia distinguible en la que puede ser preciso que decida.\footnote{%
  Elegir una estrategia es como escribir, antes de que empiece el juego, un libro de instrucciones tan detallado que un representante podría jugar en nombre del jugador sin consultarle, limitándose a buscar en él qué hacer en cada caso.
} 

En un juego simultáneo de una sola jugada, estrategia y acción coinciden. En uno secuencial no: quien mueve después debe prever una respuesta para cada cosa que pueda haber hecho quien movió antes, de modo que tiene tantas estrategias como combinaciones de respuestas.

Así, en el ejemplo, J2 tiene dos acciones, precio alto o bajo, pero cuatro estrategias, porque ha de decir qué precio fija si J1 entra y qué precio fija si no entra. Las escribimos con el precio tras la entrada en primer lugar: (Alto, Alto), (Alto, Bajo), (Bajo, Alto) y (Bajo, Bajo). J1, que mueve primero y una sola vez, tiene dos estrategias, que coinciden con sus dos acciones. El panel (a) de la figura dibuja cada estrategia de J2 como una manera de elegir una rama en cada uno de sus nodos; el panel (b) muestra que, si J2 no distinguiera sus dos nodos, tendría que elegir la misma rama en ambos y solo le quedarían dos estrategias.

\imagenfit{JE_ejemplo_estrategias.pdf}{Estrategias de J2 en el juego secuencial y en el simultáneo}{Elaboración propia.}

Por tanto, como plan completo, \textbf{una estrategia especifica acciones incluso para situaciones que no llegarán a producirse si todos la siguen}. No es un detalle formal: lo que un jugador haría en las situaciones que no se dan determina lo que los demás hacen en las que sí se dan. La amenaza de una guerra de precios si entra un competidor forma parte de la estrategia del incumbente aunque, si la amenaza funciona, nadie entre y la guerra no llegue a librarse. En el ejemplo, es la estrategia (Bajo, Alto) de J2: si J1 se la cree, no entra, y J2 obtiene $c_2$ gracias a lo que habría hecho en una situación que no llega a producirse. Si J2 es acomodaticio ($a_2>b_2$), la amenaza es además dudosa: si J1 entrara, a J2 le convendría más acomodarse que desatar la guerra. Si esas amenazas son creíbles es, como veremos, el problema que obliga a refinar el equilibrio de Nash en los juegos dinámicos.

La combinación de una estrategia por jugador se llama \textbf{perfil de estrategias}, $s=(s_1,\ldots,s_I)$, y suele escribirse $s=(s_i,s_{-i})$ para separar la estrategia del jugador $i$ de las de todos los demás, $s_{-i}$. En el ejemplo hay $2\times4=8$ perfiles; uno de ellos es (Entrar; Alto, Alto): J1 entra y J2 fija el precio alto tanto si entra como si no.

\paragraph{Taxonomía: ¿cómo influye la información sobre la estrategia?}

La información determina, ante todo, cuántas estrategias tiene cada jugador. Como una estrategia asigna una acción a cada conjunto de información, el jugador tiene tantas como combinaciones de acciones pueda formar: el producto del número de acciones disponibles en cada uno de sus conjuntos de información. Cuanto más observa un jugador, más conjuntos de información tiene y más ricas son sus estrategias, porque puede condicionar lo que hace a lo que ha visto. Las cuatro estrategias de J2 en el ejemplo son $2\times2$: dos acciones en cada uno de sus dos conjuntos de información. Si J2 no observara la decisión de J1, tendría un solo conjunto de información y dos estrategias, precio alto o bajo, porque no puede condicionar su acción a lo que no ve.

Cuando las acciones son continuas, la riqueza que da la información se ve todavía mejor: la estrategia de quien observa deja de ser un número y pasa a ser una función. En el duopolio de COURNOT (1838), la estrategia de cada empresa es una cantidad. En el de STACKELBERG (1934), la seguidora observa la cantidad de la líder antes de decidir, y su estrategia es una regla que asigna una cantidad a cada cantidad posible de la líder; la que sigue en equilibrio es su función de reacción ([\authorfont{Ver Tema 3.A.19}]). Lo mismo ocurre con la información privada: en una subasta a sobre cerrado, cada participante conoce su valoración, y su estrategia es una función que asigna una puja a cada valoración posible, $b(\theta)$: es la estrategia del juego bayesiano que resulta de la transformación de HARSANYI, como veremos en el bloque de subastas.

Una estrategia que elige un curso de acción determinado es una \textbf{estrategia pura}. Un jugador puede también elegir al azar entre varias, asignando a cada una una probabilidad: es una \textbf{estrategia mixta}, como la del lanzador de penaltis que no quiere ser previsible o la de la inspección fiscal que no revisa todas las declaraciones.\footnote{%
    Su sentido se verá al tratar el equilibrio de Nash; su definición formal, y la de la estrategia de comportamiento, está [\authorfont{Ver Anexo \ref{anx:no_cooperativos}}].
}

En los juegos dinámicos se puede aleatorizar de dos maneras, y la diferencia vuelve a ser de información. Una estrategia mixta sortea al principio uno de los planes completos: en el panel (a) de la figura, J2 asigna una probabilidad a cada una de sus cuatro estrategias puras, de $\lambda_1$ a $\lambda_4$, de modo que necesita tres números (el cuarto es lo que falta hasta uno). Una \textbf{estrategia de comportamiento} sortea la acción cada vez que toca mover, en cada conjunto de información por separado: en el panel (b), a J2 le bastan dos números, la probabilidad de bajar el precio si J1 entra, $\beta_E$, y la de bajarlo si no entra, $\beta_N$.

\imagenfit{JE_ejemplo_aleatorizar.pdf}{Estrategia mixta y estrategia de comportamiento}{Elaboración propia.}

KUHN (1953) demostró que, si ningún jugador olvida lo que sabía ni lo que hizo (\textbf{memoria perfecta}), ambas formas son equivalentes: toda estrategia mixta tiene una de comportamiento que produce la misma distribución de probabilidad sobre los resultados, y viceversa. En el ejemplo, la de comportamiento equivalente asigna a cada nodo la probabilidad conjunta de los planes que bajan el precio en él: $\beta_E=\lambda_3+\lambda_4$ y $\beta_N=\lambda_2+\lambda_4$. Por eso los juegos dinámicos se analizan con estrategias de comportamiento, mucho más manejables, sin perder generalidad; los juegos que se usan en economía suponen casi siempre memoria perfecta.

\subsubsection{Representación gráfica: ¿cómo aproximar un juego visualmente?}

Una vez definido el juego y comprendido su concepto clave, la estrategia, veamos cómo puede representarse gráficamente. Las dos formas se distinguen por cuánto conservan de la descripción: el árbol, que ya hemos usado con el juego de entrada, la conserva entera; la matriz se queda solo con las estrategias y sus pagos.

Que las estrategias contengan toda la información relevante del juego, de modo que lo que la matriz omite no importe, es una tesis atractiva, pero discutida, como veremos al comparar las dos formas.

\paragraph{Forma extensiva: ¿cuál es su utilidad?}
La \textbf{forma extensiva} representa el juego como un árbol (\textit{game tree}): cada nodo de decisión indica a quién le toca mover; cada rama, una acción posible; los nodos que pertenecen a un mismo conjunto de información se unen con una línea discontinua; las jugadas de la naturaleza llevan sus probabilidades; y en cada nodo terminal figuran los pagos de todos los jugadores. 

\imagenfit{JE_forma_extensiva.pdf}{Forma extensiva del juego de entrada}{Elaboración propia.}

La figura lo muestra con el ejemplo. En el panel (a), el árbol del juego secuencial: el nodo inicial es de J1, los dos siguientes son de J2, cada rama es una acción y los nodos finales llevan los pagos de los dos jugadores. En el panel (b), el mismo juego con una jugada previa de la naturaleza, que elige el tipo de J2 con probabilidades $p$ y $1-p$, y una línea discontinua que une los dos nodos de J1, que no sabe qué tipo le ha tocado.

Recoge, por tanto, todo lo que fijan las reglas: quién mueve, cuándo, con qué opciones y sabiendo qué. Su definición formal exige siete elementos y supone que el número de nodos es finito, lo que excluye: las \textbf{acciones continuas} (una cantidad producida, una puja), los \textbf{juegos sin final} y los de \textbf{infinitos jugadores}.\footnote{%
  Los tres casos se tratan con ajustes menores de la definición [\authorfont{Ver Anexo \ref{anx:no_cooperativos}}]
} El árbol es la forma más común de representar un juego en forma extensiva, pero no la única.


\paragraph{Forma normal: ¿cuál es su utilidad?}
La \textbf{forma normal} o estratégica prescinde del orden de las jugadas y describe el juego con tres listas: los jugadores, el conjunto de estrategias de cada uno y el pago que obtiene cada jugador con cada perfil de estrategias. Con dos jugadores y pocas estrategias es la conocida \textbf{matriz de pagos}, con las estrategias de un jugador en las filas, las del otro en las columnas y, en cada casilla, los pagos de ambos. 

\imagenfit{JE_forma_normal.pdf}{Forma normal del juego de entrada}{Elaboración propia.}

La figura lo muestra con el ejemplo. En el panel (a), la forma normal de la versión simultánea: J1 elige fila, J2 elige columna, y cada casilla recoge los pagos de uno de los cuatro desenlaces. En el panel (b), la de la versión secuencial: las columnas ya no son acciones de J2, sino sus cuatro estrategias, con el precio tras la entrada en primer lugar, y por eso los desenlaces se repiten en varias casillas.

Es la representación natural de los juegos simultáneos, en los que cada jugador decide sin saber qué han elegido los otros, pero todo juego en forma extensiva tiene su forma normal: basta con listar las estrategias, que son planes completos. Al hacerlo se ve con claridad qué es una estrategia en cada tipo de juego, y se prepara el terreno para la idea de perfección en subjuegos.

\paragraph{Equivalencias: ¿son equivalentes ambas representaciones?}
La figura compara las dos representaciones con las dos variantes del ejemplo, que solo se diferencian en una línea discontinua. En el panel (a), J2 observa si J1 ha entrado: sus dos nodos son conjuntos de información distintos, cada estrategia suya dice qué precio fija tras la entrada y cuál tras la no entrada, y su forma normal tiene cuatro columnas, aunque en cada nodo solo elija entre dos acciones. Las letras dejan ver lo que hace la matriz con el árbol: cada casilla repite uno de sus cuatro desenlaces, porque dos estrategias de J2 que solo se diferencian en lo que haría en un nodo al que el juego no llega producen el mismo resultado. En el panel (b), J2 fija el precio sin saber si J1 ha entrado: sus dos nodos forman un único conjunto de información, sus estrategias se reducen a dos y la matriz, a dos columnas, una por desenlace. El juego es, a todos los efectos, simultáneo. La diferencia no está en lo que hacen los jugadores, sino en lo que saben al hacerlo.

\imagenfit{JE_ejemplo_formas.pdf}{Forma extensiva y forma normal del juego de entrada, secuencial y simultáneo}{Elaboración propia.}

La forma normal condensa el juego y omite detalles de la forma extensiva, de modo que cabe preguntarse si esa omisión importa o si la forma normal resume toda la información estratégicamente relevante. La pregunta puede plantearse así: \textit{¿es equivalente que los jugadores escriban de antemano sus planes completos y se los entreguen a un árbitro que los ejecute, a que jueguen secuencialmente tal y como describe la forma extensiva?} 

La cuestión ha sido objeto de controversia. KOHLBERG y MERTENS (1986) defendieron que toda la información estratégicamente relevante está en la forma normal y que un concepto de solución debería dar la misma respuesta en todas las formas extensivas que comparten forma normal. Los refinamientos basados en la racionalidad secuencial, de SELTEN (1965) en adelante, se definen en cambio sobre el árbol, porque es en él donde se ve qué haría cada jugador si llegara a un punto que el equilibrio no alcanza. En la práctica, los juegos dinámicos se analizan sobre la forma extensiva y los simultáneos, sobre la normal.


\subsection{Resolución: ¿cómo resolver un juego y cuáles son los conceptos de resolución?}

Una vez descrito el juego, la pregunta es cómo se jugará: en el ejemplo, si J1 entrará y qué precio fijará J2. Un \textbf{concepto de solución} es una regla que selecciona, en cada juego, los perfiles de estrategias que cabe esperar de jugadores racionales.

El más conocido es el \textbf{equilibrio}: un perfil de estrategias que representa un punto de estabilidad, porque ningún jugador tiene incentivos para cambiar su estrategia unilateralmente si los demás mantienen las suyas. La definición encierra dos ideas: que nadie quiere desviarse, porque si quisiera lo haría, y que cada jugador, al determinar su óptimo, tiene en cuenta las estrategias de todos los demás. No todos los conceptos de solución son equilibrios en este sentido, y la forma más clara de ordenarlos es preguntarse qué exige cada uno a las expectativas de los jugadores.

Los presentaremos con el mismo juego de entrada. Cada concepto resuelve un problema que deja abierto el anterior, y ese problema aparece al cambiar el orden de los pagos del ejemplo o lo que cada jugador observa al decidir.

\subsubsection{Mapa: ¿cuáles son los conceptos de equilibrio en los juegos no cooperativos?}
Las dos distinciones con que hemos descrito los juegos, el orden de las jugadas y la información sobre los demás, ordenan también sus conceptos de solución. GIBBONS (1992) las cruzó en un mapa de cuatro casillas que sirve de guía para todo lo que sigue: cada casilla tiene su concepto de equilibrio, y se pasa de una a otra cuando el concepto anterior deja de bastar.\footnote{%
  Otras clasificaciones completan el mapa. Según la suma de los pagos, en los \textbf{juegos de suma cero} lo que gana uno lo pierde el otro. 
  
  Fueron el punto de partida de la teoría, porque para ellos VON NEUMANN (1928) demostró el Teorema del minimax: si se admiten estrategias mixtas, el pago que un jugador puede garantizarse coincide con el máximo al que el otro puede limitarle. Pero en economía casi todas las situaciones son de \textbf{suma variable}: el comercio, la colusión o la coordinación generan ganancias o pérdidas conjuntas, que el minimax no capta, y de ahí la importancia de la generalización de NASH. Según el horizonte, un juego puede jugarse una sola vez o repetirse, un número finito de veces o indefinidamente; la repetición permite sostener con la amenaza de castigos futuros una cooperación que no se sostendría en un juego de una sola vez, y es la base del análisis de la colusión tácita ([\authorfont{Ver Tema 3.A.19}]).
} 

\imagenfit{JE_mapa.pdf}{Clases de juegos no cooperativos y sus conceptos de solución}{Elaboración propia a partir de GIBBONS (1992).}

Antes de llegar al equilibrio, la \textbf{dominancia} y las \textbf{estrategias racionalizables} permiten descartar jugadas sin exigir que nadie acierte lo que harán los demás: les basta la racionalidad de los jugadores y que esta sea de conocimiento común. Por eso predicen poco, y el primer equilibrio propiamente dicho da un paso más.

En los juegos estáticos con información completa, el concepto es el \textbf{equilibrio de NASH} (EN), propuesto por NASH (1950a, 1951): exige que las conjeturas de todos sean correctas y existe siempre, al menos en estrategias mixtas, si el juego es finito.\footnote{%
  Es el concepto con el que se resuelven, entre otros, el dilema del prisionero y los oligopolios de COURNOT y de BERTRAND ([\authorfont{Ver Tema 3.A.19}]).
} En los juegos dinámicos con información completa, el equilibrio de NASH admite amenazas que nadie cumpliría llegado el caso, y SELTEN (1965) lo refinó con el \textbf{equilibrio de NASH perfecto en subjuegos} (ENPS), que se halla por inducción hacia atrás.\footnote{%
  Con él se analizan el liderazgo de STACKELBERG y la disuasión de la entrada ([\authorfont{Ver Tema 3.A.19}]).
} En los juegos estáticos con información incompleta, HARSANYI (1967-1968) extendió el equilibrio de NASH al juego bayesiano: en el \textbf{equilibrio bayesiano de NASH} (EBN), una estrategia es un plan para cada tipo, y las conjeturas correctas se refieren a esos planes.\footnote{%
  Es la herramienta de las subastas y del diseño de mecanismos, que estudiamos en el bloque siguiente.
} Y en los juegos dinámicos con información incompleta, el \textbf{equilibrio bayesiano perfecto} (EBP) añade creencias sobre lo que el jugador no ha visto, actualizadas con la regla de Bayes, y se refina con el equilibrio secuencial de KREPS y WILSON (1982a), emparentado con el equilibrio perfecto de mano temblorosa de SELTEN (1975).\footnote{%
  Es el concepto de los juegos de señalización, como la educación como señal de la productividad (SPENCE, 1973; [\authorfont{Ver Tema 3.A.13}]).
}

Cada paso responde a un problema del anterior y lo contiene como caso particular, de modo que los conceptos no son alternativos, sino acumulativos: todo ENPS es un equilibrio de NASH, el EBN es el equilibrio de NASH del juego bayesiano y todo EBP es un EBN. El juego de entrada recorre el mapa entero: sus versiones simultánea y secuencial ocupan la fila de arriba, y su versión con tipos, la de abajo.

\subsubsection{Dominancia: ¿se trata de juegos en sentido teórico?}

El concepto menos exigente es la dominancia. La característica de los juegos es que representan situaciones estratégicas, en las que la acción de cada uno depende de lo que espera que hagan los otros. Pero hay juegos en los que la configuración de pagos es tal que esa condición estratégica prácticamente se disuelve, porque un jugador tiene buenas razones para preferir una estrategia a otra al margen de lo que haga el rival. Una estrategia \textbf{domina estrictamente} a otra si le proporciona un pago mayor haga lo que haga el resto de jugadores, es decir, frente a cualquier combinación de estrategias ajenas. Es una de las conclusiones primeras e importantes de la teoría: un jugador racional nunca juega una estrategia estrictamente dominada y, para descartarla, no necesita prever nada sobre los demás: le basta con ser racional.

\imagenfit{JE_res_dominancia.pdf}{Dominancia estricta en el juego de entrada simultáneo}{Elaboración propia.}

La figura lo muestra con la versión simultánea del ejemplo, en la que J2 fija el precio sin saber si J1 ha entrado, y con un J2 acomodaticio. El precio alto le da más que el bajo tanto si J1 entra ($a_2>b_2$) como si no entra ($c_2>d_2$): Alto domina estrictamente a Bajo, y J2 puede decidir sin preguntarse qué hará J1. J1, en cambio, no tiene ninguna estrategia que domine a la otra, porque entrar le conviene si J2 mantiene el precio alto ($a_1>c_1$), pero no si lo baja ($d_1>b_1$).

¿Sigue siendo esto un juego? Formalmente, sí: tiene jugadores, reglas, estrategias y pagos, y el pago de cada uno depende de lo que hagan los demás. Pero, para el jugador que tiene una estrategia dominante, la interdependencia deja de importar a la hora de decidir. A J2 no le da igual que J1 entre o no, porque sus pagos cambian, pero su mejor opción es la misma en los dos casos, y su problema se reduce al de una decisión individual. La interdependencia sigue en los pagos, pero desaparece de la decisión.

Por la misma razón, la dominancia es un concepto de solución, pero no un equilibrio en el sentido que acabamos de ver. No pide que las estrategias de unos y otros sean coherentes entre sí, sino solo que nadie juegue una estrategia que nunca le conviene: dice qué no se jugará, y solo en algunos juegos llega a decir qué se jugará.

Si una estrategia domina a todas las demás del jugador, es una \textbf{estrategia dominante}, y si cada jugador tiene una, el juego tiene un \textbf{equilibrio en estrategias dominantes}: la predicción más robusta que puede ofrecer la teoría, y la única en la que ningún jugador necesita anticipar lo que harán los demás. En el ejemplo, ocurriría si J1 fuese tan eficiente que ganase incluso con la guerra de precios ($b_1>d_1$): entrar dominaría entonces a no entrar, y (Entrar, Alto) sería un equilibrio en estrategias dominantes.\footnote{%
  Robusta, pero no necesariamente buena para los jugadores. El \textbf{dilema del prisionero} es el ejemplo clásico. Dos detenidos, interrogados por separado, pueden confesar o callar. Si ambos callan, cada uno obtiene $R$; si ambos confiesan, $P$; si solo uno confiesa, él obtiene $T$ y el otro, $S$, con $T>R>P>S$ (una condena más larga es un pago menor). Confesar es lo mejor para cada uno haga lo que haga el otro ($T>R$ y $P>S$), pero si ambos confiesan obtienen $P$, menos que el $R$ que tendrían si ambos hubieran callado. La racionalidad individual conduce a un resultado ineficiente en el sentido de Pareto, y esa tensión entre el interés individual y el colectivo es la que aparece en la colusión entre oligopolistas, en la provisión de bienes públicos o en las guerras comerciales.
}

Cuando no todos los jugadores tienen una estrategia dominante, todavía puede avanzarse por \textbf{eliminación iterada} de estrategias estrictamente dominadas: se descartan las dominadas, se examina el juego que queda y se repite. Aquí aparece por primera vez el conocimiento común de la racionalidad. Para eliminar en la primera ronda basta con que cada jugador sea racional; para eliminar en la segunda, hace falta además que sepa que los demás lo son y que, por tanto, no jugarán sus estrategias dominadas; en la tercera, que sepa que los demás saben que él lo es, y así sucesivamente. Cada ronda exige un escalón más de conocimiento mutuo.

\imagenfit{JE_res_iterada.pdf}{Eliminación iterada de estrategias estrictamente dominadas}{Elaboración propia.}

En el ejemplo bastan dos rondas. En la primera, J2, que es racional, descarta el precio bajo, estrictamente dominado. J1 no tiene ninguna estrategia dominada en el juego inicial, pero, si sabe que J2 es racional, sabe también que no bajará el precio, y en el juego que queda entrar le da más que no entrar ($a_1>c_1$). La solución es (Entrar, Alto), con pagos $(a_1,a_2)$. Es, además, un equilibrio de NASH, como veremos: ninguno de los dos querría cambiar de estrategia dada la del otro. El razonamiento es el mismo sobre el árbol: J2 elige en su conjunto de información la rama que le da más en sus dos nodos, y J1, que lo anticipa, compara solo los desenlaces de esa rama.

Con dominancia estricta, el orden en que se eliminan las estrategias no altera el resultado final: por dondequiera que se empiece, se llega a la misma solución. Este proceso de búsqueda de la solución por eliminación de estrategias dominadas se denomina \textbf{dominación repetida} (\textit{iterated dominance}), y los juegos que se resuelven así se llaman resolubles por dominancia. Naturalmente, no todos los juegos lo son: en muchos ninguna estrategia está dominada, y la dominancia no dice nada. Es lo que ocurre en el ejemplo si J2 es agresivo: le conviene bajar el precio si J1 entra ($b_2>a_2$) y mantenerlo si no entra ($c_2>d_2$), de modo que ninguna de sus estrategias domina a la otra, y tampoco las de J1. La formalización de la dominancia, incluida la dominancia por estrategias mixtas, se recoge en el anexo ([\authorfont{Ver Anexo \ref{anx:no_cooperativos}}]).

\paragraph{Dominancia débil: }
Una versión diluida del concepto de dominancia es la \textbf{dominancia débil}. Diremos que una estrategia domina débilmente a otra si nunca le proporciona un pago menor y, en algún caso, le proporciona uno mayor. La segunda condición importa: sin ella, dos estrategias que dieran siempre lo mismo se dominarían débilmente la una a la otra.

\imagenfit{JE_res_debil.pdf}{Dominancia débil en el juego de entrada secuencial}{Elaboración propia.}

La dominancia débil aparece de forma natural en los juegos secuenciales, porque en su forma normal muchas estrategias solo se diferencian en lo que harían en nodos que el juego no alcanza. En la versión secuencial del ejemplo, con J2 acomodaticio, la estrategia (Alto, Alto) da a J2 lo mismo que (Alto, Bajo) si J1 entra, y más si no entra ($c_2>d_2$); y lo mismo que (Bajo, Alto) si J1 no entra, y más si entra ($a_2>b_2$). Domina débilmente a las dos, y estrictamente a (Bajo, Bajo). Si J2 descarta las tres, a J1 le queda comparar $a_1$ con $c_1$, y entra. Entre las estrategias descartadas está (Bajo, Alto), la amenaza de guerra de precios: la dominancia débil anticipa así lo que veremos al tratar la credibilidad.

La dominancia débil es también la que justifica que, en la subasta a sobre cerrado de segundo precio, pujar la propia valoración sea la mejor opción, como veremos al tratar las subastas. Con ella, sin embargo, la eliminación iterada es menos fiable que con la dominancia estricta.\footnote{%
    \textbf{NOTA.-} La eliminación iterada de estrategias débilmente dominadas es más difícil de justificar, por dos razones. La primera es de coherencia: descartar una estrategia débilmente dominada solo es seguro si el jugador concede alguna probabilidad positiva a todas las estrategias del rival (si no, la estrategia dominada podría ser igual de buena), pero la lógica de la eliminación iterada consiste precisamente en suponer que algunas de esas estrategias no se jugarán (SAMUELSON, 1992). La segunda es práctica: el resultado puede depender del orden en que se eliminen las estrategias, cosa que no ocurre con la dominancia estricta.
}

\paragraph{Estrategias racionalizables: ¿qué son y por qué importan?}

Entre la dominancia y el equilibrio de NASH hay un paso intermedio que introduce la idea central de toda la teoría: la \textbf{mejor respuesta} (o respuesta óptima). La mejor respuesta es la estrategia que proporciona mejores resultados que todas las demás frente a una conjetura concreta sobre lo que harán los rivales. A diferencia de la dominancia, que compara estrategias frente a \textit{todo} lo que pueden hacer los rivales, la mejor respuesta se define frente a \textit{una} conjetura dada.

\imagenfit{JE_res_racionalizables.pdf}{Mejores respuestas y estrategias racionalizables}{Elaboración propia.}

En la figura, cada pago subrayado señala la mejor respuesta de un jugador a una estrategia del otro, y cada flecha, hacia dónde se movería quien no la está jugando. En la versión simultánea del ejemplo, la mejor respuesta de J1 es entrar si espera que J2 mantenga el precio alto ($a_1>c_1$) y no entrar si espera que lo baje ($d_1>b_1$). La de J2 depende de su tipo: si es acomodaticio, mantener el precio alto haga lo que haga J1; si es agresivo, bajarlo si J1 entra ($b_2>a_2$) y mantenerlo si no entra ($c_2>d_2$).

BERNHEIM (1984) y PEARCE (1984) llamaron \textbf{racionalizable} a una estrategia que es mejor respuesta a alguna conjetura razonable sobre los rivales, entendiendo por razonable que solo les atribuya estrategias que, a su vez, sean mejores respuestas a alguna conjetura razonable, y así sucesivamente. Es exactamente lo que se puede deducir de la racionalidad y de su conocimiento común, sin más supuestos. Una estrategia que no es mejor respuesta frente a ninguna conjetura nunca se juega.\footnote{%
  En los juegos de dos jugadores, las estrategias que sobreviven a ese descarte iterado son las mismas que sobreviven a la eliminación iterada de estrategias estrictamente dominadas ([\authorfont{Ver Anexo \ref{anx:no_cooperativos}}])
}

El ejemplo muestra por qué el concepto importa y dónde se queda corto. Con J2 acomodaticio, en el panel (a), el precio bajo no es mejor respuesta a nada, de modo que J1 no puede atribuírselo a un J2 racional: su única conjetura razonable es el precio alto, y su única estrategia racionalizable, entrar. Las flechas convergen en (Entrar, Alto), y la racionalizabilidad predice lo mismo que la eliminación iterada. Con J2 agresivo, en el panel (b), todas las estrategias son racionalizables, porque cada una es mejor respuesta a una conjetura que, a su vez, puede justificarse. J1 puede entrar porque espera el precio alto; esa conjetura es razonable, porque J2 mantendría el precio alto si esperase que J1 no entra; esta también lo es, porque J1 no entraría si esperase el precio bajo; y J2 bajaría el precio si esperase la entrada. La cadena se cierra sobre sí misma, como las flechas del panel (b), y deja todo abierto.

La racionalizabilidad exige, pues, que las conjeturas sean coherentes con la racionalidad de todos, pero no que sean \textit{correctas}: cada jugador puede jugar una mejor respuesta a algo que el rival no va a hacer. Si J1 entra porque espera el precio alto y J2 mantiene el precio alto porque espera que J1 no entre, los dos han razonado bien, pero J2 se ha equivocado. Por eso a menudo predice poco. Para predecir más hay que exigir más a las expectativas, y es lo que hace el equilibrio de NASH.

\subsubsection{Equilibrio de NASH: }

La dominancia es aplicable de manera restringida, porque en muchos juegos no hay estrategias dominadas, ni estricta ni débilmente, y la racionalizabilidad deja a menudo demasiado abierto. En ese caso, ¿cómo se juega el juego? ¿Qué se considera una solución razonable? La respuesta más general la dio NASH (1950a, 1951), que generalizó a cualquier juego finito la idea de equilibrio que VON NEUMANN y MORGENSTERN (1944) habían propuesto para un ámbito muy concreto, los juegos de suma cero.

Su noción de equilibrio parte de la mejor respuesta y añade una exigencia más: que las conjeturas de todos sean \textbf{correctas}. ¿Qué significa? Cada jugador decide a partir de una conjetura sobre lo que harán los demás. Para que una estrategia sea racionalizable basta con que esa conjetura sea defendible, es decir, que atribuya a los rivales estrategias que un jugador racional podría elegir. El equilibrio de NASH exige más: que la conjetura de cada jugador coincida con lo que los demás hacen de verdad. No es un requisito sobre cómo razona cada uno por separado, sino sobre la coherencia entre todos: cada jugador razona bien y, además, nadie se equivoca sobre los demás.

\imagenfit{JE_nash_definicion.pdf}{Equilibrio de NASH en estrategias puras: conjeturas que se confirman}{Elaboración propia.}

La figura lo muestra con la versión simultánea del ejemplo; como antes, cada pago subrayado señala la mejor respuesta de un jugador a una estrategia del otro. En el panel (a), con J2 acomodaticio, J1 espera que J2 mantenga el precio alto y entra; J2 espera que J1 entre y mantiene el precio alto. Las dos conjeturas se confirman: en la casilla (Entrar, Alto) están subrayados los dos pagos, y es un equilibrio de NASH. En el panel (b), con J2 agresivo, ninguna casilla tiene los dos pagos subrayados, de modo que no hay manera de que las dos conjeturas se confirmen a la vez. Si J1 entra porque espera el precio alto y J2 acierta y espera la entrada, J2 bajará el precio y J1 se habrá equivocado; si J1 acierta y espera el precio bajo, no entrará, y entonces a J2 le convendrá mantenerlo. El juego no tiene equilibrio en estrategias puras.

Un \textbf{equilibrio de NASH} (EN) es, simplemente, una combinación de estrategias en la que la de cada jugador es una mejor respuesta a las de los demás.\footnote{%
  Analíticamente, si $s^*=(s_i^*,s_{-i}^*)$ es el perfil de equilibrio, la condición de Equilibrio de NASH en estrategias puras es que ningún jugador pueda aumentar su pago cambiando él solo de estrategia:
  $$u_i(s_i^*,s_{-i}^*)\geq u_i(s_i,s_{-i}^*)\qquad\forall s_i\in S_i,\ \ \forall i$$
}

En esta situación, como todos los jugadores utilizan sus mejores respuestas, ninguno tiene razón alguna para cambiar de estrategia: quien se desviara estaría igual o peor. Por eso se dice que esa combinación de estrategias está en equilibrio, que es estable. El equilibrio es, además, autocumplido: si cada uno espera que los demás jueguen su parte del equilibrio, jugar la propia es lo mejor que puede hacer, y las expectativas se confirman. Nótese que la condición no se exige frente a toda estrategia de los demás, como en la dominancia, sino frente a la estrategia concreta que juegan en el equilibrio.

\paragraph{Equilibrio de Nash en estrategias mixtas}

Una \textbf{estrategia mixta} es una distribución de probabilidad sobre el conjunto de estrategias puras, que son las que se han visto hasta ahora: estrategias que consisten en elegir un curso de acción determinado. La idea no es en absoluto evidente: ¿por qué elegiría al azar un jugador racional? La respuesta está en los juegos en los que ser previsible se castiga, como el ejemplo con J2 agresivo. Si J2 mantuviera siempre el precio alto, J1 entraría, y entonces a J2 le convendría bajarlo; si lo bajara siempre, J1 no entraría, y entonces le convendría mantenerlo. Ninguna casilla es un equilibrio, porque quien sale perdiendo siempre querría cambiar.\footnote{%
  Es la estructura del juego de \textbf{cara o cruz} (\textit{matching pennies}), el ejemplo clásico de juego sin equilibrio en estrategias puras: dos jugadores muestran a la vez una moneda, uno gana si coinciden y el otro si no, y en el único equilibrio cada uno elige cara o cruz con probabilidad $\tfrac{1}{2}$ ([\authorfont{Ver Anexo \ref{anx:no_cooperativos}}]).
}

\imagenfit{JE_nash_mixtas.pdf}{Equilibrio de NASH en estrategias mixtas}{Elaboración propia.}

Supongamos, como en el panel (a) de la figura, que J1 entra con probabilidad $p$ y que J2 baja el precio con probabilidad $q$. Entrar le da a J1, en promedio, $(1-q)\,a_1+q\,b_1$, y no entrar, $c_1$, de modo que J1 solo está dispuesto a mezclar si $q$ es igual al umbral que encontramos al tratar la función de pagos. Del mismo modo, mantener el precio le da a J2, en promedio, $p\,a_2+(1-p)\,c_2$, y bajarlo, $p\,b_2+(1-p)\,d_2$, y J2 solo está dispuesto a mezclar si las dos cosas le dan lo mismo. Son las condiciones de indiferencia del equilibrio en estrategias mixtas de NASH (1950a), aplicadas a la versión simultánea del ejemplo con J2 agresivo:
$$q^*=\frac{a_1-c_1}{a_1-b_1}\qquad\qquad p^*=\frac{c_2-d_2}{(b_2-a_2)+(c_2-d_2)}$$
El panel (b) dibuja las mejores respuestas de cada uno. J1 entra si la probabilidad de guerra de precios es menor que $q^*$ y no entra si es mayor; J2 baja el precio si la probabilidad de entrada es mayor que $p^*$ y lo mantiene si es menor. Las dos líneas solo se cortan en $(q^*,p^*)$, y ese es el equilibrio de NASH del juego: J1 entra con probabilidad $p^*$ y J2 baja el precio con probabilidad $q^*$. Ninguno puede mejorar, porque, frente a la mezcla del otro, cualquier elección propia le da lo mismo.

El ejemplo revela la lógica de todo equilibrio en estrategias mixtas, el \textbf{principio de indiferencia}: un jugador solo está dispuesto a mezclar entre varias estrategias puras si todas le dan el mismo pago esperado, y esa indiferencia la produce la mezcla del rival. Cada jugador, con su estrategia mixta, deja al otro indiferente entre sus estrategias puras, y solo así el otro está dispuesto a mezclar. De ahí un resultado que sorprende a primera vista: la probabilidad con que mezcla cada jugador depende de los pagos del otro, no de los suyos. En el ejemplo, $q^*$ solo depende de los pagos de J1 y $p^*$, de los de J2: si la guerra de precios fuese más rentable para J2, J2 no bajaría el precio más a menudo; sería J1 quien entraría menos.\footnote{%
  Dos precisiones corrigen una lectura habitual. Frente a una estrategia mixta cualquiera no todas las estrategias son mejores respuestas: solo lo son frente a la mezcla concreta que deja indiferente al jugador. Y el equilibrio en estrategias mixtas no exige que mezclen todos: hay juegos con equilibrios en los que un jugador mezcla y otro juega una estrategia pura, aunque en los juegos de dos por dos con pagos genéricos los equilibrios en mixtas son completamente mixtos. La condición formal de indiferencia y el cálculo del ejemplo están en el anexo ([\authorfont{Ver Anexo \ref{anx:no_cooperativos}}]).
}

Las estrategias mixtas incluyen diferentes cursos de acción, cada uno con una probabilidad determinada, y amplían así el conjunto de estrategias entre las que puede elegir el jugador\footnote{%
    Las estrategias mixtas tienen una lectura inmediata, por ejemplo, en juegos en los que cabe combinar las estrategias puras en determinadas proporciones, como el de pares o nones; en juegos que se repiten muchas veces, en los que se puede elegir unas veces una estrategia y otras veces otra; o en juegos que permiten dividir lo que se apuesta y colocar una parte sobre una estrategia pura y el resto sobre otra.
}. Gracias a esa ampliación aparecen equilibrios en juegos que, examinados solo con estrategias puras, no los tenían. Al mezclar, el jugador elige probabilidades, y su pago pasa a ser la media de los pagos de las estrategias puras ponderada por esas probabilidades; de ahí que los pagos deban ser utilidades de VON NEUMANN-MORGENSTERN.

¿Qué significa en la realidad una estrategia mixta? Desde una interpretación no racional, sus probabilidades pueden entenderse simplemente como las frecuencias con que se han elegido en el pasado las estrategias puras, o como las proporciones de una población que juegan cada una; las estrategias mixtas serían solo regularidades estocásticas. La interpretación más interesante, debida a HARSANYI (1973), no obliga a abandonar el supuesto de racionalidad, pero exige considerar que la información no es completa. En el ejemplo, J2 no tiene por qué sortear nada: basta con que J1 no conozca con exactitud sus pagos, porque dependen de algo que solo J2 sabe, como sus costes. J2 conoce sus pagos y elige una estrategia pura, pero, para J1, que no tiene esa información, la elección de J2 se presenta como una estrategia mixta: con cierta probabilidad sus pagos serán unos y bajará el precio, y con la probabilidad restante serán otros y lo mantendrá. La estrategia mixta se traduce así en incertidumbre de un jugador sobre los verdaderos pagos, o sobre la información privada, del rival, que juega estrategias puras pero se le presenta como si mezclara. HARSANYI (1973) demostró que casi todo equilibrio en estrategias mixtas puede verse como el límite de equilibrios en estrategias puras de juegos con pagos ligeramente perturbados por información privada (la \textit{purificación}): es, como veremos, la lógica del equilibrio bayesiano de NASH.

\paragraph{Existencia: ¿siempre existe un Equilibrio de NASH?}

Al analizar un juego en forma normal se buscan todas las combinaciones de estrategias que son equilibrios de NASH, recorriendo las casillas y comprobando si las estrategias son simultáneamente mejores respuestas. Hay juegos con un único equilibrio, como el ejemplo simultáneo con J2 acomodaticio; juegos sin ninguno en estrategias puras, como el ejemplo con J2 agresivo; y juegos con varios.

\imagenfit{JE_nash_varios.pdf}{Varios equilibrios de NASH en el juego de entrada secuencial}{Elaboración propia.}

La figura muestra el tercer caso con la versión secuencial del ejemplo y un J2 acomodaticio. Los pagos subrayados vuelven a señalar las mejores respuestas, y en su forma normal hay tres equilibrios de NASH en estrategias puras, recuadrados. En (Entrar; Alto, Alto) y en (Entrar; Alto, Bajo), J1 entra y J2 mantiene el precio alto; las dos estrategias de J2 solo se diferencian en lo que haría si J1 no entrase, algo que no llega a ocurrir. En (No entrar; Bajo, Alto), J2 amenaza con bajar el precio si J1 entra, J1 se lo cree y no entra, y como la entrada no se produce, la amenaza no le cuesta nada a J2: su estrategia es una mejor respuesta a que J1 no entre. La multiplicidad es frecuente y no es un problema menor: si hay varios equilibrios, la teoría no predice cuál se jugará sin un criterio adicional, y de ahí nacen los refinamientos y las teorías de selección de equilibrios. En el ejemplo, el criterio que falta es la credibilidad de la amenaza.\footnote{%
  Tampoco hay que confundir la ausencia de estrategias dominadas con la ausencia de equilibrio. En la \textbf{batalla de los sexos}, en la que dos personas quieren pasar la tarde juntas pero prefieren planes distintos, ninguna estrategia está dominada y hay dos equilibrios en estrategias puras, uno en cada plan. En los juegos de dos jugadores con dos estrategias cada uno, que un jugador tenga una estrategia estrictamente dominante sí garantiza que haya un equilibrio en estrategias puras; la ausencia de estrategias dominadas es, por tanto, una condición necesaria para que no lo haya, pero no suficiente.
}

La teoría de juegos no puede conformarse con constatar la ausencia de equilibrio, porque eso supondría reconocer que es incapaz de predecir cómo se jugará un juego bien definido, cuando aspira a determinar en todo juego posible qué cuenta como elección racional. NASH (1950a) resolvió el problema al demostrar que todo juego \textbf{finito} (con un número finito de jugadores y de estrategias puras) tiene al menos un equilibrio de NASH si se admiten estrategias mixtas, como el que acabamos de hallar en el ejemplo con J2 agresivo.\footnote{%
  Con infinitas estrategias, como cuando las empresas eligen cantidades o precios en un continuo, existe un equilibrio en estrategias puras si los conjuntos de estrategias son compactos y convexos y el pago de cada jugador es continuo y cuasicóncavo en su propia estrategia (DEBREU, 1952; GLICKSBERG, 1952; FAN, 1952), condiciones que cumple, por ejemplo, el duopolio de COURNOT con demanda lineal; si fallan, el equilibrio puede no existir. La demostración del teorema de NASH y estas condiciones están en el anexo ([\authorfont{Ver Anexo \ref{anx:no_cooperativos}}]).
}

\paragraph{Interpretación del equilibrio de Nash}

Conviene subrayar que la noción de equilibrio de NASH no se basa en ninguna teoría sobre cómo los jugadores llegan al equilibrio. Lo único que establece es que, si los jugadores eligen estrategias que son mutuamente mejores respuestas, ninguno tendrá incentivos para cambiar; pero no dice nada sobre cómo llegan a seleccionarlas. Es la otra cara de su exigencia: pide que las conjeturas sean correctas, pero no explica de dónde procede esa corrección.

La literatura ha ofrecido varias razones para esperar que se juegue un equilibrio de NASH. Puede ser la consecuencia de la inferencia racional, si cada jugador es capaz de deducir lo que harán los demás, que es la interpretación racionalista del propio NASH (1950b). Puede ser también la condición necesaria de cualquier predicción única, porque un resultado que todos esperan y del que alguien quisiera desviarse no sería esperado; un \textbf{punto focal} que destaca por razones culturales o de contexto (SCHELLING, 1960); un \textbf{acuerdo que se cumple solo}, el único que pueden alcanzar jugadores que hablan antes de jugar pero no pueden comprometerse; o una \textbf{convención social estable}, el punto de reposo al que lleva la experiencia cuando el juego se repite, que es la lectura más cercana al uso tradicional del equilibrio en economía. Todas comparten una idea: cualquier resultado que vaya a jugarse de forma estable ha de ser un equilibrio de NASH, porque, si no lo fuera, alguien se desviaría. Ninguna dice, en cambio, cuál se jugará cuando hay varios.\footnote{%
  Las dos interpretaciones de NASH (1950b), las cinco razones que reúnen MAS-COLELL, WHINSTON y GREEN (1995) y por qué la convención social es la más cercana a la economía se desarrollan en el anexo ([\authorfont{Ver Anexo \ref{anx:no_cooperativos}}]).
}

En el ejemplo secuencial, ninguna de estas razones explica por sí sola por qué esperar que J1 entre y no que se quede fuera ante la amenaza de J2. Para elegir entre esos equilibrios hay que exigir algo más a las conjeturas: que no descansen en amenazas que nadie cumpliría llegado el caso.


\subsubsection{Credibilidad: ¿qué es el ENPS y qué completa?}

En los juegos dinámicos, el equilibrio de NASH no basta. Lo vimos con la versión secuencial del ejemplo: con J2 acomodaticio, su forma normal tiene tres equilibrios de NASH, y en uno de ellos, (No entrar; Bajo, Alto), J2 amenaza con una guerra de precios y J1, que se lo cree, no entra.

\imagenfit{JE_enps.pdf}{Credibilidad e inducción hacia atrás en el juego de entrada secuencial}{Elaboración propia.}

Ese perfil cumple la definición porque, como J1 no entra, la amenaza nunca se pone a prueba y no le cuesta nada a J2: el equilibrio de NASH solo exige que las estrategias sean óptimas en la senda que efectivamente se recorre, y deja libre lo que se haría fuera de ella. Pero la amenaza no es \textbf{creíble}: si J1 entrara, a J2 le convendría acomodarse ($a_2>b_2$). Por tanto, un concepto de solución que admite amenazas que ningún jugador racional cumpliría predice mal.

La idea de credibilidad es una de las aportaciones más fecundas de la teoría de juegos a la economía. Explica el \textbf{valor del compromiso}: un jugador puede salir ganando si limita sus propias opciones, porque así hace creíble una amenaza que de otro modo no lo sería, como el incumbente que invierte en capacidad antes de que llegue el entrante o el líder que se compromete con su cantidad ([\authorfont{Ver Tema 3.A.19}]). En el ejemplo, a J2 le convendría comprometerse de antemano con la guerra de precios, por ejemplo con una capacidad que hiciera rentable bajar el precio: se comportaría como si fuera agresivo, y J1 no entraría.\footnote{%
  SCHELLING (1960) lo ilustró con el general que quema los puentes tras de sí. La misma lógica, trasladada a la política económica, es la de la \textbf{inconsistencia temporal}: el plan óptimo anunciado hoy deja de serlo mañana, una vez que los agentes han actuado, y los agentes, que lo anticipan, no se lo creen (KYDLAND y PRESCOTT, 1977). De ahí el valor de las reglas y de instituciones como la independencia del banco central ([\authorfont{Ver Tema 3.A.36}]).
}

Para resolver esta situación, SELTEN (1965) propuso exigir el \textbf{principio de racionalidad secuencial}: las estrategias de equilibrio deben especificar una conducta óptima desde cualquier punto del juego en adelante, incluidos los que no se alcanzan. Lo formalizó con la noción de \textbf{subjuego}, la parte del árbol que empieza en un nodo que forma por sí solo un conjunto de información y contiene todos sus sucesores, sin partir ningún conjunto de información; visto aisladamente, un subjuego es un juego en sí mismo. Un perfil de estrategias es un \textbf{equilibrio de Nash perfecto en subjuegos} (ENPS) si induce un equilibrio de Nash en cada subjuego. Como el propio juego es un subjuego, todo ENPS es un equilibrio de NASH; lo contrario no es cierto. En el ejemplo hay tres subjuegos: el juego entero y los dos que empiezan en los nodos de J2, recuadrados en la figura. En el que sigue a la entrada, el único equilibrio es acomodarse; en el que sigue a la no entrada, mantener el precio alto. Eso elimina dos de los tres equilibrios de NASH: el de la amenaza, que prescribe bajar el precio tras la entrada, y (Entrar; Alto, Bajo), que prescribe bajarlo sin competidor, algo que J2 nunca haría si llegara a ese nodo.

En los juegos finitos con información perfecta, el ENPS se halla por \textbf{inducción hacia atrás} (retroinducción): se resuelve la última decisión, se sustituye por su resultado y se retrocede hasta el principio. El panel (a) de la figura lo muestra: J2 elige en cada nodo la rama que más le da, en trazo grueso, y J1, que lo anticipa, compara $a_1$ con $c_1$ y entra. El ENPS es (Entrar; Alto, Alto). El panel (b) muestra que la conclusión depende del orden de los pagos. Si J2 es agresivo ($b_2>a_2$), la guerra de precios es su mejor respuesta a la entrada, y J1, que compara $b_1$ con $c_1$, no entra: el ENPS es (No entrar; Bajo, Alto). La misma estrategia de J2 que antes era una amenaza vacía es ahora una amenaza creíble.\footnote{%
  El resultado, que se remonta a ZERMELO (1913) y que KUHN (1953) generalizó, es que todo juego finito de información perfecta tiene un ENPS en estrategias puras, que es único si ningún jugador obtiene el mismo pago en dos nodos terminales distintos. Con información imperfecta, el procedimiento se generaliza sustituyendo cada subjuego final por los pagos de uno de sus equilibrios de NASH. La formalización y otros ejemplos están en el anexo ([\authorfont{Ver Anexo \ref{anx:no_cooperativos}}]).
}

En los juegos de horizonte infinito no hay una última decisión desde la que empezar y la inducción hacia atrás no es aplicable, de modo que el análisis se apoya en otros métodos, como comprobar que ningún jugador gana desviándose en un solo periodo. El caso más importante es el de los \textbf{juegos repetidos}.\footnote{%
  Si un juego con un único equilibrio de NASH se repite un número finito de veces, la inducción hacia atrás implica que el único ENPS consiste en jugar ese equilibrio en cada etapa: en el dilema del prisionero repetido diez veces, los jugadores confiesan siempre, y en el juego de entrada repetido en veinte mercados, el incumbente se acomoda siempre, por más que la intuición sugiera que le convendría luchar al principio para ganarse una reputación (la paradoja de la cadena de tiendas de SELTEN, 1978).
} ¿Por qué? Porque si el juego se repite indefinidamente y los jugadores valoran lo suficiente el futuro, la amenaza de castigar cualquier desviación con la vuelta permanente al equilibrio de la etapa puede sostener la cooperación como ENPS (FRIEDMAN, 1971).\footnote{%
    Esta es la base de la teoría de la colusión tácita [\authorfont{Ver Tema 3.A.19}].
  
    \textbf{NOTA.-} La inducción hacia atrás exige que los jugadores sigan suponiendo que los demás son racionales incluso después de observar jugadas que lo contradicen. En el juego del ciempiés (ROSENTHAL, 1981), en el que dos jugadores pueden ir pasándose un bote que crece o quedárselo, predice que el primero se lo quedará de inmediato, pero en los experimentos de MCKELVEY y PALFREY (1992) casi nadie lo hace. Es una de las críticas más citadas a la racionalidad secuencial.
}.

\subsubsection{Información privada: el equilibrio bayesiano de Nash}

Cuando la información es incompleta, el equilibrio de NASH no puede aplicarse tal cual, porque un jugador no conoce los pagos del rival y, por tanto, no sabe qué le conviene a este. La transformación de HARSANYI (1967-1968), que vimos al tratar la información, convierte el problema en un juego bayesiano, y su concepto de solución es el \textbf{equilibrio bayesiano de NASH} (EBN): un perfil de estrategias, que especifican qué hace cada tipo de cada jugador, en el que cada tipo elige la mejor respuesta a las estrategias de los demás, ponderando lo que haría cada tipo ajeno por su probabilidad. Las conjeturas siguen teniendo que ser correctas, como en el equilibrio de NASH, pero ahora son conjeturas sobre la \textit{regla} con que decide cada tipo del rival, no sobre una acción concreta, y se combinan con una distribución de probabilidad de los tipos que todos conocen. En rigor, el EBN es el equilibrio de NASH del juego bayesiano, en el que cada tipo se trata como si fuera un jugador distinto.

\imagenfit{JE_ebn.pdf}{Equilibrio bayesiano de NASH en el juego de entrada con información incompleta}{Elaboración propia a partir de HARSANYI (1967-1968).}

La figura recoge la versión simultánea del ejemplo con información incompleta. J1 no sabe si J2 es acomodaticio, con probabilidad $p$, o agresivo, con probabilidad $1-p$; J2 conoce su tipo, pero fija el precio sin saber si J1 ha entrado. Supongamos que, con costes bajos, a J2 le conviene el precio bajo tanto si J1 entra como si no ($\hat{b}_2>\hat{a}_2$ y $\hat{d}_2>\hat{c}_2$). El panel (a) muestra las dos matrices entre las que J1 no puede distinguir; el panel (b), el mismo juego tras la transformación de HARSANYI.

El juego se resuelve en dos pasos. J2, que conoce su tipo, tiene una estrategia dominante en cada caso, subrayada en el panel (a): mantener el precio alto si es acomodaticio ($a_2>b_2$ y $c_2>d_2$) y bajarlo si es agresivo. Su estrategia es, por tanto, una regla: Alto si es de tipo $\theta_{ac}$ y Bajo si es de tipo $\theta_{ag}$. J1 no sabe de qué tipo es J2, pero conoce esa regla y la probabilidad de cada tipo, de modo que espera el precio alto con probabilidad $p$ y el bajo con probabilidad $1-p$. Entrar le da en promedio $p\,a_1+(1-p)\,b_1$, y le conviene si ese pago supera a $c_1$. Es la condición del equilibrio bayesiano de NASH (HARSANYI, 1967-1968) en la versión simultánea del ejemplo: J1 entra si la probabilidad de que J2 sea agresivo es menor que el umbral que ya conocemos,
$$1-p<q^*=\frac{a_1-c_1}{a_1-b_1}$$
El equilibrio bayesiano de NASH es, pues, que J2 siga esa regla y que J1 entre si $1-p<q^*$ y no entre si $1-p>q^*$. Lo que hace J1 depende de su creencia sobre el tipo de J2, y el umbral, de los pagos: cuanto más le cueste a J1 la guerra de precios (cuanto menor sea $b_1$), menos probable ha de ser un incumbente agresivo para que se arriesgue a entrar. La probabilidad de guerra que, al tratar la función de pagos, era una conjetura de J1 sale ahora de la distribución de los tipos de J2, que todos conocen.\footnote{%
  El EBN es la herramienta con la que estudiaremos las subastas a sobre cerrado y el diseño de mecanismos, que busca reglas de juego cuyos equilibrios bayesianos produzcan el resultado deseado aunque el diseñador no conozca la información privada de los participantes. Además, es la lógica que permite reinterpretar las estrategias mixtas como incertidumbre sobre el tipo del rival (HARSANYI, 1973). Su definición formal está en el anexo ([\authorfont{Ver Anexo \ref{anx:no_cooperativos}}]).
}


\subsubsection{Creencias: el equilibrio bayesiano perfecto}

Cuando el juego es dinámico y la información es imperfecta o incompleta, la perfección en subjuegos puede no servir de nada.\footnote{%
  Si un jugador no observa la jugada anterior, o si la naturaleza ha elegido al principio un tipo que no todos ven, puede que el único subjuego sea el juego entero, y entonces todo equilibrio de Nash es perfecto en subjuegos.

  El anexo lo formaliza con la misma variante del juego de entrada que usamos a continuación ([\authorfont{Ver Anexo \ref{anx:no_cooperativos}}]).
} ¿Por qué? Esencialmente, porque para aplicar la racionalidad secuencial hace falta saber qué cree cada jugador sobre lo que ha ocurrido cuando le toca mover sin haberlo visto.

\paragraph{Creencias: ¿por qué no basta la perfección en subjuegos?}

\imagenfit{JE_ebp_creencias.pdf}{Creencias y racionalidad secuencial: entrada de dos maneras}{Elaboración propia a partir de MAS-COLELL, WHINSTON y GREEN (1995).}

Una primera variante del ejemplo lo muestra. Supongamos que J1 puede entrar de dos maneras, imitando el producto de J2 o diferenciándose de él, y que J2 ve que J1 ha entrado, pero no de qué manera; si J1 no entra, J2 sigue como monopolista y el desenlace es $(c_1,c_2)$. Los dos desenlaces nuevos, tras diferenciarse, llevan las letras $e$ y $f$. Los dos nodos de J2 forman un único conjunto de información, de modo que el único subjuego es el juego entero, y la perfección en subjuegos no descarta nada. En el panel (a) de la figura, el perfil (No entrar, Bajo) es un equilibrio de NASH y, por tanto, perfecto en subjuegos: si J1 cree en la guerra de precios, no entra ($c_1>b_1$ y $c_1>f_1$), y como no entra, a J2 no le cuesta nada amenazar. Pero si J2 prefiere acomodarse sea cual sea la manera en que J1 haya entrado ($a_2>b_2$ y $e_2>f_2$), la amenaza no es creíble, aunque no haya un subjuego en el que comprobarlo.

En contextos como este, la solución pasa por hacer explícitas las \textbf{creencias}: en cada conjunto de información con varios nodos, el jugador que mueve asigna una probabilidad a cada uno de ellos, es decir, a cada historia que ha podido conducir hasta allí. Un \textbf{equilibrio bayesiano perfecto} combina dos exigencias. La primera es la racionalidad secuencial: en cada conjunto de información, la acción elegida debe ser óptima dadas las creencias del jugador y las estrategias de los demás. La segunda es la coherencia de las creencias: deben deducirse de las estrategias de equilibrio con la regla de Bayes siempre que sea posible. \textit{Siempre que sea posible} es la clave. En los conjuntos de información que se alcanzan con probabilidad positiva, la regla de Bayes fija las creencias; en los que el equilibrio no alcanza, la probabilidad de llegar a ellos es nula y la regla no dice nada.\footnote{%
  El \textbf{equilibrio bayesiano perfecto débil} deja libres las creencias fuera de la senda de equilibrio; el equilibrio bayesiano perfecto les impone restricciones adicionales, y el \textbf{equilibrio secuencial} de KREPS y WILSON (1982a) exige que puedan justificarse como el límite de lo que se creería si los jugadores cometieran errores con una probabilidad muy pequeña, una idea emparentada con el equilibrio perfecto de mano temblorosa de SELTEN (1975). Las tres definiciones están en el anexo ([\authorfont{Ver Anexo \ref{anx:no_cooperativos}}]).
}

En el ejemplo, J2 asigna una probabilidad $\mu$ a que J1 haya imitado y $1-\mu$ a que se haya diferenciado. Mantener el precio le da $\mu\,a_2+(1-\mu)\,e_2$ y bajarlo, $\mu\,b_2+(1-\mu)\,f_2$; como acomodarse es mejor en los dos nodos, lo es para cualquier creencia. En el panel (b), la racionalidad secuencial obliga a J2 a acomodarse, y J1, que lo anticipa, entra de la manera que más le convenga: la amenaza queda descartada. Si J2 prefiriera la guerra frente a un imitador ($b_2>a_2$), lo óptimo dependería de $\mu$, y la coherencia de las creencias pasaría a importar: $\mu$ tendría que deducirse, con la regla de Bayes, de la manera en que J1 entra en equilibrio.

\paragraph{Señalización: ¿qué aprende un jugador de lo que hace el otro?}

La segunda variante cambia el orden del juego. J2, que conoce su coste, fija su precio antes de que J1 decida, y J1 ve el precio, pero no el coste: es el modelo de precio límite de MILGROM y ROBERTS (1982a). Ahora la jugada de J2 puede revelar su tipo, y J1 ha de interpretarla.

\imagenfit{JE_ebp_senales.pdf}{Señalización en el juego de entrada: equilibrio separador y agrupador}{Elaboración propia a partir de MILGROM y ROBERTS (1982a).}

En la figura, la naturaleza (N) elige el tipo de J2 y J2 fija el precio. J1 tiene dos conjuntos de información, tras el precio alto y tras el bajo, y en cada uno una creencia sobre la probabilidad de que J2 sea agresivo, $\mu_A$ y $\mu_B$; en cada uno decide si entra (E) o no entra (N). Como lo que interesa es la decisión de J1, en los nodos finales figuran solo sus pagos: entrar le da $a_1$ frente a un J2 acomodaticio, que se acomodará, y $b_1$ frente a uno agresivo, que desatará la guerra; no entrar le da $c_1$.

La regla de Bayes ata esas creencias a la estrategia de J2. En un equilibrio \textbf{separador}, como el del panel (a), cada tipo fija un precio distinto: el agresivo, el bajo, y el acomodaticio, el alto. El precio revela el tipo ($\mu_A=0$ y $\mu_B=1$), y J1 entra solo tras el precio alto. Para que se sostenga, al tipo acomodaticio no le ha de compensar imitar al agresivo: lo que ganaría disuadiendo la entrada ha de ser menor que lo que le cuesta bajar el precio. En un equilibrio \textbf{agrupador}, como el del panel (b), los dos tipos fijan el precio bajo, y J1 no aprende nada: su creencia sigue siendo la de partida ($\mu_B=1-p$), y decide como en el equilibrio bayesiano de NASH, de modo que no entra si $1-p>q^*$. El precio alto no se observa en ese equilibrio, y la regla de Bayes no dice nada sobre $\mu_A$: la creencia tras él es libre, y se elige para que ningún tipo quiera desviarse; en el panel (b), J1 entraría tras un precio alto. De ahí que estos juegos tengan muchos equilibrios, que criterios como el intuitivo de CHO y KREPS (1987) tratan de reducir descartando las creencias poco razonables.\footnote{%
  Es el concepto con el que se analizan los \textbf{juegos de señalización}, en los que un jugador informado actúa primero y el no informado infiere su tipo de lo que observa: además del precio límite, la educación como señal de la productividad (SPENCE, 1973; [\authorfont{Ver Tema 3.A.13}]) o la lucha temprana con que se construye una reputación de duro, que resuelve la paradoja de la cadena de tiendas cuando hay una pequeña probabilidad de que el incumbente sea de un tipo que siempre lucha (KREPS y WILSON, 1982b; MILGROM y ROBERTS, 1982b).
}

El recorrido ilustra la tensión con la que empezaba el bloque. Cada refinamiento elimina predicciones poco razonables a cambio de exigir más a la racionalidad y a las expectativas de los jugadores, y ninguno acaba del todo con la multiplicidad de equilibrios. La evidencia experimental, además, muestra que los jugadores reales se apartan de las predicciones con más frecuencia cuanto más larga es la cadena de razonamiento que exige el concepto. Por eso el equilibrio de Nash y sus refinamientos se entienden mejor como condiciones de coherencia que debe cumplir cualquier predicción razonable que como una descripción literal de cómo razonan las personas.











\clearpage
\section{Diseño de Mecanismos: Subastas}
\puntosclave{
    Una de las aplicaciones más sobresalientes de la teoría de juegos en la microeconomía es la teoría del diseño de mecanismos, con gran repercusión en la economía del bienestar.

    Las subastas se pueden contemplar como un tipo de mecanismos de esta teoría en el que los pujadores tienen información incompleta acerca de las valoraciones de los demás pujadores acerca del bien subastado. Por ello, para el estudio de las subastas, se acude al empleo del concepto de equilibrio Bayesiano de Nash.

    La taxonomía básica de los modelos teóricos sobre subastas distingue entre un modelo básico de valor privado, un modelo básico de valor común, y un modelo general que comprende los anteriores.

    Los tipos de subasta \textit{estándar} son (4) la ascendente, la descendente, la de primer precio a sobre cerrado, y la de segundo precio a sobre cerrado.

    En subastas de un único objeto existe una equivalencia entre la subasta descendente y la subasta de primer precio a sobre cerrado. En cuanto a la subasta ascendente y la subasta de segundo precio a sobre cerrado, también son equivalentes bajo muchos supuestos (y tienen propiedades similares, en general).

    El principal resultado de la teoría de subastas es el teorema de la equivalencia de ingresos. De hecho, gran parte de la teoría de subastas puede entenderse en términos de este teorema y de cómo sus resultados se ven afectados por la relajación de sus supuestos.
}

\textcolor{magenta}{En este apartado se aborda una de las aplicaciones más sobresalientes de la teoría de juegos: el diseño de mecanismos.

La teoría de diseño de mecanismos pretende permitir al usuario desarrollar o construir \textbf{instituciones}, conjuntos de reglas o regulaciones que gobiernan las interacciones entre los participantes de un sistema, de forma que se consigan los objetivos perseguidos. Por tanto, nos centraremos en aplicar los elementos de la teoría presentados con el fin de crear un juego.

Intuitivamente podemos pensar que esta rama de la Teoría de Juegos tiene el objetivo de permitirnos crear una máquina que utiliza como inputs información privada y cuyo output es el objetivo que nos hemos marcado, por ejemplo, podría ser la equidad, o cualquier otra cosa (máximo revenue, eficiencia, mínimo gasto, etc.). 

Por lo tanto, esta máquina que diseñaremos será en realidad un marco normativo (número de jugadores, estrategias, conjuntos de información, etc.) que, dadas las preferencias desconocidas a priori, la información privada de los jugadores, nos conduzca a un objetivo deseado mediante las técnicas tradicionales de resolución de juegos que hemos estudiado. 

El quid de la cuestión será precisamente encontrar el sistema de resultados óptimo dada esa información privada que desconocemos (podemos tener una idea abstracta, una estimación, etc.) que responde precisamente a las preferencias de los jugadores. Por lo tanto, el objetivo del diseñador del mecanismo será diseñar esta máquina de manera que, independientemente de cuáles sean las preferencias o deseos de los jugadores, el objetivo marcado se consiga. La importancia de esto es clave ya que precisamente, el objetivo depende de la información privada de los participantes, por tanto el planificador debe ser capaz de diseñar un mecanismo tal que el objetivo se cumpla independientemente de las preferencias de los jugadores.

Ejemplos de aplicación de esta teoría son el diseño de mecanismos de votación, la redacción de contratos entre partes que tendrán información privada y la elaboración de procedimientos para decidir acerca de proyectos públicos o estándares medioambientales (estos temas son objeto de estudio en la parte del temario dedicada a la economía del bienestar)\footnote{%
    En este sentido, el enfoque de diseño de mecanismos considera los mercados competitivos como una mera institución particular dentro de una familia más amplia de mecanismos dirigidos por planificadores benevolentes.
}.

Otro campo natural en el que se aplica la metodología del diseño de mecanismos es la teoría de subastas. En este sentido, se pueden definir las subastas como mecanismos mediante los cuales un principal trata de emplear la competencia entre distintos agentes para reducir las rentas de información que tiene que ceder al agente con el cual está negociando. Ello requiere una modelización de la relación entre los pujadores (los agentes) que pujan bajo información imperfecta acerca de las valoraciones de los demás pujadores acerca del bien subastado (o contrato). Por ello, la clarificación del concepto de Equilibrio Bayesiano de Nash por parte de HARSANYI (1967-68) fue el detonante del desarrollo posterior que vivió la teoría de las subastas.

Con arreglo a este enfoque, la estrategia de cada jugador es una función de su información y maximiza su pago esperado dadas las estrategias de los demás jugadores y dadas sus creencias acerca de la información de los demás jugadores.

El mecanismo nos permitirá revelar las preferencias de los jugadores, de manera que, si estos juegan el Juego revelarán indirectamente sus preferencias mediante el desarrollo del Juego.

(Repensar, reescribir: mucha repetición, poca coherencia y rigor, escritura muy confusa. Didáctico, ligado con lo anterior y objetivo del presente bloque: presentar el diseño de mecanismos (desde una perspectiva concreta utilizando los conceptos presentados) como paso previo a la aplicación práctica en el ámbito de las subastas)}

\subsection{Estructura General de un Mecanismo de Diseño}

\textcolor{magenta}{(Repensar, reescribir y reorganizar el contenido de este apartado. Objetivo: presentar la teoría del diseño de mecanismos (de manera concreta y no muy general) como vertiente que utiliza los conceptos presentados para conseguir una cosa. 

Además, libertad total para transferir todo el actual contenido a el anexo correspondiente así como para repensar, reorganizar, redesarrollar y reescribir todo el anexo: Objetivo: presentar formalmente la teoría del diseño de mecanismos, sus conceptos, su desarrollo, su vertientes, extensiones, limitaciones, aplicaciones, etc. (Mejor eso antes que esto, puesto que de allí se recuperará lo clave para construir este apartado, cuyo objetivo es mucho más concreto: solo presentación como previo en abstracto para la aplicación concreta a subastas) )}


\subsection{Subastas: ¿qué es y cómo diseñarlas?}
\textcolor{magenta}{
  
(Este apartado presenta dos graves problemas:
  1. Desde unaperspectiva formal: El desarrollo es muy deficiente. Muchas omisiones, faltas de desarrollo, incoherencias (se presentan tres y luego se explican dos, etc.). Quiero que tdoso el desarrollo formal de la teoría de las Subastas lo incluyas en el anexo del diseño de mecanismos una vez hayas desarrollado TODO lo relativo al diseño de mecanismos en abstracto. Dale profundidad y rigor (didáctico, ilustrativo, analítico y riguroso). 
  2. Desde una perspectiva expositiva: Todo el apartado tiene un enfoque manifiestamente deficiente. Quiero repensar, redesarrollar y reescribir todo el apartado con el objetivo de poder dar a la exposición un tono didáctico y sencillo. He pensado en lo siguiente: 
    3.1. Subastas: ¿qué son y cómo podemos diseñarlas?. Aquí incluir la introducción de la teoría de las subastas, ligado con lo expuesto en el apartado anterior. Presentar sus objetivos y la taxonomía de subastas. 
    3.2. Primer precio sobre cerrado
    3.3. Segundo precio sobre cerrado. 
    Estas dos últimos apartados deben presentarse con un enfoque eminentemente aplicado, como hicimos en el bloque anterior. La idea no es presentar un ejemplo, no, es mantener el rigor pero siguiendo un hilo basado en un ejemplo práctico que permita exponer los principales resultados, desarrollos, conceptos e implicaciones.
    
He pensado que, para la materialización del 3.2 y 3.3. podría ser una buena idea considerar la emisión de deuda por el tesoro, se trata en otras partes del temario por lo que es escalable, y el Tribunal estará familiarizado con ello. Lo importante es mantener el balance entre ejemplo y concepto que le dimos en el bloque anterior. Evidentemente, la idea es plantearlo de forma gráfica para evitar el uso intensivo de instrumental matemático que es más tedioso de escribir en la pizarra. Si crees que esto no es recomendable, cambiamos enfoque, si crees que sí, seguimos con ello: libertad para criticar. 

¿Sería conviente preservar el 3.4 teorema de la equivalencia?)



A pesar de las diferentes clasificaciones que a continuación presentaremos, las subastas comparten un marco común, en tanto necesitan dos componentes principales: (1) una regla de decisión: ¿quién gana?; y, (2) una regla de pago: ¿cuánto pagaré?

La primera clasificación, relativa a la información que cada Jugador conoce de las estrategias de sus rivales y sobre el desarrollo del Juego, podemos hacerla como sigue:

\begin{la}
    \item Subastas de información abierta:
    \begin{itemize}
        \item Subasta de precio ascendente (\textit{english auctions}). La subasta empieza a partir de un precio de reserva (normalmente) y los participantes aumentan el precio de manera sucesiva hasta que solo queda uno;
        \item Subasta de precio descendente (\textit{dutch auctions}). La subasta empieza a partir de un precio elevado y desciende paulatinamente hasta que algún participante acepte el precio propuesto, en cuyo caso se lo llevará.
    \end{itemize}
    \item Subastas de información privada o sobre cerrado:
    \begin{itemize}
        \item Subastas a primer precio. Se trata de un juego simultáneo en el que todos los participantes realizan su puja, en sobre cerrado, y gana aquel que ofrezca mejor precio.
        \item Subastas a segundo precio (Vickrey). También consiste en un juego simultáneo en el que todos los participantes realizan su puja de forma secreta y gana aquel que ofrezca mejor precio pero pagará el precio que haya estipulado aquel que haya quedado en segundo lugar (es decir, el primero por arriba o el último por abajo, de entre los perdedores).
        \item Subastas a tercer precio. Muy similar al anterior, el ganador es el mejor ofertante pero paga el precio del tercero por arriba (el segundo de los perdedores).
        \item Subastas a todo-pago (\textit{all-pay auctions}). Todo el mundo paga, independientemente de si se gana o se pierde, solo el mejor ofertante es el que recibe el objeto.
    \end{itemize}
\end{la}

Además, podemos distinguir también tres tipos principales de marcos teóricos, en función de la información sobre el objeto de la subasta que cada participante tiene:

\begin{lnum}
    \item \textit{Modelo básico de valor privado}. Cada pujador conoce en cuánto valora el/los objeto/s a la venta, pero su valor es información privada suya;
    \item \textit{Modelo de valor común puro}. El valor es el mismo para todos los pujadores, pero los pujadores tienen información privada distinta acerca de cuál es dicho valor. Por ejemplo, el valor del arrendamiento para explotar petróleo en un terreno depende de cuánto petróleo albergue y los pujadores pueden tener acceso a distintas señales geológicas acerca de dicha cantidad. En este caso un pujador modificaría su estimación del valor si conociese la señal de otro pujador (en contraposición al caso de valor privado, en el que su valor no se vería afectado por conocer las preferencias o información de cualquier otro pujador)\footnote{%
    Un aspecto clave en las pujas de subastas con componentes de valores comunes es la maldición del ganador: todo pujador debe ser consciente de que se adjudica el objeto solo si tiene la señal más alta (en equilibrio simétrico). El hecho de no tener en cuenta las malas noticias acerca de las señales de los demás que conlleva toda victoria, conducirá a que el adjudicatario pague, de media, más de lo que vale el objeto, y se dice que ello a menudo ocurre en la práctica (el artículo empírico clásico sobre esta maldición se debe a CAPEN, CLAPP y CAMPBELL (1971) en relación con las subastas de derechos de explotación petrolífera en el Golfo de México). En equilibrio, los pujadores deben ajustar hacia abajo sus pujas de modo acorde.
}.
\end{lnum}

Algunos marcos comparten elementos de ambos, por ejemplo, las subastas en el mercado inmobiliario contienen tanto información objetiva como información subjetiva relativa al objeto de la subasta (todos comparten un mínimo de información común atendiendo a las características del inmueble, la situación, la localización, etc.). Por ello, a veces se distingue como otro tipo diferente: \textit{Modelo general}. Supone que cada pujador recibe una señal de información privada, pero permite que el valor de cada pujador sea una función general de todas las señales. Otro ejemplo sería el valor que se otorga a una pintura, que puede depender de la información privada propia (cuánto nos gusta) pero también en cierto modo de la información privada de los demás (cuánto gusta a los demás) porque ello afecta al valor de reventa y/o el prestigio que confiere ser su propietario.

\subsubsection{Modelo general: Valor privado independiente}

Podemos ver la definición formal del modelo general, aplicable a cualquier modelo de subastas de información privada aquí: \url{https://www.youtube.com/watch?v=puv3pfCVvNA&list=PLcrc6i6xwaQT52YtKdIpdARVZ4BTFOrQP&index=2}

Nuestro estudio de las subastas se centrará en exponer los tipos de subastas más frecuentemente utilizadas, con especial referencia a:

\begin{lnum}
    \item La subasta de primer precio a sobre cerrado, y
    \item La subasta de segundo precio a sobre cerrado.
\end{lnum}

\subsubsection{Primer precio a sobre cerrado}

En el modelo de subasta a primer precio con sobre cerrado debemos empezar por remarcar que ningún Jugador tendrá incentivos para plantear su estrategia igual que su valor de reserva (tipo o señal, todos estos conceptos son aquí equivalentes) ya que la Utilidad Esperada en caso de ganar será igual a cero. Por lo tanto, la primera diferencia clave respecto al modelo anterior es el tipo de Equilibrio que intentaremos hallar.

Todos los jugadores realizan su puja de manera simultánea y el ganador será aquel con la puja más alta, que será precisamente el pago que deba desembolsar. Por lo tanto, la pregunta será si existe algún Equilibrio en el que los jugadores puedan esperar un pago positivo al realizar su puja. Buscaremos entonces el Equilibrio de Bayes-Nash porque ninguna estrategia será dominante respecto a las demás y, dicho de alguna forma, necesitaremos un concepto más débil de solución / Equilibrio: EBN\footnote{%
    Recordar que en el EBN tomamos la Esperanza de la Utilidad, esto es, tomamos la Esperanza de la Función de Utilidad de Bernoulli sobre la distribución de los tipos de los jugadores.
}.

Nos centraremos en resolver el EBN simétrico: la distribución de probabilidad de las estrategias es común para todos los jugadores; en este caso, una Función de Probabilidad (Variable Aleatoria) cuyo dominio es el conjunto de tipos posibles y recorrido la recta real (en este caso, las estrategias de los jugadores ya que estas son sus \textit{pujas}). Formalmente el problema se define como:

\textbf{Teorema.-} Sea una subasta a primer precio (sobre cerrado) tal que:
\begin{enumerate}
    \item[(i)] $\sigma=\{s_1,s_2,\ldots,s_n\}\equiv$ el perfil de estrategias de los $n$ jugadores,
    \item[(ii)] $S\equiv$ el conjunto de estrategias posibles: $S=S_1\times S_2\times\cdots\times S_n$ (denso en $\mathbb{R}^n$)
    \item[(iii)] $\Theta\equiv$ el conjunto de tipos posibles: $\Theta=\Theta_1\times\Theta_2\times\cdots\times\Theta_n$ (denso en $\mathbb{R}^n$)
    \item[(iv)] $b:\Theta\rightarrow S\equiv$ la función de estrategias cuyo $\mathrm{Dom}(b)$ es el conjunto de tipos (señales o valores reserva) y $\mathrm{Rec}(b)$ el conjunto de estrategias, tal que:
    $$b_i(\cdot)=b(\cdot),\ \forall i=1,2,\ldots,n\quad(\text{estrategias simétricas})\ (*)$$
    \begin{itemize}
        \item[$\rightarrow$] Estrictamente creciente, monótona y diferenciable;
        \item[$\rightarrow$] Los jugadores $i\neq j$ utilizan una estrategia idéntica: $s_j=b(\theta_j)$ (la estrategia de $j$, $b_j$, se define mediante la distribución común $b(\cdot)$, dado su tipo $\theta_j$).
    \end{itemize}
    \item[(v)] La Función de Utilidad Esperada viene definida por:
    $$u_i(s_i,s_{-i},\theta_i)=(\theta_i-s_i)\times P\left[s_j=b(\theta_j)\leq s_i,\ \forall j\neq i\right]$$
\end{enumerate}
Entonces $\exists$ EBN (simétrico) cuyo perfil de estrategias $\sigma$ viene definido por:
$$\max_{\{b_i\}}\ (\theta_i-s_i)\,F^{n-1}\left[b^{-1}(s_i)\right]$$

Nótese que, en lo señalado $(*)$, simplemente estamos restringiendo las dimensiones sobre las que aplica $b$: $b:\Theta_1\times\Theta_2\times\cdots\times\Theta_n\rightarrow S_1\times S_2\times\cdots\times S_n$ si $b_i(\cdot)=b(\cdot)$ $\forall i\in N$ $\Longrightarrow b:\Theta\rightarrow S$.

Es importantísimo entender que, en realidad, la función $b$ se puede entender como la definición de un Vector de Variables Aleatorias Independientes, idénticamente distribuidas (por las suposiciones que hacemos en el punto (iv)). La existencia del Equilibrio Bayes-Nash y, en concreto, la equivalencia entre la probabilidad de que la puja (o estrategia) del Jugador $i$ sea mayor o igual a la puja (o estrategia) de cualquier otro Jugador y el producto de las Funciones de Distribución de la inversa de la Variable Aleatoria se puede realizar por el hecho de que este vector inicial de variables aleatorias idénticamente distribuidas, restringido con posterioridad a una variable aleatoria por la propiedad de independencia que, en nuestro caso, viene dado por la Hipótesis de Independencia en la Valoración:

$$P\left[s_j=b(\theta_j)\leq s_i\ \ \forall j\neq i\right]\sim F^{n-1}\left[b^{-1}(s_i)\right];\quad\text{¿por qué?}$$
Sea $(s_j=)\,b(\theta_j)\leq s_i$; sabemos que $b(\cdot)\in\mathcal{C}^k$, $k\geq 2$, entonces $\theta_j\leq b^{-1}(s_i)$, ninguno sabe el tipo/señal real de otro pero, todos saben que $\forall i,j\in N$, $s_i$ o $s_j\sim F(\cdot)$ (se distribuyen como una Función de Probabilidad) (nótese que, prescindiendo de notación, aquí la Variable Aleatoria es $b$).

$\Longrightarrow$ Por la hipótesis de Valoración Independiente (tipos o señales independientes) ($\sim$ V.A.\ i.i.d.):
$$P\left[\theta_j\leq b^{-1}(s_i)\right]=F^{n-1}\left[b^{-1}(s_i)\right]\quad(\text{para V.A. i.i.d.})$$

Entonces, como cualquier problema de optimización, para encontrar el Equilibrio deberemos primero obtener las CPO y, de estas, obtener una Ecuación Diferencial de Primer Orden:

Desde
$$\max_{\{b_i\}}\ (\theta_i-s_i)\cdot F^{n-1}\left[b^{-1}(s_i)\right]\qquad\left[\text{Recuérdese que }\forall s,\ s=b(\theta)\right]$$
CPO:
$$\begin{aligned}\frac{d}{ds_i}(\theta_i-s_i)\,F^{n-1}\left[b^{-1}(s_i)\right]&=(\theta_i-s_i)(n-1)F^{n-2}\left[b^{-1}(s_i)\right]\cdot\underbrace{f\left(b^{-1}(s_i)\right)\frac{1}{b'\left(b^{-1}(s_i)\right)}}_{\text{Recordar: }\frac{d}{dx}\text{CDF}=\text{PDF}}\\&\quad-F^{n-1}\left(b^{-1}(s_i)\right)=0,\qquad\forall i\in N\end{aligned}$$
donde $b'\left(b^{-1}(s_i)\right)=\frac{d}{dx}b\left(b^{-1}(s_i)\right)$.

En el EBN (simétrico) las respuestas óptimas del Jugador $i$ dado que las estrategias de todo el resto se distribuyen atendiendo a $b(\cdot)$ (y suponemos que $s_i\neq b(\theta_i)$) queremos probar que también se distribuyen como $b(\cdot)$:
$$(\theta-s)(n-1)F^{n-2}[\theta]\,f(\theta)\frac{1}{b'(\theta)}=F^{n-1}[\theta]\ \Longrightarrow\ (\theta-s)(n-1)\frac{1}{b'(\theta)}=\frac{F(\theta)}{f(\theta)}$$
(Obtenemos una ecuación diferencial de primer orden) $\Longrightarrow$
$$b'(s)=\frac{f(\theta)}{F(\theta)}(\theta-s)(n-1)$$

Al resolver esta ecuación diferencial ordinaria de primer orden (\textit{1st-order ODE}), obtendremos la solución al sistema dado cualquier tipo de Jugador. Es decir, será el Equilibrio simétrico de Bayes-Nash, de tal forma que ningún Jugador puede mejorar su estrategia, dadas las de sus rivales (respuesta óptima) pero, además, ninguna solución pasará nunca por el mismo punto, es decir, habrá una asignación clara y diferente para cada Jugador (en función de la condición inicial: su \textit{tipo, señal} o \textit{valor reserva}), definido por:
$$b(\theta)=\theta-\frac{\int_{\underline{\theta}}^{\theta}F^{n-1}(t)\,dt}{F^{n-1}(\theta)}$$

La derivación analítica de esta fórmula se encuentra al final del tema, en las últimas notas\footnote{%
    Dado $b'(\theta)=\frac{f(\theta)}{F(\theta)}(\theta-s)(n-1)$, podemos desarrollarlo t.q.:
    $$F(\theta)\,b'(\theta)+b(\theta)(n-1)f(\theta)=\theta(n-1)f(\theta)\qquad[s=b(\theta)]$$
    Multiplicando por $F^{n-2}(\theta)$:
    $$\underbrace{F^{n-1}(\theta)\,b'(\theta)+b(\theta)(n-1)f(\theta)F^{n-2}(\theta)}_{=\frac{d}{d\theta}\left[F^{n-1}(\theta)\,b(\theta)\right]}=\theta(n-1)f(\theta)F^{n-2}(\theta)$$
    entonces tenemos:
    $$\frac{d}{d\theta}\left[F^{n-1}(\theta)\,b(\theta)\right]=\theta(n-1)f(\theta)F^{n-2}(\theta)$$
    $$\int\frac{d}{d\theta}\left[F^{n-1}(\theta)\,b(\theta)\right]=\int\theta(n-1)f(\theta)F^{n-2}(\theta)$$
    Como las señales o tipos $\theta\in[\underline{\theta},\bar{\theta}]$, tenemos,
    $$\int_{\underline{\theta}}^{\theta}\frac{d}{d\theta}\left[F^{n-1}(\theta)\,b(\theta)\right]=\int_{\underline{\theta}}^{\theta}\theta(n-1)f(\theta)F^{n-2}(\theta);\quad\text{sea }t=\theta$$
    tenemos,
    $$\int_{\underline{\theta}}^{\theta}\frac{d}{dt}\left[F^{n-1}(t)\,b(t)\right]dt=\int_{\underline{\theta}}^{\theta}t\,(n-1)f(t)F^{n-2}(t)\,dt$$
    donde el lado izquierdo es $F^{n-1}(\theta)b(\theta)-\underbrace{F^{n-1}(\underline{\theta})b(\underline{\theta})}_{=0}$ y, para el lado derecho, sea $u=t$ y $dv=(n-1)f(t)F^{n-2}(t)\Rightarrow v=F^{n-1}(t)$, por tanto, $\int u\,dv=uv-\int v\,du\Rightarrow tF^{n-1}(t)-\int_{\underline{\theta}}^{\theta}F^{n-1}(t)\,dt$ (con $t=\theta$) $=\theta F^{n-1}(\theta)-\int_{\underline{\theta}}^{\theta}F^{n-1}(t)\,dt$.

    Por lo tanto,
    $$F^{n-1}(\theta)\,b(\theta)=\theta F^{n-1}(\theta)-\int_{\underline{\theta}}^{\theta}F^{n-1}(t)\,dt\ \Longrightarrow\ b(\theta)=\theta-\underbrace{\frac{\int_{\underline{\theta}}^{\theta}F^{n-1}(t)\,dt}{F^{n-1}(\theta)}}_{(*)}$$
    Nótese que $(*)$ es siempre positivo por lo que, en equilibrio, los jugadores siempre pujarán algo ligeramente inferior a su tipo o valor reserva.
}. Nótese que este será siempre ligeramente inferior al \textit{valor de reserva} correspondiente al tipo de Jugador.

\subsubsection{Segundo precio a sobre cerrado}

Una subasta Vickrey (por WILLIAM VICKREY, 1961) es un tipo de subasta de puja sellada, donde los oferentes presentan ofertas simultáneas por escrito sin conocer la oferta de las otras personas en la subasta, y en la que gana el postor más alto, pero el precio que paga este es la segunda oferta más alta\footnote{%
    Si hay dos pujas iguales, imaginemos que tiramos una moneda y el ganador se decide por puro azar. La segunda \textit{puja} más alta será en este caso particular la misma que la ganadora.
}.

\textbf{Teorema.-} En una subasta a segundo precio, siempre hay una estrategia débilmente dominante para \textit{pujar} por el valor que uno atribuye al objeto: $b_{i}(\theta_{i})=\theta_{i}$.

Sea una subasta de $n$ jugadores:
\begin{itemize}
    \item[] $b_i\equiv$ la estrategia del Jugador $i$; en este caso, la puja: $b_i\in\mathbb{R}_+$ o $b_i\in\mathbb{R}_{[\underline{\theta},\bar{\theta}]}$; donde $\theta\equiv$ el intervalo de señales, siendo $\underline{\theta}\equiv$ ínfimo y $\bar{\theta}\equiv$ supremo.
    \item[] $\theta_i\equiv$ el tipo del Jugador $i$, $\theta_i\in[\underline{\theta},\bar{\theta}]$; como la señal del Jugador (el tipo) se define sobre el conjunto de señales, hay infinitos tipos para $i$ (y para cualquier otro Jugador, p.e. $j:\ \theta_j\in[\underline{\theta},\bar{\theta}]$).
    \item[] $u_i\equiv$ Una función de Utilidad tal que:
\end{itemize}
$$u_i(\underbrace{b_i,b_{-i}}_{(1)},\theta_i)=\begin{cases}\underbrace{v(\theta_i)}_{(2)}-\underbrace{\max_{\{j\neq i\}}b_j}_{(3)}, & \text{si } b_i>\max_{\{j\neq i\}}b_j,\ \forall j\in N\\[4pt] 0, & \text{de otra forma}\end{cases}$$

Donde (1) es el perfil de estrategias del Jugador $i$ (su estrategia, dada la de sus oponentes, \textit{optimal response}); (2) es la Utilidad que le proporciona por ser del tipo $i$, esto es, el tipo estará asociado en este caso particular a una valoración subjetiva del objeto; y, (3) el valor de la segunda puja más alta, siempre y cuando su estrategia (su \textit{puja}) sea de hecho la más alta, si no le reportará cero. ¿Qué quiere decir? En una subasta del tipo Vickrey existe un Equilibrio por estrategias dominadas que consiste en que todos los jugadores pujen por el valor que cada uno atribuye al objeto.

Nótese que, aunque se trata de un problema de información imperfecta el Teorema no nos habla de un Equilibrio de Bayes-Nash (que sería lo normal en este caso) sino de una estrategia que domina débilmente (es decir, que independientemente de qué hagan los otros Jugadores, el Jugador $i$ siempre prefiere pujar por su valor de reserva o tipo que, en definitiva, nos conducirá a un Equilibrio por estrategias dominadas\footnote{%
    Por lo que tampoco será necesario calcular la Utilidad Esperada, que no tiene cabida en el Equilibrio por estrategia débilmente dominante, ya que esta será tal para cualquier distribución de probabilidad.
}), porque precisamente el Equilibrio por estrategias dominadas (aunque débilmente) es \textit{más fuerte} que el Equilibrio de Bayes-Nash. Dicho de otra forma, el Equilibrio que vamos a encontrar es también un Equilibrio de Bayes-Nash. La prueba formal de este Teorema es muy sencilla:

\textbf{Teorema.-} Supongamos una subasta a segundo precio (sobre cerrado), entonces tenemos que $\forall i\in N$, la estrategia $b_i(\theta_i)=\theta_i$ domina débilmente a todas las demás estrategias de $i$, dado su tipo/señal $\theta_i$.

\textbf{Prueba:} Supongamos que el tipo/señal (o valor reserva) de $i$ es $\theta_i$, y este tiene que planificar su estrategia. Denotemos por $h\equiv$ el valor de la puja de los otros jugadores t.q. $h=\max_{\{j\neq i\}}b_j$, entonces, imaginemos dos escenarios:

\begin{enumerate}
    \item[(1)] $b_i<s_i$, en cuyo caso tendríamos:
    \begin{enumerate}
        \item[(1.1)] $h<b_i<s_i$: $\quad u_i(b_i,b_{-i},\theta_i)\simeq s_i-h>0$
        \item[(1.2)] $b_i<h<s_i$: $\quad u_i(b_i,b_{-i},\theta_i)=0\ (*)$ \quad(pierde)
        \item[(1.3)] $b_i<s_i<h$: $\quad u_i(b_i,b_{-i},\theta_i)=0$ \quad(pierde)
    \end{enumerate}
    \item[(2)] $s_i<b_i$, en cuyo caso tendríamos:
    \begin{enumerate}
        \item[(2.1)] $h<s_i<b_i$: $\quad u_i(b_i,b_{-i},\theta_i)\simeq s_i-h>0$
        \item[(2.2)] $s_i<h<b_i$: $\quad u_i(b_i,b_{-i},\theta_i)\simeq s_i-h<0\ (*)$
        \item[(2.3)] $s_i<b_i<h$: $\quad u_i(b_i,b_{-i},\theta_i)=0$ \quad(pierde)
    \end{enumerate}
\end{enumerate}

Los dos casos señalados $(*)$ son los que nos indican la dominancia de $b_i(\theta_i)=\theta_i$, ya que en ambos la $u_i$ hubiese sido estrictamente mayor con dicha estrategia. Con cualquier otro escenario, hubiese sido igual (por eso débilmente).

\subsubsection{Teorema de equivalencia de ingresos (\textit{revenue equivalence theorem})}

Antes de cerrar este capítulo, veamos uno de los resultados más importantes dentro de la Teoría de Subastas: el Teorema de Equivalencia de Resultados.

\textbf{Teorema.-} Cualquier subasta con Información Privada Independiente, con una distribución de probabilidades común (esto es, que la Variable Aleatoria de asignación de \textit{tipos}, \textit{señales} o \textit{valores de reserva}, forme un Vector de V.A. independientes, idénticamente distribuidas, i.i.d.) en la que:

\begin{lnum}
    \item El número de jugadores sea el mismo ($N$) y todos ellos sean neutrales al riesgo;
    \item El objeto de la subasta siempre lo gane el Jugador cuya \textit{puja} (o estrategia) sea la más alta; y,
    \item El Jugador con el \textit{valor de reserva} más bajo tenga un pago esperado de cero (\textit{zero-surplus})\footnote{%
    Esta propiedad implica que el Jugador cuyo valor de reserva sea el más bajo no tendrá ni un beneficio esperado positivo, ni un beneficio esperado negativo.
}.
\end{lnum}

Entonces, tendrán siempre el mismo \textit{pago} esperado.

La intuición detrás de este Teorema es que toda subasta en la que la \textit{función de estrategia} (o función de distribución de \textit{tipos}) sea creciente e idénticamente distribuida para todos los jugadores, podremos obviar estas ya que dicha función será lineal y por tanto, equivalente a calcular el valor esperado (\textit{expected revenue}) de la esperanza de la distribución de la función de estrategia\footnote{%
    Si los jugadores no son todos neutrales al riesgo, no podemos garantizar que la función de estrategia sea lineal, por tanto no aplicará el resultado del Teorema.
}.

Veamos este Teorema aplicado a las dos subastas que hemos estudiado. El único Equilibrio simétrico de Bayes-Nash (o el único Equilibrio Secuencial) en la subasta de precio ascendente conlleva que cada Jugador abandone el Juego cuando alcanza su \textit{valor reserva}. Análogamente, si pensamos en el contexto de una subasta Vickrey podemos decir que el Juego termina cuando el Jugador con el segundo \textit{valor de reserva} más alto abandona el Juego, por lo que el Jugador ganador paga una suma equivalente al segundo mayor \textit{valor de reserva}.

\textbf{Proposición.-} La subasta a primer precio y la subasta a segundo precio (Subasta Vickrey) conducen al mismo Pago Esperado en términos de Beneficio para el vendedor.

\paragraph{Beneficio Esperado}

Sea una disposición de Juego tal que hay $N$ jugadores, un Vector de V.A. que asigna un tipo a cada Jugador (para cada jugador) y con V.A. i.i.d. sobre un intervalo concreto (conocido, supongamos, pero no es importante). Sea $Y$ una V.A. que denota el segundo \textit{valor reserva} (o tipo) más elevado. Sabemos que el Beneficio del vendedor es el pago del Jugador ganador. Por tanto:

Sea
\begin{enumerate}
    \item[(i)] $Y$ una v.a.: $Y\sim[\underline{\theta},\bar{\theta}]$ relativa al segundo pago (por arriba);
    \item[(ii)] $X$ un vector de v.a. i.i.d.: $F_{X_i}\sim[\underline{\theta},\bar{\theta}]$, entonces $F_X\sim[\underline{\theta},\bar{\theta}]$;
    \item[(iii)] $G$ una v.a. $\equiv$ Probabilidad de que el tipo $\theta$ gane. Entonces
    $$F_G(\theta)=P\left[Y\leq\theta\right]\ \Longrightarrow\ F_G(\theta)=F^{N-1}(\theta);\qquad f_g(\theta)=\frac{d}{d\theta}F_G(\theta)$$
\end{enumerate}

Definimos el pago esperado de un jugador $i$ cuyo tipo es $\theta_i^*$ (realizado):
$$m(\theta_i^*)=\int_{\underline{\theta}}^{\theta_i^*}\theta\,f_g(\theta)\,d\theta$$
Por otro lado el Beneficio Esperado \textit{ex ante} de un jugador (cualquiera) vendrá definido como:
$$m=\int_{\underline{\theta}}^{\bar{\theta}}m(\theta)\,f_X(\theta)\,d\theta;$$
Y, el Beneficio Esperado del Juego será entonces:
$$ER=N\,m$$

Si aplicamos esto ante un marco de subasta de tipo Vickrey, obtendremos el mismo Beneficio Esperado que si lo planteamos ante un marco de tipo subasta a primer precio. [Para calcularlo deberíamos definir un intervalo de \textit{tipo} o valores de reserva y una función de distribución para los \textit{tipos}, básicamente]. De hecho, encontraremos en este caso que el Beneficio Esperado depende (normalmente) del número de jugadores y tiende al máximo del intervalo de \textit{tipos} o \textit{valores reserva} si el número de participantes o jugadores tiende a infinito.}




\subsection{Implicaciones de Política Económica}



\textcolor{magenta}{Desarrollar}






\subsubsection{Limitaciones}



\textcolor{magenta}{Desarrollar}















\clearpage
\section{Conclusión} 




\subsection{Recapitulación}



\subsection{Extensiones y relación con otras partes del Temario}



\subsection{Opinión}

\subsubsection{Idea final (Salida o cierra)}






\clearpage
\pagenumbering{gobble}
\MyBibliography

\textbf{Investigación}

\noindent\hangindent=1.5em\hangafter=1 Aumann, R. J. (1976). Agreeing to disagree. \textit{The Annals of Statistics}, 4(6), 1236-1239. \href{https://doi.org/10.1214/aos/1176343654}{\nolinkurl{https://doi.org/10.1214/aos/1176343654}}

\noindent\hangindent=1.5em\hangafter=1 Aumann, R. J. y Shapley, L. S. (1976). Long-term competition: A game-theoretic analysis. Manuscrito, publicado en N. Megiddo (Ed.), \textit{Essays in game theory in honor of Michael Maschler} (1994, pp. 1-15). Springer.

\noindent\hangindent=1.5em\hangafter=1 Benoit, J.-P. y Krishna, V. (1985). Finitely repeated games. \textit{Econometrica}, 53(4), 905-922.

\noindent\hangindent=1.5em\hangafter=1 Bernheim, B. D. (1984). Rationalizable strategic behavior. \textit{Econometrica}, 52(4), 1007-1028. \href{https://doi.org/10.2307/1911196}{\nolinkurl{https://doi.org/10.2307/1911196}}

\noindent\hangindent=1.5em\hangafter=1 Binmore, K., Rubinstein, A. y Wolinsky, A. (1986). The Nash bargaining solution in economic modelling. \textit{The RAND Journal of Economics}, 17(2), 176-188.

\noindent\hangindent=1.5em\hangafter=1 Bondareva, O. N. (1963). Some applications of linear programming methods to the theory of cooperative games. \textit{Problemy Kibernetiki}, 10, 119-139 [en ruso].

\noindent\hangindent=1.5em\hangafter=1 Camerer, C. F., Ho, T.-H. y Chong, J.-K. (2004). A cognitive hierarchy model of games. \textit{The Quarterly Journal of Economics}, 119(3), 861-898. \href{https://doi.org/10.1162/0033553041502225}{\nolinkurl{https://doi.org/10.1162/0033553041502225}}

\noindent\hangindent=1.5em\hangafter=1 Capen, E. C., Clapp, R. V. y Campbell, W. M. (1971). Competitive bidding in high-risk situations. \textit{Journal of Petroleum Technology}, 23(6), 641-653.

\noindent\hangindent=1.5em\hangafter=1 Cho, I.-K. y Kreps, D. M. (1987). Signaling games and stable equilibria. \textit{The Quarterly Journal of Economics}, 102(2), 179-221. \href{https://doi.org/10.2307/1885060}{\nolinkurl{https://doi.org/10.2307/1885060}}

\noindent\hangindent=1.5em\hangafter=1 Cournot, A. A. (1838). \textit{Recherches sur les principes mathématiques de la théorie des richesses}. Hachette.

\noindent\hangindent=1.5em\hangafter=1 Dasgupta, P. y Maskin, E. (1986). The existence of equilibrium in discontinuous economic games, I: Theory. \textit{The Review of Economic Studies}, 53(1), 1-26.

\noindent\hangindent=1.5em\hangafter=1 Debreu, G. (1952). A social equilibrium existence theorem. \textit{Proceedings of the National Academy of Sciences}, 38(10), 886-893. \href{https://doi.org/10.1073/pnas.38.10.886}{\nolinkurl{https://doi.org/10.1073/pnas.38.10.886}}

\noindent\hangindent=1.5em\hangafter=1 Fan, K. (1952). Fixed-point and minimax theorems in locally convex topological linear spaces. \textit{Proceedings of the National Academy of Sciences}, 38(2), 121-126. \href{https://doi.org/10.1073/pnas.38.2.121}{\nolinkurl{https://doi.org/10.1073/pnas.38.2.121}}

\noindent\hangindent=1.5em\hangafter=1 Friedman, J. W. (1971). A non-cooperative equilibrium for supergames. \textit{The Review of Economic Studies}, 38(1), 1-12. \href{https://doi.org/10.2307/2296617}{\nolinkurl{https://doi.org/10.2307/2296617}}

\noindent\hangindent=1.5em\hangafter=1 Fudenberg, D. y Maskin, E. (1986). The folk theorem in repeated games with discounting or with incomplete information. \textit{Econometrica}, 54(3), 533-554.

\noindent\hangindent=1.5em\hangafter=1 Fudenberg, D. y Tirole, J. (1991). \textit{Game theory}. MIT Press.

\noindent\hangindent=1.5em\hangafter=1 Gibbons, R. (1992). \textit{Game theory for applied economists}. Princeton University Press.

\noindent\hangindent=1.5em\hangafter=1 Glicksberg, I. L. (1952). A further generalization of the Kakutani fixed point theorem, with application to Nash equilibrium points. \textit{Proceedings of the American Mathematical Society}, 3(1), 170-174. \href{https://doi.org/10.2307/2032478}{\nolinkurl{https://doi.org/10.2307/2032478}}

\noindent\hangindent=1.5em\hangafter=1 Harsanyi, J. C. (1967-1968). Games with incomplete information played by Bayesian players, I-III. \textit{Management Science}, 14(3), 159-182; 14(5), 320-334; 14(7), 486-502.

\noindent\hangindent=1.5em\hangafter=1 Harsanyi, J. C. (1973). Games with randomly disturbed payoffs: A new rationale for mixed-strategy equilibrium points. \textit{International Journal of Game Theory}, 2(1), 1-23. \href{https://doi.org/10.1007/BF01737554}{\nolinkurl{https://doi.org/10.1007/BF01737554}}

\noindent\hangindent=1.5em\hangafter=1 Kakutani, S. (1941). A generalization of Brouwer's fixed point theorem. \textit{Duke Mathematical Journal}, 8(3), 457-459. \href{https://doi.org/10.1215/S0012-7094-41-00838-4}{\nolinkurl{https://doi.org/10.1215/S0012-7094-41-00838-4}}

\noindent\hangindent=1.5em\hangafter=1 Kalai, E. y Smorodinsky, M. (1975). Other solutions to Nash's bargaining problem. \textit{Econometrica}, 43(3), 513-518.

\noindent\hangindent=1.5em\hangafter=1 Keynes, J. M. (1936). \textit{The general theory of employment, interest and money}. Macmillan.

\noindent\hangindent=1.5em\hangafter=1 Kohlberg, E. y Mertens, J.-F. (1986). On the strategic stability of equilibria. \textit{Econometrica}, 54(5), 1003-1037. \href{https://doi.org/10.2307/1912320}{\nolinkurl{https://doi.org/10.2307/1912320}}

\noindent\hangindent=1.5em\hangafter=1 Kreps, D. M., Milgrom, P., Roberts, J. y Wilson, R. (1982). Rational cooperation in the finitely repeated prisoners dilemma. \textit{Journal of Economic Theory}, 27(2), 245-252.

\noindent\hangindent=1.5em\hangafter=1 Kreps, D. M. y Wilson, R. (1982a). Sequential equilibria. \textit{Econometrica}, 50(4), 863-894.

\noindent\hangindent=1.5em\hangafter=1 Kreps, D. M. y Wilson, R. (1982b). Reputation and imperfect information. \textit{Journal of Economic Theory}, 27(2), 253-279. \href{https://doi.org/10.1016/0022-0531(82)90030-8}{\nolinkurl{https://doi.org/10.1016/0022-0531(82)90030-8}}

\noindent\hangindent=1.5em\hangafter=1 Kuhn, H. W. (1953). Extensive games and the problem of information. En H. W. Kuhn y A. W. Tucker (Eds.), \textit{Contributions to the theory of games} (Vol. II, pp. 193-216). Princeton University Press. \href{https://doi.org/10.1515/9781400881970-012}{\nolinkurl{https://doi.org/10.1515/9781400881970-012}}

\noindent\hangindent=1.5em\hangafter=1 Kydland, F. E. y Prescott, E. C. (1977). Rules rather than discretion: The inconsistency of optimal plans. \textit{Journal of Political Economy}, 85(3), 473-491. \href{https://doi.org/10.1086/260580}{\nolinkurl{https://doi.org/10.1086/260580}}

\noindent\hangindent=1.5em\hangafter=1 Lewis, D. K. (1969). \textit{Convention: A philosophical study}. Harvard University Press.

\noindent\hangindent=1.5em\hangafter=1 Mas-Colell, A., Whinston, M. D. y Green, J. R. (1995). \textit{Microeconomic theory}. Oxford University Press.

\noindent\hangindent=1.5em\hangafter=1 Maschler, M., Peleg, B. y Shapley, L. S. (1979). Geometric properties of the kernel, nucleolus, and related solution concepts. \textit{Mathematics of Operations Research}, 4(4), 303-338.

\noindent\hangindent=1.5em\hangafter=1 McKelvey, R. D. y Palfrey, T. R. (1992). An experimental study of the centipede game. \textit{Econometrica}, 60(4), 803-836. \href{https://doi.org/10.2307/2951567}{\nolinkurl{https://doi.org/10.2307/2951567}}

\noindent\hangindent=1.5em\hangafter=1 Milgrom, P. y Roberts, J. (1982a). Limit pricing and entry under incomplete information: An equilibrium analysis. \textit{Econometrica}, 50(2), 443-459. \href{https://doi.org/10.2307/1912637}{\nolinkurl{https://doi.org/10.2307/1912637}}

\noindent\hangindent=1.5em\hangafter=1 Milgrom, P. y Roberts, J. (1982b). Predation, reputation, and entry deterrence. \textit{Journal of Economic Theory}, 27(2), 280-312. \href{https://doi.org/10.1016/0022-0531(82)90031-X}{\nolinkurl{https://doi.org/10.1016/0022-0531(82)90031-X}}

\noindent\hangindent=1.5em\hangafter=1 Milgrom, P. y Roberts, J. (1990). Rationalizability, learning, and equilibrium in games with strategic complementarities. \textit{Econometrica}, 58(6), 1255-1277.

\noindent\hangindent=1.5em\hangafter=1 Milgrom, P. R. y Weber, R. J. (1985). Distributional strategies for games with incomplete information. \textit{Mathematics of Operations Research}, 10(4), 619-632.

\noindent\hangindent=1.5em\hangafter=1 Myerson, R. B. (1991). \textit{Game theory: Analysis of conflict}. Harvard University Press.

\noindent\hangindent=1.5em\hangafter=1 Nagel, R. (1995). Unraveling in guessing games: An experimental study. \textit{The American Economic Review}, 85(5), 1313-1326.

\noindent\hangindent=1.5em\hangafter=1 Nash, J. F. (1950a). Equilibrium points in n-person games. \textit{Proceedings of the National Academy of Sciences}, 36(1), 48-49.

\noindent\hangindent=1.5em\hangafter=1 Nash, J. F. (1950b). \textit{Non-cooperative games} [Tesis doctoral, Princeton University].

\noindent\hangindent=1.5em\hangafter=1 Nash, J. F. (1950c). The bargaining problem. \textit{Econometrica}, 18(2), 155-162.

\noindent\hangindent=1.5em\hangafter=1 Nash, J. (1951). Non-cooperative games. \textit{Annals of Mathematics}, 54(2), 286-295.

\noindent\hangindent=1.5em\hangafter=1 Pearce, D. G. (1984). Rationalizable strategic behavior and the problem of perfection. \textit{Econometrica}, 52(4), 1029-1050. \href{https://doi.org/10.2307/1911197}{\nolinkurl{https://doi.org/10.2307/1911197}}

\noindent\hangindent=1.5em\hangafter=1 Rasmusen, E. (2007). \textit{Games and information: An introduction to game theory} (4.ª ed.). Blackwell.

\noindent\hangindent=1.5em\hangafter=1 Rosenthal, R. W. (1981). Games of perfect information, predatory pricing and the chain-store paradox. \textit{Journal of Economic Theory}, 25(1), 92-100. \href{https://doi.org/10.1016/0022-0531(81)90018-1}{\nolinkurl{https://doi.org/10.1016/0022-0531(81)90018-1}}

\noindent\hangindent=1.5em\hangafter=1 Rubinstein, A. (1979). Equilibrium in supergames with the overtaking criterion. \textit{Journal of Economic Theory}, 21(1), 1-9.

\noindent\hangindent=1.5em\hangafter=1 Rubinstein, A. (1982). Perfect equilibrium in a bargaining model. \textit{Econometrica}, 50(1), 97-109.

\noindent\hangindent=1.5em\hangafter=1 Samuelson, L. (1992). Dominated strategies and common knowledge. \textit{Games and Economic Behavior}, 4(2), 284-313. \href{https://doi.org/10.1016/0899-8256(92)90020-S}{\nolinkurl{https://doi.org/10.1016/0899-8256(92)90020-S}}

\noindent\hangindent=1.5em\hangafter=1 Schelling, T. C. (1960). \textit{The strategy of conflict}. Harvard University Press.

\noindent\hangindent=1.5em\hangafter=1 Schmeidler, D. (1969). The nucleolus of a characteristic function game. \textit{SIAM Journal on Applied Mathematics}, 17(6), 1163-1170.

\noindent\hangindent=1.5em\hangafter=1 Selten, R. (1965). Spieltheoretische Behandlung eines Oligopolmodells mit Nachfrageträgheit. \textit{Zeitschrift für die gesamte Staatswissenschaft}, 121(2), 301-324; 121(4), 667-689.

\noindent\hangindent=1.5em\hangafter=1 Selten, R. (1975). Reexamination of the perfectness concept for equilibrium points in extensive games. \textit{International Journal of Game Theory}, 4(1), 25-55. \href{https://doi.org/10.1007/BF01766400}{\nolinkurl{https://doi.org/10.1007/BF01766400}}

\noindent\hangindent=1.5em\hangafter=1 Selten, R. (1978). The chain store paradox. \textit{Theory and Decision}, 9(2), 127-159. \href{https://doi.org/10.1007/BF00131770}{\nolinkurl{https://doi.org/10.1007/BF00131770}}

\noindent\hangindent=1.5em\hangafter=1 Shapley, L. S. (1953). A value for n-person games. En H. W. Kuhn y A. W. Tucker (Eds.), \textit{Contributions to the theory of games} (Vol. II, pp. 307-317). Princeton University Press.

\noindent\hangindent=1.5em\hangafter=1 Shapley, L. S. (1967). On balanced sets and cores. \textit{Naval Research Logistics Quarterly}, 14(4), 453-460.

\noindent\hangindent=1.5em\hangafter=1 Shapley, L. S. (1971). Cores of convex games. \textit{International Journal of Game Theory}, 1(1), 11-26.

\noindent\hangindent=1.5em\hangafter=1 Spence, M. (1973). Job market signaling. \textit{The Quarterly Journal of Economics}, 87(3), 355-374. \href{https://doi.org/10.2307/1882010}{\nolinkurl{https://doi.org/10.2307/1882010}}

\noindent\hangindent=1.5em\hangafter=1 Stahl, D. O. y Wilson, P. W. (1995). On players' models of other players: Theory and experimental evidence. \textit{Games and Economic Behavior}, 10(1), 218-254. \href{https://doi.org/10.1006/game.1995.1031}{\nolinkurl{https://doi.org/10.1006/game.1995.1031}}

\noindent\hangindent=1.5em\hangafter=1 Thompson, F. B. (1952). \textit{Equivalence of games in extensive form} (Research Memorandum RM-759). RAND Corporation.

\noindent\hangindent=1.5em\hangafter=1 Vickrey, W. (1961). Counterspeculation, auctions, and competitive sealed tenders. \textit{The Journal of Finance}, 16(1), 8-37.

\noindent\hangindent=1.5em\hangafter=1 Von Neumann, J. (1928). Zur Theorie der Gesellschaftsspiele. \textit{Mathematische Annalen}, 100(1), 295-320. \href{https://doi.org/10.1007/BF01448847}{\nolinkurl{https://doi.org/10.1007/BF01448847}}

\noindent\hangindent=1.5em\hangafter=1 Von Neumann, J. y Morgenstern, O. (1944). \textit{Theory of games and economic behavior}. Princeton University Press.

\noindent\hangindent=1.5em\hangafter=1 Von Stackelberg, H. (1934). \textit{Marktform und Gleichgewicht}. Julius Springer.

\noindent\hangindent=1.5em\hangafter=1 Zermelo, E. (1913). Über eine Anwendung der Mengenlehre auf die Theorie des Schachspiels. En E. W. Hobson y A. E. H. Love (Eds.), \textit{Proceedings of the Fifth International Congress of Mathematicians} (Vol. II, pp. 501-504). Cambridge University Press.

\textbf{Complementario}

\noindent\hangindent=1.5em\hangafter=1 \href{https://www.gametheory.net/news/news.pl?Concept=Auc&highlight=AUC}{GameTheory.net. Noticias sobre subastas.}





\clearpage
\anexo{Juegos no-cooperativos: Desarrollo formal}\label{anx:no_cooperativos}

Este anexo desarrolla formalmente, de principio a fin y en el mismo orden, la teoría de juegos no cooperativos que el bloque de principales conceptos presenta de forma intuitiva. Primero define con precisión qué es un juego: sus elementos, sus dos representaciones y las estrategias. Después formaliza los conceptos de solución de los juegos estáticos (la dominancia, la mejor respuesta y las estrategias racionalizables, el equilibrio de NASH y el equilibrio bayesiano de NASH) y, a continuación, los refinamientos que exigen los juegos dinámicos: la racionalidad secuencial y el equilibrio de NASH perfecto en subjuegos, los juegos repetidos y, por último, las creencias, con el equilibrio bayesiano perfecto, el equilibrio secuencial, el equilibrio perfecto de mano temblorosa y los juegos de señalización. Va fuera del cuerpo porque la notación no añade intuición a la exposición, pero da la precisión necesaria para manejar los resultados con seguridad y comprobar los ejemplos del texto, que se retoman aquí con la misma notación de letras.\footnote{%
  Se sigue la notación de MAS-COLELL, WHINSTON y GREEN (1995), salvo en los ejemplos, que conservan la del cuerpo del tema.
}

Como la mayor parte de la literatura, el anexo trata los refinamientos como condiciones adicionales sobre el equilibrio de NASH: el objetivo es siempre alcanzar un equilibrio de NASH y exigirle, además, que sea razonable fuera de la senda de equilibrio o que descanse en creencias coherentes. No se desarrolla, en cambio, la inducción hacia adelante, que descarta creencias atendiendo a lo que una desviación revela sobre las intenciones de quien se desvía (KOHLBERG y MERTENS, 1986); solo se menciona al tratar el criterio intuitivo.

\subsection{Representación formal de un juego}

Como vimos en el cuerpo, describir un juego exige responder a cuatro preguntas: quién participa, qué puede hacer cada uno y cuándo, qué sabe al hacerlo y cómo valora los resultados. Formalmente, las respuestas se organizan en cuatro piezas que se construyen en orden: los jugadores, con sus preferencias y su información; la forma extensiva, que fija las reglas sobre un árbol; las estrategias, que se definen sobre esa forma extensiva; y la forma normal, que se deriva de las estrategias. Este apartado las define en ese orden y las ilustra con el juego de entrada del texto: J1 decide si entra (Entrar o No entrar) y J2 fija un precio alto o bajo (Alto o Bajo), con pagos $(a_1,a_2)$, $(b_1,b_2)$, $(c_1,c_2)$ y $(d_1,d_2)$ en los desenlaces Entrar-Alto, Entrar-Bajo, No entrar-Alto y No entrar-Bajo.

\subsubsection{Jugadores, preferencias e información}

Un juego tiene un conjunto finito de jugadores, $\{1,\ldots,I\}$, al que puede añadirse la naturaleza, el jugador $0$, que no tiene preferencias y elige según probabilidades conocidas por todos. En el juego de entrada, $I=2$.

\paragraph{Preferencias y pagos}

Cada combinación de decisiones lleva a un resultado, que puede ser aleatorio si interviene la naturaleza o algún jugador aleatoriza. Se supone que las preferencias de cada jugador sobre las loterías de resultados cumplen los axiomas de VON NEUMANN y MORGENSTERN (1944) (completitud, transitividad, continuidad e independencia), de modo que pueden representarse por la utilidad esperada de una función de utilidad de Bernoulli, $u_i$, definida sobre los resultados ciertos. En el teorema de la utilidad esperada de VON NEUMANN y MORGENSTERN (1944), una lotería que da el resultado $z_k$ con probabilidad $\pi_k$ vale para el jugador $i$:
$$U_i=\sum_k\pi_k\,u_i(z_k)$$
Esta representación es única salvo transformaciones afines positivas: $\tilde{u}_i=\alpha\,u_i+\beta$, con $\alpha>0$, describe las mismas preferencias y, por tanto, el mismo juego. Los pagos de un juego son, por tanto, cardinales, pero solo hasta ese punto: importan sus distancias relativas, no su nivel ni su unidad. El umbral que el cuerpo obtiene al tratar la función de pagos lo ilustra. $q^*=(a_1-c_1)/(a_1-b_1)$ no cambia si todos los pagos de J1 se multiplican por $\alpha>0$ y se les suma $\beta$, porque el factor y la constante se cancelan; sí cambia, en cambio, si se altera la distancia entre ellos sin alterar su orden.

\paragraph{Racionalidad y conocimiento común}

Un jugador es \textbf{racional} si elige una estrategia que maximiza su pago esperado dada una conjetura, es decir, una distribución de probabilidad sobre las estrategias de los demás; la noción se precisa al definir la mejor respuesta. El \textbf{conocimiento común} se formaliza, siguiendo a AUMANN (1976), con un espacio de estados $\Omega$ y, para cada jugador, una partición $\mathcal{P}_i$ de $\Omega$: en el estado $\omega$, el jugador $i$ solo sabe que el estado verdadero pertenece al elemento de su partición que contiene a $\omega$, $\mathcal{P}_i(\omega)$. El jugador $i$ \textit{sabe} el suceso $E\subseteq\Omega$ en $\omega$ si $\mathcal{P}_i(\omega)\subseteq E$, y el conjunto de estados en que lo sabe es $K_i(E)=\{\omega\in\Omega:\mathcal{P}_i(\omega)\subseteq E\}$. Todos lo saben en $K(E)=\bigcap_iK_i(E)$, y $E$ es de conocimiento común en $\omega$ si
$$\omega\in K^n(E)\qquad\forall n\geq1,$$
donde $K^n$ es la aplicación reiterada $n$ veces del operador $K$ (definición de AUMANN, 1976): todos lo saben, todos saben que todos lo saben, y así indefinidamente.\footnote{%
  Con esta formalización, AUMANN (1976) demostró que dos jugadores con una probabilidad a priori común no pueden estar de acuerdo en discrepar: si sus probabilidades a posteriori de un suceso son de conocimiento común, han de ser iguales.
} La teoría supone que la estructura del juego y la racionalidad de todos son de conocimiento común.

\paragraph{Información: tipos y probabilidad a priori común}

Lo que cada jugador sabe de los demás se formaliza con los \textbf{tipos} de HARSANYI (1967-1968). Cada jugador $i$ tiene un conjunto de tipos posibles, $\Theta_i$, y su tipo, $\theta_i\in\Theta_i$, resume toda su información privada; sus pagos pueden depender de él, $u_i(\cdot,\theta_i)$. Los perfiles de tipos, $\theta=(\theta_1,\ldots,\theta_I)\in\Theta=\Theta_1\times\cdots\times\Theta_I$, se distribuyen según una distribución conjunta $F$, la \textbf{probabilidad a priori común}, que es de conocimiento común. La información es completa si cada $\Theta_i$ tiene un solo elemento, e incompleta en otro caso.

La \textbf{transformación de HARSANYI} convierte el juego de información incompleta en uno de información imperfecta sobre una jugada de la naturaleza: la naturaleza elige primero $\theta$ según $F$, cada jugador observa su componente $\theta_i$, y los nodos que solo se distinguen por los tipos ajenos, $\theta_{-i}$, pertenecen a un mismo conjunto de información del jugador $i$. Al observar su tipo, cada jugador actualiza su creencia sobre los ajenos con la regla de Bayes, $F(\theta_{-i}\mid\theta_i)$. En el ejemplo del texto, $\Theta_2=\{\theta_{ac},\theta_{ag}\}$, $\Theta_1$ tiene un único elemento y $F$ asigna probabilidad $p$ a $\theta_{ac}$; como J1 solo tiene un tipo, su creencia sobre el de J2 es la probabilidad a priori.

\subsubsection{Forma extensiva}

La forma extensiva recoge las jugadas que puede hacer cada jugador, el momento en que puede hacerlas, lo que sabe cuando las hace, los resultados en función de las acciones de todos y los pagos de cada uno. Siguiendo a MAS-COLELL, WHINSTON y GREEN (1995), se define mediante siete elementos:

\begin{enumerate}
    \item[(i)] Un conjunto finito de nodos $\mathcal{X}$, un conjunto finito de acciones $\mathcal{A}$ y un conjunto finito de jugadores $\{1,\ldots,I\}$.
    \item[(ii)] Una función $p:\mathcal{X}\rightarrow\{\mathcal{X}\cup\emptyset\}$ que especifica un antecesor único de cada nodo $x$; $p(x)$ es no vacío para todo $x\in\mathcal{X}$ salvo para uno, designado como inicial, $x_0$. Los sucesores inmediatos de $x$ son entonces $s(x)=p^{-1}(x)$, y el conjunto de todos los predecesores y de todos los sucesores del nodo $x$ se obtiene iterando $p(x)$ y $s(x)$. Para que haya una estructura de árbol, esos dos conjuntos han de ser disjuntos (un antecesor del nodo $x$ no puede ser también sucesor suyo). El conjunto de nodos terminales es $T=\{x\in\mathcal{X}: s(x)=\emptyset\}$; los demás nodos, $\mathcal{X}\setminus T$, son nodos de decisión.
    \item[(iii)] Una función $\alpha:\mathcal{X}\setminus\{x_0\}\rightarrow\mathcal{A}$ que da la acción que lleva a cada nodo no inicial $x$ desde su predecesor inmediato $p(x)$, y que asigna acciones distintas a sucesores distintos de un mismo nodo. El conjunto de acciones disponibles en $x$ es:
    $$c(x)=\{a\in\mathcal{A}: a=\alpha(x') \text{ para algún } x'\in s(x)\}$$
    \item[(iv)] Una colección de conjuntos de información $\mathcal{H}$ y una función $H:\mathcal{X}\rightarrow\mathcal{H}$ que asigna cada nodo de decisión $x$ a un conjunto de información $H(x)\in\mathcal{H}$, de modo que los conjuntos de información forman una partición de los nodos de decisión. Todos los nodos de un mismo conjunto de información han de tener las mismas acciones disponibles: $c(x)=c(x')$ si $H(x)=H(x')$. Las acciones disponibles en un conjunto de información $H$ se denotan $C(H)=\{a\in\mathcal{A}: a\in c(x) \text{ para } x\in H\}$.
    \item[(v)] Una función $\iota:\mathcal{H}\rightarrow\{0,1,\ldots,I\}$ que asigna cada conjunto de información al jugador (o a la naturaleza, formalmente el jugador $0$) que mueve en sus nodos. La colección de conjuntos de información del jugador $i$ es $\mathcal{H}_i=\{H\in\mathcal{H}: i=\iota(H)\}$.
    \item[(vi)] Una función $\rho:\mathcal{H}_0\times\mathcal{A}\rightarrow[0,1]$ que asigna probabilidades a las acciones de los conjuntos de información en que mueve la naturaleza, con $\rho(H,a)=0$ si $a\notin C(H)$ y $\sum_{a\in C(H)}\rho(H,a)=1$ para todo $H\in\mathcal{H}_0$.
    \item[(vii)] Una colección de funciones de pagos $u=\{u_1(\cdot),\ldots,u_I(\cdot)\}$ que asignan a cada jugador una utilidad en cada nodo terminal, $u_i:T\rightarrow\mathbb{R}$; pueden entenderse como funciones de utilidad de Bernoulli.
\end{enumerate}

Un juego en forma extensiva queda especificado, en la notación de MAS-COLELL, WHINSTON y GREEN (1995), por la colección:
$$\Gamma_E=\left\{\mathcal{X},\mathcal{A},I,p(\cdot),\alpha(\cdot),\mathcal{H},H(\cdot),\iota(\cdot),\rho(\cdot),u\right\}$$

Caben dos precisiones. La primera: los elementos (i) a (vi) especifican por sí mismos la estructura del juego en forma extensiva; al añadir el elemento (vii) se tiene un \textit{juego representado en forma extensiva}. La segunda: la definición encierra tres supuestos de finitud, que conviene identificar porque en economía son frecuentes los juegos que no los cumplen. Primero, las acciones disponibles en cada nodo forman un conjunto finito, lo que excluye que un jugador escoja cualquier número de un intervalo real $[a,b]$, como una cantidad o un precio; admitir infinitas acciones exige infinitos nodos, pero los elementos (i) a (vii) siguen siendo los básicos de la representación. Segundo, el juego termina: hay juegos sin nodos terminales, como los que se repiten indefinidamente, en los que los pagos se asocian a las sendas completas del juego en lugar de a los nodos terminales, con pequeños ajustes en la definición. Tercero, el número de jugadores es finito; también puede admitirse un número infinito de jugadores ajustando la definición de forma natural. Las tres extensiones exigen relajar el supuesto de que el número de nodos es finito, y por eso los juegos con un número finito de nodos se llaman \textit{juegos finitos}. La definición se extiende a los casos infinitos sin mucha dificultad, aunque los resultados pueden diferir de forma importante entre unos y otros.

\paragraph{Información perfecta y memoria perfecta}

Con esta notación, las distinciones del cuerpo se definen de forma precisa. Un juego es de \textbf{información perfecta} si todos sus conjuntos de información tienen un único nodo, $|H|=1$ para todo $H\in\mathcal{H}$; en otro caso es de información imperfecta. Un juego tiene \textbf{memoria perfecta} (KUHN, 1953) si ningún jugador olvida lo que sabía ni lo que hizo: para todo jugador $i$ y todo par de nodos $x$ y $x'$ de un mismo conjunto de información de $i$, la sucesión de conjuntos de información de $i$ que atraviesa el camino desde $x_0$ hasta $x$, y las acciones que $i$ elige en ellos, coinciden con las del camino hasta $x'$. Si no fuera así, $i$ podría distinguir $x$ de $x'$ recordando su propia historia, y no estarían en el mismo conjunto de información. Todos los juegos de este tema tienen memoria perfecta.

\paragraph{El juego de entrada en forma extensiva}

La figura escribe el juego de entrada con esta notación. Sus nodos son $\mathcal{X}=\{x_0,x_1,x_2,t_1,t_2,t_3,t_4\}$: el nodo inicial $x_0$, en el que decide J1; los nodos de decisión de J2, $x_1$ tras la entrada y $x_2$ tras la no entrada; y los nodos terminales $T=\{t_1,t_2,t_3,t_4\}$. La función de predecesores es $p(x_1)=p(x_2)=x_0$, $p(t_1)=p(t_2)=x_1$ y $p(t_3)=p(t_4)=x_2$; las acciones que llevan a cada nodo son $\alpha(x_1)=$ Entrar, $\alpha(x_2)=$ No entrar, $\alpha(t_1)=\alpha(t_3)=$ Alto y $\alpha(t_2)=\alpha(t_4)=$ Bajo.\footnote{%
  En rigor, las acciones de conjuntos de información distintos son acciones distintas, Alto tras la entrada y Alto tras la no entrada, aunque las llamemos igual. En el juego simultáneo, en cambio, son la misma acción, porque se eligen en un único conjunto de información, y por eso $c(x_1)=c(x_2)$, como exige el elemento (iv).
} En la versión secuencial, los conjuntos de información son $\mathcal{H}=\{\{x_0\},\{x_1\},\{x_2\}\}$, con $\iota(\{x_0\})=1$ e $\iota(\{x_1\})=\iota(\{x_2\})=2$, y el juego es de información perfecta. En la simultánea, $\mathcal{H}=\{\{x_0\},\{x_1,x_2\}\}$, y el juego es de información imperfecta. No hay jugadas de la naturaleza, $\mathcal{H}_0=\emptyset$, y los pagos son $u_1(t_1)=a_1$, $u_2(t_1)=a_2$, y así sucesivamente.

\begin{center}\begin{minipage}{\linewidth}\centering {\small\textbf{El juego de entrada en la notación de la forma extensiva}}\par\vspace{0.3em} \includegraphics{JE_anx_arbol.pdf}\par\vspace{0.3em} {\small\textit{Fuente: }Elaboración propia a partir de MAS-COLELL, WHINSTON y GREEN (1995).}\end{minipage}\end{center}

Con información incompleta, la transformación de HARSANYI añade un nodo inicial de la naturaleza, $x_N$, con $\iota(\{x_N\})=0$ y $\rho(\{x_N\},\theta_{ac})=p$, y duplica el resto del árbol, una copia por tipo de J2. Los dos nodos en que decide J1, uno por cada tipo de J2, forman un único conjunto de información de J1, porque solo se diferencian en el tipo ajeno.

\subsubsection{Estrategias}

Sea $\mathcal{H}_i$ la colección de conjuntos de información del jugador $i$, $\mathcal{A}$ el conjunto de acciones posibles y $C(H)\subset\mathcal{A}$ el conjunto de acciones posibles en el conjunto de información $H$. En la notación de MAS-COLELL, WHINSTON y GREEN (1995), una \textbf{estrategia pura} del jugador $i$ es una función que asigna a cada uno de sus conjuntos de información una de las acciones disponibles en él:
$$s_i:\mathcal{H}_i\rightarrow\mathcal{A},\quad\text{con}\quad s_i(H)\in C(H)\ \ \forall H\in\mathcal{H}_i$$
El número de estrategias puras es el producto del número de acciones disponibles en cada conjunto de información del jugador. En el juego de entrada del texto, J2 tiene dos conjuntos de información con dos acciones cada uno y, por tanto, $2\times2=4$ estrategias; si no observa la decisión de J1, un solo conjunto de información y dos estrategias. En general, el jugador $i$ tiene
$$|S_i|=\prod_{H\in\mathcal{H}_i}|C(H)|$$
estrategias puras (MAS-COLELL, WHINSTON y GREEN, 1995), de modo que su número crece con la información: cada conjunto de información adicional multiplica el número de planes posibles. La misma definición cubre los dos casos que el cuerpo presenta como estrategias que son funciones. Con acciones continuas, el conjunto de información de quien observa una variable continua es uno para cada valor observado, y su estrategia es una función: la seguidora de STACKELBERG elige una regla $s_2:Q_1\rightarrow Q_2$ que asigna una cantidad a cada cantidad de la líder. Con información incompleta, cada tipo es un conjunto de información distinto del jugador, y su estrategia es una \textbf{regla de decisión}, $s_i:\Theta_i\rightarrow S_i$, que asigna una acción a cada tipo, como la función de puja de una subasta, $b:\Theta_i\rightarrow\mathbb{R}_+$.

Un \textbf{perfil de estrategias} es un vector $s=(s_1,\ldots,s_I)$ con una estrategia por jugador. Se escribe $s=(s_i,s_{-i})$ cuando interesa separar la del jugador $i$ de las de los demás, $s_{-i}=(s_1,\ldots,s_{i-1},s_{i+1},\ldots,s_I)$, y $S=S_1\times\cdots\times S_I$ es el conjunto de perfiles.

Dado un conjunto finito de estrategias puras $S_i$, una \textbf{estrategia mixta} del jugador $i$ es una distribución de probabilidad $\sigma_i\in\Delta(S_i)$, que asigna a cada estrategia pura $s_i$ una probabilidad $\sigma_i(s_i)\geq0$, con $\sum_{s_i\in S_i}\sigma_i(s_i)=1$. Si los jugadores aleatorizan de forma independiente, el pago esperado del jugador $i$ con el perfil de estrategias mixtas $\sigma=(\sigma_1,\ldots,\sigma_I)$ es, en la notación de MAS-COLELL, WHINSTON y GREEN (1995):
$$u_i(\sigma)=\sum_{s\in S}\left[\sigma_1(s_1)\,\sigma_2(s_2)\cdots\sigma_I(s_I)\right]u_i(s)$$
El juego en que se admiten estrategias mixtas es la \textit{extensión mixta} del juego original, $\Gamma_N=\left[I,\{\Delta(S_i)\},\{u_i(\cdot)\}\right]$.

En un juego en forma extensiva, una \textbf{estrategia de comportamiento} del jugador $i$ asigna a cada conjunto de información $H\in\mathcal{H}_i$ una distribución de probabilidad sobre las acciones disponibles en él, $\beta_i(\cdot\mid H)\in\Delta(C(H))$: el jugador aleatoriza cada vez que mueve, y no una sola vez al principio. Las dos formas de aleatorizar tienen dimensiones muy distintas. Una estrategia mixta es un punto del símplice de $\prod_{H\in\mathcal{H}_i}|C(H)|$ estrategias puras y necesita, por tanto, $\prod_{H}|C(H)|-1$ números; una de comportamiento solo necesita $\sum_{H}\left(|C(H)|-1\right)$. Dos estrategias del jugador $i$ son \textit{equivalentes en resultados} si, frente a cualquier perfil de estrategias de los demás, inducen la misma distribución de probabilidad sobre los nodos terminales.

\textbf{Teorema} (KUHN, 1953).- En un juego finito con memoria perfecta, para toda estrategia mixta de un jugador existe una estrategia de comportamiento equivalente en resultados, y viceversa.

La memoria perfecta es imprescindible: si un jugador puede olvidar lo que hizo, una estrategia mixta puede correlacionar sus decisiones en distintos momentos de un modo que ninguna estrategia de comportamiento, que aleatoriza en cada conjunto de información por separado, puede reproducir. El juego de entrada lo ilustra. J2 tiene cuatro estrategias puras, que numeramos (Alto, Alto), (Alto, Bajo), (Bajo, Alto) y (Bajo, Bajo), con el precio tras la entrada en primer lugar. Una estrategia mixta $\lambda=(\lambda_1,\lambda_2,\lambda_3,\lambda_4)$ necesita tres números; una de comportamiento, $\beta=(\beta_E,\beta_N)$, con las probabilidades de bajar el precio tras la entrada y tras la no entrada, solo dos. La estrategia de comportamiento equivalente a $\lambda$ asigna a cada conjunto de información la probabilidad de los planes que bajan el precio en él,
$$\beta_E=\lambda_3+\lambda_4,\qquad\beta_N=\lambda_2+\lambda_4,$$
y, a la inversa, una estrategia mixta equivalente a $\beta$ es la que trata los dos sorteos como independientes,
$$\lambda=\left((1-\beta_E)(1-\beta_N),\ (1-\beta_E)\beta_N,\ \beta_E(1-\beta_N),\ \beta_E\beta_N\right),$$
aunque no es la única: cualquier $\lambda$ con las mismas probabilidades marginales induce los mismos resultados (aplicación del teorema de KUHN, 1953, al ejemplo del texto). Por eso los juegos dinámicos con memoria perfecta se analizan con estrategias de comportamiento sin pérdida de generalidad.

\subsubsection{Forma normal}

Para un juego de $I$ jugadores, la representación en forma normal, $\Gamma_N$, especifica para cada jugador $i$ un conjunto de estrategias $S_i$ (con $s_i\in S_i$) y una función de pagos $u_i(s_1,s_2,\ldots,s_I)$, que da los niveles de utilidad de VON NEUMANN-MORGENSTERN asociados al resultado, posiblemente aleatorio, del perfil de estrategias $(s_1,s_2,\ldots,s_I)$. En la notación de MAS-COLELL, WHINSTON y GREEN (1995), el juego en forma normal es:
$$\Gamma_N=\left[I,\{S_i\},\{u_i(\cdot)\}\right]$$
Cuando el número de estrategias es finito, suele numerarse cada una, escribiendo el conjunto de estrategias del jugador $i$ como $S_i=\{s_{1i},s_{2i},\ldots\}$ y refiriéndose a cada estrategia por su número.

Todo juego en forma extensiva tiene su forma normal. Cada perfil de estrategias puras $s$ induce una distribución de probabilidad sobre los nodos terminales, $P(t\mid s)$, que es degenerada si no interviene la naturaleza y que, si interviene, se obtiene multiplicando las probabilidades $\rho$ de sus jugadas a lo largo de cada camino. El pago de la forma normal es la utilidad esperada de esa distribución:
$$u_i(s)=\sum_{t\in T}P(t\mid s)\,u_i(t)$$
(paso de la forma extensiva a la forma normal, en la notación de MAS-COLELL, WHINSTON y GREEN, 1995). En la versión secuencial del juego de entrada, con las filas para las estrategias de J1 y las columnas para las de J2, se obtiene la matriz $2\times4$ del cuerpo:
$$\begin{array}{c|cccc}
 & (\text{Alto},\text{Alto}) & (\text{Alto},\text{Bajo}) & (\text{Bajo},\text{Alto}) & (\text{Bajo},\text{Bajo})\\\hline
\text{Entrar} & (a_1,a_2) & (a_1,a_2) & (b_1,b_2) & (b_1,b_2)\\
\text{No entrar} & (c_1,c_2) & (d_1,d_2) & (c_1,c_2) & (d_1,d_2)
\end{array}$$
Cada fila repite dos veces cada desenlace, porque las estrategias de J2 que coinciden en el conjunto de información que se alcanza dan el mismo resultado.

\paragraph{Equivalencia de representaciones}

La forma normal puede tener estrategias redundantes. Dos estrategias puras de un jugador son \textit{equivalentes} si dan a todos los jugadores los mismos pagos frente a cualquier perfil de los demás; eliminarlas hasta dejar una de cada clase da la \textbf{forma normal reducida}. Ocurre, por ejemplo, cuando un jugador mueve dos veces y su primera acción hace inalcanzable la segunda decisión: los planes que solo difieren en esa decisión son equivalentes.\footnote{%
  En el juego de entrada no hay estrategias equivalentes: (Alto, Alto) y (Alto, Bajo) dan lo mismo si J1 entra, pero no si no entra, y J1 puede no entrar.
} Juegos en forma extensiva distintos pueden tener la misma forma normal reducida; THOMPSON (1952)\textcolor{orange}{\textsuperscript{[Probable]}} caracterizó las transformaciones del árbol que la dejan invariante, como intercambiar el orden de jugadas simultáneas o fundir dos decisiones consecutivas de un mismo jugador en una. KOHLBERG y MERTENS (1986) propusieron que un concepto de solución dependa solo de la forma normal reducida, de modo que dé la misma respuesta en todas las formas extensivas que la comparten. El equilibrio de NASH cumple esa invariancia; el equilibrio de NASH perfecto en subjuegos, no, porque los subjuegos son una propiedad del árbol y no de la matriz, y es lo que el cuerpo resume al decir que los juegos dinámicos se analizan sobre la forma extensiva.

\subsection{Dominancia}

Los conceptos de solución de los juegos estáticos se ordenan, como en el cuerpo, por lo que exigen a las expectativas de los jugadores: la dominancia no exige ninguna conjetura sobre los demás; la racionalizabilidad, conjeturas coherentes con la racionalidad de todos; el equilibrio de NASH, conjeturas correctas; y el equilibrio bayesiano de NASH lleva esa última exigencia a los juegos con información incompleta. Este apartado y los tres siguientes los formalizan en ese orden sobre la forma normal, $\Gamma_N$, y los aplican a la versión simultánea del juego de entrada.

En un juego $\Gamma_N=\left[I,\{S_i\},\{u_i(\cdot)\}\right]$, y en la notación de MAS-COLELL, WHINSTON y GREEN (1995), la estrategia $s_i\in S_i$ está \textbf{estrictamente dominada} si existe otra estrategia $s_i'\in S_i$ que le da al jugador $i$ un pago mayor frente a cualquier combinación de estrategias de los demás:
$$u_i(s_i',s_{-i})>u_i(s_i,s_{-i})\qquad\forall s_{-i}\in S_{-i}$$
Está \textbf{débilmente dominada} si existe $s_i'\in S_i$ tal que $u_i(s_i',s_{-i})\geq u_i(s_i,s_{-i})$ para toda $s_{-i}\in S_{-i}$, con desigualdad estricta para alguna $s_{-i}$. Una estrategia es \textbf{estrictamente dominante} si domina estrictamente a todas las demás estrategias del jugador. Si cada jugador tiene una, el perfil que forman es un equilibrio en estrategias dominantes y es, necesariamente, el único equilibrio de Nash del juego, porque frente a cualquier conjetura la única mejor respuesta de cada jugador es su estrategia dominante.

En la versión simultánea del juego de entrada, Alto domina estrictamente a Bajo para J2 si y solo si $u_2(\text{Alto},s_1)>u_2(\text{Bajo},s_1)$ para las dos estrategias de J1, es decir, si $a_2>b_2$ y $c_2>d_2$: es el caso del J2 acomodaticio. J1 no tiene ninguna estrategia dominante con los pagos de partida, porque $a_1>c_1$ pero $d_1>b_1$. Si J1 ganase incluso con la guerra de precios ($b_1>d_1$), Entrar dominaría estrictamente a No entrar, y (Entrar, Alto) sería el equilibrio en estrategias dominantes del juego.

En la extensión mixta del juego, una estrategia pura puede estar dominada por una estrategia mixta sin estarlo por ninguna pura. Supongamos que J1 tiene tres estrategias, cuyos pagos frente a las dos de J2 son $(x,0)$, $(0,x)$ e $(y,y)$, con $0<y<x/2$. Ni la primera ni la segunda dominan a la tercera, porque cada una le da $0<y$ frente a una de las estrategias de J2; pero la mezcla de ambas con probabilidad $\tfrac{1}{2}$ le da $x/2$ frente a cualquiera de las estrategias de J2, más que el $y$ que le da la tercera. Si $y>x/2$, en cambio, la tercera estrategia no está dominada por nada, aunque sea la peor frente a cada estrategia pura de J2. Para comprobar si una estrategia está estrictamente dominada en la extensión mixta basta con compararla frente a las estrategias puras de los demás (MAS-COLELL, WHINSTON y GREEN, 1995, proposición 8.B.1): $\sigma_i$ domina estrictamente a $\sigma_i'$ si y solo si $u_i(\sigma_i,s_{-i})>u_i(\sigma_i',s_{-i})$ para toda $s_{-i}\in S_{-i}$. De ahí el procedimiento que describen los propios autores:

\begin{cita}[MAS-COLELL, WHINSTON y GREEN][1995][Microeconomic Theory]
En resumen, para hallar el conjunto de estrategias estrictamente dominadas del jugador $i$ en $\Gamma_N=[I,\{\Delta(S_i)\},\{u_i(\cdot)\}]$, podemos eliminar primero las estrategias puras que están estrictamente dominadas aplicando el criterio de la Proposición 8.B.1. Llamemos $S_i^u\subset S_i$ al conjunto de estrategias puras no dominadas del jugador $i$. A continuación, eliminamos las estrategias mixtas del conjunto $\Delta(S_i^u)$ que estén dominadas. El conjunto de estrategias no dominadas (puras y mixtas) del jugador $i$ es exactamente el de las estrategias que quedan en $\Delta(S_i^u)$.

Como cuando solo considerábamos estrategias puras, podemos llevar más lejos la lógica de la eliminación de estrategias estrictamente dominadas en el juego $\Gamma_N=[I,\{\Delta(S_i)\},\{u_i(\cdot)\}]$ mediante la eliminación iterativa. Lo anterior implica que este procedimiento iterativo puede llevarse a cabo en dos etapas: primero, se eliminan iterativamente las estrategias puras dominadas con el criterio de la Proposición 8.B.1, aplicado en cada etapa al conjunto de estrategias puras que queda. Llamemos $\{\bar{S}_1^u,\ldots,\bar{S}_I^u\}$ a los conjuntos de estrategias puras restantes. Después, se eliminan las estrategias mixtas de los conjuntos $\{\Delta(\bar{S}_1^u),\ldots,\Delta(\bar{S}_I^u)\}$ que estén dominadas.
\end{cita}

La \textbf{eliminación iterada} de estrategias estrictamente dominadas define una sucesión de juegos reducidos: se parte de $S_i^0=S_i$ y, en cada ronda $k$, el conjunto $S_i^{k}$ contiene las estrategias de $S_i^{k-1}$ que no están estrictamente dominadas en el juego en que cada jugador $j$ dispone solo de $S_j^{k-1}$:
$$S_i^k=\left\{s_i\in S_i^{k-1}:\ \nexists\,\sigma_i\in\Delta(S_i^{k-1})\ \text{tal que}\ u_i(\sigma_i,s_{-i})>u_i(s_i,s_{-i})\ \ \forall s_{-i}\in S_{-i}^{k-1}\right\}$$
(eliminación iterada de estrategias estrictamente dominadas, en la notación de MAS-COLELL, WHINSTON y GREEN, 1995). En un juego finito, la sucesión es decreciente y se estabiliza en un número finito de rondas en $S_i^\infty=\bigcap_kS_i^k$, y su resultado no depende del orden en que se eliminen las estrategias. Un juego es \textbf{resoluble por dominancia} si $S^\infty$ contiene un único perfil. Cada ronda exige un escalón más de conocimiento mutuo de la racionalidad: la $k$-ésima solo está justificada si la racionalidad es de conocimiento mutuo de orden $k-1$, y la eliminación completa, si es de conocimiento común.

En la versión simultánea del ejemplo con J2 acomodaticio, la primera ronda deja $S_1^1=\{\text{Entrar},\text{No entrar}\}$ y $S_2^1=\{\text{Alto}\}$; en la segunda, frente a $S_2^1$, Entrar domina a No entrar ($a_1>c_1$), y $S^2=\{(\text{Entrar},\text{Alto})\}$. El juego es resoluble por dominancia en dos rondas. Con J2 agresivo, en cambio, $S^\infty=S$: ninguna estrategia está dominada, porque cada una es la mejor frente a alguna estrategia del rival.

Con dominancia débil, el resultado sí puede depender del orden de eliminación. Sea un juego en el que J1 elige entre $U$, $M$ y $D$ y J2 entre $L$ y $R$, con pagos $(x,x)$ y $(0,0)$ para $U$, $(x,x)$ e $(y,x)$ para $M$ y $(0,0)$ e $(y,x)$ para $D$, frente a $L$ y $R$ respectivamente, y $x>y>0$. $M$ domina débilmente tanto a $U$ como a $D$. Si se elimina primero $U$, $R$ domina débilmente a $L$ para J2 y el resultado es $(y,x)$; si se elimina primero $D$, es $L$ la que domina débilmente a $R$ y el resultado es $(x,x)$. Por eso la eliminación iterada de estrategias débilmente dominadas no define un único conjunto de soluciones, salvo que se fije el orden, por ejemplo eliminando en cada ronda todas las estrategias débilmente dominadas de todos los jugadores a la vez.

\subsection{Mejor respuesta y estrategias racionalizables}

La dominancia compara estrategias frente a todo lo que pueden hacer los demás; la mejor respuesta, frente a una conjetura concreta. Es la pieza con la que se construyen los demás conceptos de solución, y con ella se define el conjunto de estrategias que pueden deducirse solo de la racionalidad y de su conocimiento común.

La \textbf{correspondencia de mejor respuesta} del jugador $i$ asigna a cada combinación de estrategias (mixtas) de los demás, $\sigma_{-i}$, el conjunto de estrategias que maximizan su pago esperado frente a ella:
$$MR_i(\sigma_{-i})=\left\{\sigma_i\in\Delta(S_i):\ u_i(\sigma_i,\sigma_{-i})\geq u_i(\sigma_i',\sigma_{-i})\ \ \forall\sigma_i'\in\Delta(S_i)\right\}$$
La mejor respuesta se define frente a \textit{una} conjetura concreta; la dominancia compara frente a \textit{todas}. Una estrategia que no es mejor respuesta frente a ninguna conjetura \textit{nunca es mejor respuesta}. En la formulación de BERNHEIM (1984) y PEARCE (1984), las \textbf{estrategias racionalizables} son las que sobreviven a la eliminación iterada de estrategias que nunca son mejor respuesta: en cada ronda se descartan las estrategias que no son mejor respuesta a ninguna conjetura sobre las estrategias de los demás que han sobrevivido a la ronda anterior.

Formalmente, sean $R_i^0=S_i$ y, para $k\geq1$,
$$R_i^k=\left\{s_i\in R_i^{k-1}:\ \exists\,\mu_i\in\Delta\!\left(R_{-i}^{k-1}\right)\ \text{tal que}\ s_i\in\arg\max_{s_i'\in R_i^{k-1}}\ \sum_{s_{-i}}\mu_i(s_{-i})\,u_i(s_i',s_{-i})\right\},$$
donde $\mu_i$ es la conjetura del jugador $i$, una distribución de probabilidad sobre los perfiles de los demás que han sobrevivido a la ronda anterior (estrategias racionalizables de BERNHEIM, 1984, y PEARCE, 1984, en la versión que admite conjeturas correlacionadas). El conjunto de estrategias racionalizables es $R_i=\bigcap_kR_i^k$. De forma equivalente, un producto de conjuntos $B_1\times\cdots\times B_I$ tiene la \textit{propiedad de mejor respuesta} si toda estrategia de cada $B_i$ es mejor respuesta a alguna conjetura sobre $B_{-i}$, y $R$ es el mayor producto con esa propiedad.

Una estrategia estrictamente dominada nunca es mejor respuesta, de modo que las estrategias racionalizables son un subconjunto de las que sobreviven a la eliminación iterada de estrategias estrictamente dominadas. El recíproco es el lema de PEARCE (1984): en un juego finito de dos jugadores, una estrategia nunca es mejor respuesta si y solo si está estrictamente dominada, quizá por una estrategia mixta, y por eso en esos juegos ambos conjuntos coinciden. Con más jugadores, el de las racionalizables puede ser menor si se exige que las conjeturas sobre los distintos rivales sean independientes, porque una estrategia puede ser mejor respuesta a una conjetura correlacionada y a ninguna independiente. Toda estrategia que se juega en un equilibrio de NASH es racionalizable, pero no al revés.

En la versión simultánea del ejemplo, las mejores respuestas a estrategias puras son $MR_1(\text{Alto})=\text{Entrar}$ y $MR_1(\text{Bajo})=\text{No entrar}$, por $a_1>c_1$ y $d_1>b_1$. Con J2 acomodaticio, $MR_2(s_1)=\text{Alto}$ para las dos estrategias de J1, de modo que Bajo nunca es mejor respuesta y desaparece en la primera ronda; en la segunda desaparece No entrar, y $R=\{(\text{Entrar},\text{Alto})\}$. Con J2 agresivo, $MR_2(\text{Entrar})=\text{Bajo}$ y $MR_2(\text{No entrar})=\text{Alto}$: el producto $\{\text{Entrar},\text{No entrar}\}\times\{\text{Alto},\text{Bajo}\}$ tiene la propiedad de mejor respuesta, y todas las estrategias son racionalizables. Frente a una conjetura mixta, en la que J2 baja el precio con probabilidad $q$, la mejor respuesta de J1 es entrar si $q<q^*=(a_1-c_1)/(a_1-b_1)$, no entrar si $q>q^*$ y cualquier mezcla si $q=q^*$; esta correspondencia es la que se corta con la de J2 en el equilibrio en estrategias mixtas.

\subsection{Equilibrio de Nash}

El equilibrio de NASH añade a la racionalidad la exigencia de que las conjeturas sean correctas. Formalmente, un perfil de estrategias mixtas $\sigma^*$ es un equilibrio si existe un perfil de conjeturas $(\mu_1,\ldots,\mu_I)$ tal que, para todo jugador, (i) $\sigma_i^*$ es mejor respuesta a $\mu_i$ y (ii) $\mu_i$ coincide con lo que los demás hacen de verdad, $\mu_i=\sigma_{-i}^*$, con las estrategias de los demás elegidas de forma independiente. La condición (i) es la racionalidad, que la racionalizabilidad también exige; la (ii), la corrección de las conjeturas, es la que distingue el equilibrio, porque la racionalizabilidad solo pide que $\mu_i$ atribuya a los demás estrategias racionalizables. Sustituyendo (ii) en (i) se obtiene la definición habitual.

Un perfil de estrategias $s^*=(s_1^*,\ldots,s_I^*)$ es un \textbf{equilibrio de Nash} en estrategias puras del juego $\Gamma_N$ si, para todo jugador $i$ (NASH, 1950a, 1951, en la notación de MAS-COLELL, WHINSTON y GREEN, 1995):
$$u_i(s_i^*,s_{-i}^*)\geq u_i(s_i,s_{-i}^*)\qquad\forall s_i\in S_i$$
o, de forma equivalente, si $s_i^*\in MR_i(s_{-i}^*)$ para todo $i$: el equilibrio es un punto fijo de la correspondencia de mejor respuesta. Un perfil de estrategias mixtas $\sigma^*=(\sigma_1^*,\ldots,\sigma_I^*)$ es un equilibrio de Nash de la extensión mixta si, para todo $i$:
$$u_i(\sigma_i^*,\sigma_{-i}^*)\geq u_i(\sigma_i,\sigma_{-i}^*)\qquad\forall\sigma_i\in\Delta(S_i)$$

\subsubsection{El principio de indiferencia}

Sea $S_i^+\subset S_i$ el conjunto de estrategias puras a las que $\sigma_i^*$ asigna una probabilidad positiva (su soporte). Según la proposición 8.D.1 de MAS-COLELL, WHINSTON y GREEN (1995), el perfil $\sigma^*$ es un equilibrio de Nash si y solo si, para todo jugador $i$, todas las estrategias puras del soporte le dan el mismo pago esperado frente a $\sigma_{-i}^*$, y ese pago es mayor o igual que el de cualquier estrategia pura fuera del soporte:
$$u_i(s_i,\sigma_{-i}^*)=u_i(s_i',\sigma_{-i}^*)\quad\forall s_i,s_i'\in S_i^+;\qquad u_i(s_i,\sigma_{-i}^*)\geq u_i(s_i'',\sigma_{-i}^*)\quad\forall s_i\in S_i^+,\ s_i''\notin S_i^+$$
Como el pago esperado es una media de los pagos de las estrategias puras ponderada por sus probabilidades, si alguna estrategia del soporte rindiera más que otra, convendría desplazar hacia ella toda la probabilidad. Por eso los equilibrios en estrategias mixtas se calculan igualando pagos: las probabilidades de cada jugador se eligen para que el \textit{otro} quede indiferente.

\subsubsection{Existencia}

\textbf{Teorema} (NASH, 1950a, 1951).- Todo juego $\Gamma_N=\left[I,\{\Delta(S_i)\},\{u_i(\cdot)\}\right]$ en el que los conjuntos $S_1,\ldots,S_I$ tienen un número finito de elementos tiene al menos un equilibrio de Nash.

La demostración aplica un teorema de punto fijo a la correspondencia de mejor respuesta conjunta.

\textbf{Teorema} (KAKUTANI, 1941).- Sea $A\subset\mathbb{R}^N$ un conjunto no vacío, compacto y convexo, y sea $\phi:A\rightrightarrows A$ una correspondencia con valores no vacíos y convexos y con grafo cerrado. Entonces existe $x^*\in A$ tal que $x^*\in\phi(x^*)$.

Basta comprobar que la mejor respuesta conjunta cumple esas condiciones (NASH, 1950a, en la presentación de MAS-COLELL, WHINSTON y GREEN, 1995). Sea $\Sigma=\Delta(S_1)\times\cdots\times\Delta(S_I)$ el conjunto de perfiles de estrategias mixtas y sea $MR:\Sigma\rightrightarrows\Sigma$, con
$$MR(\sigma)=MR_1(\sigma_{-1})\times\cdots\times MR_I(\sigma_{-I})$$
Primero, $\Sigma$ es un producto de símplices de dimensión finita y, por tanto, no vacío, compacto y convexo. Segundo, cada $MR_i(\sigma_{-i})$ es no vacío, porque $u_i(\cdot,\sigma_{-i})$ es continua en el compacto $\Delta(S_i)$ y alcanza en él un máximo, por el teorema de WEIERSTRASS. Tercero, es convexo, porque $u_i$ es lineal en las probabilidades propias: los maximizadores de una función lineal sobre un símplice forman una de sus caras, que es convexa. Cuarto, el grafo de $MR$ es cerrado, porque los pagos son continuos: si $\sigma^n\to\sigma$, $\tau^n\to\tau$ y $\tau^n\in MR(\sigma^n)$ para todo $n$, entonces $\tau\in MR(\sigma)$, por el teorema del máximo de BERGE. El teorema de KAKUTANI garantiza entonces un $\sigma^*\in MR(\sigma^*)$, y ese punto fijo es, por definición, un equilibrio de NASH. NASH (1951) dio una segunda demostración, con el teorema de punto fijo de BROUWER aplicado a una función continua construida a partir de las ganancias de desviarse.

\begin{center}\begin{minipage}{\linewidth}\centering {\small\textbf{El teorema de KAKUTANI y el equilibrio de NASH}}\par\vspace{0.3em} \includegraphics{JE_anx_kakutani.pdf}\par\vspace{0.3em} {\small\textit{Fuente: }Elaboración propia a partir de KAKUTANI (1941) y NASH (1950a).}\end{minipage}\end{center}

La figura muestra el papel de cada condición. En el panel (a), una correspondencia del intervalo $[0,1]$ en sí mismo con grafo cerrado y valores convexos corta necesariamente la diagonal: en el salto, su valor es todo un segmento, que la diagonal atraviesa. En el panel (b), si el valor en el salto no es convexo, porque contiene solo los dos extremos, la correspondencia puede saltar por encima de la diagonal sin cortarla, y no hay punto fijo. Es lo que ocurre con las mejores respuestas en estrategias puras del ejemplo con J2 agresivo, y lo que las estrategias mixtas resuelven: al admitirlas, la mejor respuesta de J1 cuando $q=q^*$ es todo el intervalo de mezclas, y en el panel (c) las dos correspondencias se cortan en $(q^*,p^*)$.

Con conjuntos de estrategias infinitos, DEBREU (1952), GLICKSBERG (1952) y FAN (1952) garantizan la existencia de un equilibrio en estrategias puras si cada $S_i$ es un subconjunto no vacío, compacto y convexo de un espacio euclídeo y cada $u_i$ es continua en el perfil $s$ y cuasicóncava en la propia estrategia $s_i$. La cuasiconcavidad hace aquí el papel de la aleatorización: garantiza que el conjunto de mejores respuestas sea convexo sin necesidad de mezclar. Si los pagos son discontinuos, como en la competencia en precios con capacidades limitadas, el equilibrio en estrategias puras puede no existir, y DASGUPTA y MASKIN (1986)\textcolor{orange}{\textsuperscript{[Probable]}} dan condiciones para que exista en estrategias mixtas.

\subsubsection{Interpretación}

El artículo original de NASH (1950a) no decía nada sobre cómo se logra el equilibrio. En su tesis doctoral, de la que extrajo el artículo, NASH (1950b) sí ofreció dos interpretaciones. La primera ve el equilibrio como consecuencia de la inferencia racional: como todo jugador puede razonar sobre las consideraciones estratégicas a las que se enfrentan sus rivales, la racionalidad bastaría para que cada uno prediga correctamente lo que harán los demás; en palabras del propio NASH, \textit{se trata de una interpretación racionalista e idealizadora}. La segunda, la de \textit{acción de masas}, imagina poblaciones de jugadores que se encuentran repetidamente y que, sin conocer la estructura completa del juego, acumulan experiencia sobre qué estrategias rinden más; si ese aprendizaje se estabiliza, lo hace en un equilibrio de NASH.

A partir de estas ideas, MAS-COLELL, WHINSTON y GREEN (1995) reúnen cinco razones por las que cabe esperar que se juegue un equilibrio de NASH. La primera es la inferencia racional que acabamos de ver.

La segunda es que el equilibrio es una condición necesaria si un juego tiene un único resultado esperado: los jugadores racionales lo entenderán, y para que ninguno quiera desviarse ese resultado ha de ser un equilibrio de NASH; evidentemente, el argumento solo vale si hay un único resultado esperado.

La tercera son los \textbf{puntos focales} de SCHELLING (1960): algunos resultados destacan por razones culturales o de contexto, y todos esperan que se jueguen. SCHELLING preguntó dónde y cuándo intentarían encontrarse en Nueva York dos personas que no han podido acordarlo, y la mayoría eligió el mismo sitio y la misma hora, la estación Grand Central a mediodía. La implicación importante es que esa llamada solo puede convertir un resultado en el resultado esperado si es un equilibrio de NASH, porque si no alguien querría desviarse.

La cuarta ve el equilibrio como un \textbf{acuerdo que se cumple solo} (\textit{self-enforcing}): si los jugadores pueden hablar antes de jugar, pero no vincularse a lo que pacten, cualquier acuerdo que alcancen ha de ser un equilibrio de NASH, porque de lo contrario alguien lo incumpliría.

La quinta ve el equilibrio como una \textbf{convención social estable}: si el juego se repite, una manera de jugarlo puede predominar con el tiempo y volverse focal, y el equilibrio aparece como el punto de reposo de un proceso de ajuste dinámico.\footnote{%
  Un buen ejemplo es el juego al que juegan cada día los neoyorquinos: caminar por el centro de Manhattan. Cada día, quienes van andando al trabajo tienen que decidir por qué lado de la acera caminan. Con el tiempo, la convención social estable es que todo el mundo camina por la derecha, una convención que se hace cumplir porque cualquier persona que se desvíe unilateralmente de ella acabará, con seguridad, gravemente pisoteada. Por supuesto, en un día cualquiera es \textit{posible} que una persona decida caminar por la izquierda conjeturando que todos los demás esperan de repente que la convención cambie. Aun así, la predicción de que permaneceremos en el equilibrio de NASH en el que todo el mundo camina por la derecha parece razonable en este caso. Obsérvese que, si un resultado ha de convertirse en una convención social estable, debe ser un equilibrio de NASH. Si no lo fuera, los individuos se desviarían de él en cuanto empezara a surgir.
} No es fácil modelizar formalmente la aparición de convenciones sociales estables. Una dificultad es que el juego repetido de un día puede verse a su vez como un juego dinámico más amplio. Así, cuando consideramos jugadores racionales que eligen sus estrategias en este juego global, volvemos simplemente a nuestro dilema original: ¿por qué deberíamos esperar un equilibrio de NASH en este juego más amplio? Una respuesta a esta dificultad que recibe actualmente cierta atención formal imagina que los jugadores siguen reglas prácticas sencillas sobre el comportamiento probable de sus oponentes en situaciones en las que el juego se repite (obsérvese que esto implica renunciar en cierta medida al supuesto de racionalidad completa). Por ejemplo, un jugador podría conjeturar que lo que hicieron ayer sus oponentes se repetirá hoy. Si es así, cada día los jugadores jugarán una mejor respuesta a lo que se jugó ayer. Si surge una combinación de estrategias que es un punto estacionario de este proceso (es decir, lo que se juega hoy es lo mismo que se jugó ayer), debe ser un equilibrio de NASH. Sin embargo, está menos claro que, desde cualquier posición inicial, el proceso vaya a converger a un resultado estacionario; la convergencia resulta depender del juego.\footnote{%
    Este enfoque se remonta en realidad al procedimiento de ajuste miope de COURNOT (1838). Un ejemplo reciente se encuentra en MILGROM y ROBERTS (1990). Curiosamente, este trabajo explica el resultado ultrarracional de NASH \textit{relajando} el supuesto de racionalidad. También puede utilizarse para tratar de identificar la probabilidad de que surjan los distintos equilibrios de NASH cuando existen varios.
}

Esta última es la justificación más cercana a la economía, porque recoge el uso tradicional de la noción de equilibrio en la disciplina. En economía, un equilibrio no se ha entendido como el resultado de que cada agente deduzca lo que harán los demás, sino como el estado de reposo de un proceso de ajuste. El precio de un mercado competitivo es de equilibrio porque, una vez alcanzado, nadie tiene motivos para cambiar su conducta, y se llega a él por tanteo, a medida que los excesos de oferta y de demanda corrigen los precios. Los agentes no necesitan conocer la estructura entera del mercado ni prever lo que harán los demás: les basta reaccionar a lo que observan. La convención social lee el equilibrio de NASH de la misma manera: es el punto en el que se detiene un proceso en el que los jugadores, sin deducir nada, ajustan su conducta a lo que han visto hacer a los demás, y en el que las expectativas formadas con la experiencia se confirman. Por eso exige mucha menos racionalidad que la interpretación racionalista y une el equilibrio de NASH con el equilibrio competitivo y con el ajuste miope que ya proponía COURNOT (1838). A cambio, hereda su debilidad: un punto de reposo solo es una buena predicción si el proceso de ajuste converge a él, y eso, como acabamos de ver, depende del juego.

\subsubsection{Cálculo de los equilibrios en estrategias mixtas de los ejemplos del texto}

\textit{Juego de entrada simultáneo con J2 agresivo}. Sea $p$ la probabilidad con que J1 entra y $q$ la probabilidad con que J2 baja el precio, con $a_1>c_1=d_1>b_1$, $b_2>a_2$ y $c_2>d_2$. El pago esperado de J1 es $(1-q)\,a_1+q\,b_1$ si entra y $c_1$ si no entra; su indiferencia exige
$$q^*=\frac{a_1-c_1}{a_1-b_1},$$
que está entre $0$ y $1$ porque $a_1>c_1>b_1$. El pago esperado de J2 es $p\,a_2+(1-p)\,c_2$ si mantiene el precio y $p\,b_2+(1-p)\,d_2$ si lo baja; su indiferencia exige $p\,(b_2-a_2)=(1-p)\,(c_2-d_2)$, es decir,
$$p^*=\frac{c_2-d_2}{(b_2-a_2)+(c_2-d_2)},$$
también entre $0$ y $1$, porque $b_2>a_2$ y $c_2>d_2$. En el equilibrio, J1 entra con probabilidad $p^*$ y J2 baja el precio con probabilidad $q^*$; J1 obtiene en promedio $c_1$, lo mismo que si no entrase, y J2, $p^*a_2+(1-p^*)\,c_2$. Obsérvese que la probabilidad con que mezcla cada jugador depende de los pagos del otro, no de los suyos.

\textit{Cara o cruz}. Sea $q$ la probabilidad con que J2 juega Cara. El pago esperado de J1 es $-q+(1-q)=1-2q$ si juega Cara y $q-(1-q)=2q-1$ si juega Cruz; J1 solo está dispuesto a mezclar si ambos son iguales, lo que exige $q=\tfrac{1}{2}$. Del mismo modo, si $p$ es la probabilidad con que J1 juega Cara, el pago esperado de J2 es $2p-1$ con Cara y $1-2p$ con Cruz, y su indiferencia exige $p=\tfrac{1}{2}$. El único equilibrio es $p=q=\tfrac{1}{2}$, con un pago esperado nulo para ambos.


\subsection{Juegos bayesianos y equilibrio bayesiano de Nash}

Con información incompleta, el equilibrio de NASH se aplica al juego que resulta de la transformación de HARSANYI. Este apartado define ese juego, el juego bayesiano, y su concepto de solución, justifica que exista, lo aplica al ejemplo del texto y formaliza la reinterpretación de las estrategias mixtas que el cuerpo anticipa.

Un \textbf{juego bayesiano} (HARSANYI, 1967-1968) se describe, en la notación de MAS-COLELL, WHINSTON y GREEN (1995), por los jugadores $i=1,\ldots,I$; un conjunto de tipos $\Theta_i$ para cada jugador; una distribución de probabilidad conjunta de los tipos, $F(\theta_1,\ldots,\theta_I)$, que es de conocimiento común; un conjunto de acciones $S_i$ para cada jugador; y unas funciones de pagos $u_i(s_i,s_{-i},\theta_i)$, que dependen de las acciones de todos y del tipo propio. La naturaleza elige el perfil de tipos $\theta=(\theta_1,\ldots,\theta_I)$ según $F$, cada jugador observa su tipo y todos eligen a la vez. Una estrategia pura del jugador $i$ es una \textbf{regla de decisión} que asigna una acción a cada tipo, $s_i:\Theta_i\rightarrow S_i$, y su pago esperado \textit{ex ante} con el perfil de reglas $(s_1(\cdot),\ldots,s_I(\cdot))$ es, en la formulación de MAS-COLELL, WHINSTON y GREEN (1995):
$$\tilde{u}_i\left(s_i(\cdot),s_{-i}(\cdot)\right)=E_\theta\left[u_i\left(s_i(\theta_i),s_{-i}(\theta_{-i}),\theta_i\right)\right]$$
Un perfil de reglas de decisión $(s_1(\cdot),\ldots,s_I(\cdot))$ es un \textbf{equilibrio bayesiano de Nash} si es un equilibrio de Nash del juego en forma normal con pagos $\tilde{u}_i$, es decir, si ninguna regla alternativa da a ningún jugador un pago esperado mayor. Cuando todos los tipos tienen probabilidad positiva, esa condición equivale a que cada tipo de cada jugador elija una mejor respuesta dadas sus creencias, condicionadas a su propio tipo, sobre los tipos de los demás (MAS-COLELL, WHINSTON y GREEN, 1995, proposición 8.E.1):
$$E_{\theta_{-i}}\left[u_i\left(s_i(\theta_i),s_{-i}(\theta_{-i}),\theta_i\right)\mid\theta_i\right]\geq E_{\theta_{-i}}\left[u_i\left(s_i',s_{-i}(\theta_{-i}),\theta_i\right)\mid\theta_i\right]\qquad\forall s_i'\in S_i,\ \forall\theta_i\in\Theta_i,\ \forall i$$
Si cada conjunto de tipos tiene un solo elemento, el juego bayesiano es un juego de información completa y el equilibrio bayesiano de NASH coincide con el equilibrio de NASH.

\paragraph{La forma de agentes y la existencia}

La condición anterior permite leer el equilibrio bayesiano de NASH como el equilibrio de NASH de un juego en el que cada par $(i,\theta_i)$ es un jugador distinto, el \textit{agente} del tipo $\theta_i$, cuyo pago es el esperado condicionado a su tipo. HARSANYI (1967-1968) lo llamó juego de SELTEN\textcolor{orange}{\textsuperscript{[Probable]}}. Si los conjuntos de jugadores, tipos y acciones son finitos, ese juego es finito y, por el teorema de NASH, tiene al menos un equilibrio en estrategias mixtas, que es un equilibrio bayesiano de NASH del juego original. Con tipos continuos, como en las subastas, la existencia exige condiciones adicionales sobre la distribución de los tipos y la continuidad de los pagos (MILGROM y WEBER, 1985\textcolor{orange}{\textsuperscript{[Probable]}}); en los modelos de subastas del anexo de diseño de mecanismos, el equilibrio se obtiene directamente, resolviendo la condición de mejor respuesta.

\paragraph{El ejemplo del texto}

En la versión simultánea del juego de entrada con información incompleta, $\Theta_1$ tiene un solo elemento, $\Theta_2=\{\theta_{ac},\theta_{ag}\}$ y la probabilidad a priori asigna $p$ a $\theta_{ac}$. Las acciones son $S_1=\{\text{Entrar},\text{No entrar}\}$ y $S_2=\{\text{Alto},\text{Bajo}\}$, y los pagos de J2 son $(a_2,b_2,c_2,d_2)$ si es acomodaticio y $(\hat{a}_2,\hat{b}_2,\hat{c}_2,\hat{d}_2)$ si es agresivo, con $a_2>b_2$, $c_2>d_2$, $\hat{b}_2>\hat{a}_2$ y $\hat{d}_2>\hat{c}_2$. Una regla de J2 es un par $\left(s_2(\theta_{ac}),s_2(\theta_{ag})\right)$, de modo que tiene cuatro. Cada tipo tiene una acción estrictamente dominante, y en todo equilibrio bayesiano de NASH la regla de J2 es $s_2^*(\theta_{ac})=\text{Alto}$ y $s_2^*(\theta_{ag})=\text{Bajo}$. Conocida esa regla, el pago esperado de J1 es $p\,a_1+(1-p)\,b_1$ si entra y $c_1$ si no entra, y la condición de equilibrio para J1 es la del cuerpo:
$$s_1^*=\begin{cases}\text{Entrar}&\text{si }1-p<q^*\\ \text{No entrar}&\text{si }1-p>q^*\end{cases}\qquad\text{con}\qquad q^*=\frac{a_1-c_1}{a_1-b_1}$$
(equilibrio bayesiano de NASH del ejemplo; si $1-p=q^*$, cualquier mezcla de J1 lo es). La probabilidad de guerra de precios, que en la función de pagos era una conjetura de J1, es ahora la probabilidad a priori del tipo agresivo, que es de conocimiento común.

\paragraph{Purificación de las estrategias mixtas}

HARSANYI (1973) demostró que, en casi todos los juegos finitos, todo equilibrio en estrategias mixtas es el límite de equilibrios bayesianos en estrategias puras de juegos cuyos pagos se perturban con una pequeña componente privada de cada jugador, cuando la perturbación tiende a cero. En el juego perturbado, cada tipo juega una estrategia pura y casi siempre estricta, y la mezcla aparente del equilibrio original es la distribución de esas estrategias puras que induce la distribución de los tipos. En el ejemplo con J2 agresivo, si los pagos de cada jugador incluyen una pequeña componente privada, cada uno entra, o baja el precio, según el valor de esa componente, y las probabilidades con que lo hace, vistas por el rival, tienden a $p^*$ y $q^*$ cuando la perturbación se anula. La estrategia mixta no describe a un jugador que sortea, sino la incertidumbre del rival sobre su tipo.

\subsection{Racionalidad secuencial}

En los juegos dinámicos, el equilibrio de NASH puede descansar en amenazas que nadie cumpliría, porque solo exige que las estrategias sean óptimas en la senda de equilibrio. Este apartado formaliza el problema, define el subjuego y el equilibrio de NASH perfecto en subjuegos de SELTEN (1965), presenta los métodos para hallarlo (la inducción hacia atrás de ZERMELO, 1913, y KUHN, 1953, y su generalización) y lo aplica a los juegos repetidos, finitos e infinitos.

\subsubsection{Amenazas no creíbles}

La versión secuencial del juego de entrada, con J2 acomodaticio ($a_2>b_2$ y $c_2>d_2$) y los pagos de partida de J1 ($a_1>c_1=d_1>b_1$), tiene tres equilibrios de NASH en estrategias puras, que se comprueban en la matriz $2\times4$ de su forma normal:
$$(\text{Entrar};\ \text{Alto},\text{Alto}),\qquad(\text{Entrar};\ \text{Alto},\text{Bajo}),\qquad(\text{No entrar};\ \text{Bajo},\text{Alto})$$
En el tercero, bajar el precio en $x_1$ no es óptimo para J2, porque $b_2<a_2$, pero $x_1$ no se alcanza, porque J1 no entra, y lo que J2 haga allí no afecta a su pago. En el segundo, bajar el precio en $x_2$ tampoco es óptimo, porque $d_2<c_2$, y $x_2$ no se alcanza, porque J1 entra. El equilibrio de NASH solo exige optimalidad donde el juego llega: como se formaliza en el apartado de creencias, $\sigma$ es un equilibrio de NASH si y solo si es secuencialmente racional en todo conjunto de información $H$ que se alcanza con probabilidad positiva, $P(H\mid\sigma)>0$, y deja libre lo que se hace en los demás. Las acciones de los nodos no alcanzados importan, sin embargo, porque sostienen las decisiones de los demás: en el tercer equilibrio, J1 no entra precisamente porque espera la guerra de precios.

El \textbf{principio de racionalidad secuencial} exige más: la estrategia de cada jugador debe especificar acciones óptimas desde cualquier punto del árbol en adelante, se alcance o no, dadas las estrategias de los demás. Una amenaza es \textbf{creíble} si llevarla a cabo es óptimo cuando llega el momento.

\subsubsection{Subjuegos}

Para aplicar el principio, SELTEN (1965) se apoyó en la noción de subjuego. En la notación de MAS-COLELL, WHINSTON y GREEN (1995), un \textbf{subjuego} de un juego en forma extensiva $\Gamma_E$ es un subconjunto del juego con dos propiedades:

\begin{lnum}
    \item Comienza en un conjunto de información que contiene un único nodo de decisión, contiene todos los nodos de decisión que le suceden, inmediatos o posteriores, y solo estos nodos;
    \item Si el nodo de decisión $x$ está en el subjuego, todo $x'\in H(x)$ está también, donde $H(x)$ es el conjunto de información que contiene a $x$; es decir, no hay conjuntos de información \textit{truncados}.
\end{lnum}

El juego entero es un subjuego; los demás se llaman \textit{propios}. El rasgo fundamental de un subjuego es que, visto aisladamente, constituye en sí mismo un juego, con sus jugadores, su información y sus pagos, de modo que puede exigírsele que se juegue de acuerdo con un equilibrio de NASH.

\begin{center}\begin{minipage}{\linewidth}\centering {\small\textbf{Subjuegos del juego de entrada}}\par\vspace{0.3em} \includegraphics{JE_anx_subjuegos.pdf}\par\vspace{0.3em} {\small\textit{Fuente: }Elaboración propia.}\end{minipage}\end{center}

En la versión secuencial del juego de entrada hay tres subjuegos: el juego entero y los dos que empiezan en $x_1$ y en $x_2$, cada uno con una sola decisión de J2. En la versión simultánea, el único subjuego es el juego entero: un subjuego que empezase en $x_1$ partiría el conjunto de información $\{x_1,x_2\}$, y la segunda condición lo impide. Lo mismo ocurre tras la transformación de HARSANYI: los dos nodos en que decide J1, uno por cada tipo de J2, forman un único conjunto de información, y ningún subjuego propio puede empezar en ellos. Por eso la perfección en subjuegos pierde fuerza con información imperfecta o incompleta, el problema que resuelven las creencias.

\subsubsection{Equilibrio de Nash Perfecto en Subjuegos}

Un perfil de estrategias $\sigma=(\sigma_1,\ldots,\sigma_I)$ de un juego en forma extensiva de $I$ jugadores es un \textbf{equilibrio de NASH perfecto en subjuegos} (ENPS) si induce un equilibrio de NASH en cada subjuego del juego (SELTEN, 1965, en la notación de MAS-COLELL, WHINSTON y GREEN, 1995).\footnote{%
  Obsérvese que todo Equilibrio de Nash Perfecto en Subjuegos (ENPS) es un Equilibrio de Nash (EN) (dado que el juego en su conjunto es un subjuego), pero no todo EN es perfecto en subjuegos.
} En la versión secuencial del ejemplo con J2 acomodaticio, el único equilibrio del subjuego que empieza en $x_1$ es Alto ($a_2>b_2$), y el del que empieza en $x_2$, también Alto ($c_2>d_2$); de los tres equilibrios de NASH, solo (Entrar; Alto, Alto) induce un equilibrio en los dos, y es el único ENPS. Con J2 agresivo ($b_2>a_2$), el del subjuego de $x_1$ es Bajo, y el ENPS es (No entrar; Bajo, Alto): la misma amenaza es ahora creíble. En la versión simultánea, el único subjuego es el juego entero, y los ENPS coinciden con los equilibrios de NASH.

\paragraph{Inducción hacia atrás}

El primer resultado sobre la racionalidad secuencial es el teorema de ZERMELO (1913), que demostró, para el ajedrez, que en un juego finito de suma cero con información perfecta uno de los dos jugadores puede forzar la victoria o ambos pueden forzar las tablas. KUHN (1953) lo generalizó a cualquier juego finito de información perfecta.

\textbf{Teorema} (ZERMELO, 1913; KUHN, 1953).- Todo juego finito de información perfecta tiene un equilibrio de NASH perfecto en subjuegos en estrategias puras, que puede obtenerse por inducción hacia atrás. Si ningún jugador obtiene el mismo pago en dos nodos terminales distintos, ese ENPS es único.

La demostración es constructiva y procede por inducción sobre el número de nodos de decisión. En cada nodo cuyos sucesores son todos terminales, el jugador que mueve elige una acción que maximice su pago, que existe porque el número de acciones es finito; se sustituye el nodo por el pago resultante, y queda un juego finito de información perfecta con menos nodos de decisión, al que se aplica el mismo paso. El perfil así construido induce un equilibrio de NASH en cada subjuego, porque en cada uno se ha construido de la misma manera, desde el final. Si no hay empates, cada maximización tiene una única solución, y el ENPS es único.

\begin{center}\begin{minipage}{\linewidth}\centering {\small\textbf{Inducción hacia atrás en tres etapas}}\par\vspace{0.3em} \includegraphics{JE_anx_retroinduccion.pdf}\par\vspace{0.3em} {\small\textit{Fuente: }Elaboración propia.}\end{minipage}\end{center}

La figura extiende el ejemplo con una tercera etapa para ver qué significa que la estrategia especifique acciones en nodos que no se alcanzan, también para el propio jugador. Si J2 baja el precio tras la entrada, J1 decide si se queda, con desenlace $(b_1,b_2)$, o sale del mercado, con desenlace $(e_1,e_2)$. Supongamos que, en plena guerra de precios, a J1 le conviene salir ($e_1>b_1$), pero que entrar para salir después le deja peor que no haber entrado ($c_1>e_1$). La inducción hacia atrás empieza por esa última decisión: J1 sale. J2, que lo anticipa, compara en $x_1$ acomodarse, $a_2$, con desatar una guerra que expulsa al entrante, $e_2$.

En el panel (a), $a_2>e_2$: J2 se acomoda, J1 entra, y el ENPS es (Entrar, Salir; Alto, Alto). La senda de equilibrio es Entrar-Alto, pero la estrategia de J1 especifica también qué haría en el nodo de la guerra, que no se alcanza, y la de J2, qué precio fijaría sin competidor; sin esas acciones, el perfil no sería una especificación completa del equilibrio, porque es lo que J1 haría si llegara a la guerra lo que hace que a J2 no le compense provocarla. En el panel (b), $e_2>a_2$: la guerra le compensa a J2, porque expulsa al entrante y le devuelve el mercado, J1 lo anticipa y no entra, y el ENPS es (No entrar, Salir; Bajo, Alto). Es la lógica de la \textbf{predación}: una guerra de precios es creíble si el beneficio de expulsar al rival compensa su coste.

\paragraph{Inducción hacia atrás generalizada}

El problema de la inducción hacia atrás es que solo sirve para juegos con información perfecta. Si hay conjuntos de información con varios nodos, ir viendo en cada nodo qué acción es óptima deja de ser factible, porque el jugador no sabe en qué nodo está. El procedimiento se generaliza sustituyendo cada subjuego final por los pagos de uno de sus equilibrios de NASH (\textit{generalized backward induction procedure}, MAS-COLELL, WHINSTON y GREEN, 1995):

\begin{lnum}
    \item Empezar al final del \textit{árbol} de juego, e identificar el Equilibrio de Nash para cada subjuego final (esto es, todos los que no tienen otro subjuego dentro suyo);
    \item Escoger un Equilibrio de Nash de cada subjuego final y derivar la reducción de la forma extensiva en la que estos subjuegos finales son reemplazados por los pagos que resultan en estos Equilibrios de Nash cuando los jugadores utilizan esas estrategias de equilibrio;
    \item Repetir las acciones anteriores para el juego reducido. Continuar el proceso hasta que cualquier movimiento en el Juego de forma extensiva esté determinado. Esta colección de acciones en los diferentes conjuntos de información del Juego constituye el perfil de estrategias de ENPS;
    \item Si no encontramos algún equilibrio múltiple en alguno de los pasos, este perfil de estrategias es el único ENPS. Si encontramos algún equilibrio, el conjunto total de ENPS se identifica repitiendo el proceso para cada posible equilibrio que pueda resultar del subjuego en cuestión.
\end{lnum}

Este proceso de retroinducción generalizado aplica tanto a juegos de información perfecta como a juegos de información imperfecta. Como todo juego finito tiene un equilibrio de NASH, al menos en estrategias mixtas, cada subjuego final tiene uno, y el procedimiento demuestra que todo juego finito en forma extensiva tiene al menos un ENPS, quizá en estrategias mixtas.

\begin{center}\begin{minipage}{\linewidth}\centering {\small\textbf{Inducción hacia atrás generalizada: competencia en precios tras la entrada}}\par\vspace{0.3em} \includegraphics{JE_anx_inducgen.pdf}\par\vspace{0.3em} {\small\textit{Fuente: }Elaboración propia a partir de MAS-COLELL, WHINSTON y GREEN (1995).}\end{minipage}\end{center}

La figura lo aplica a otra variante del ejemplo. Si J1 entra, las dos empresas fijan el precio a la vez; si no entra, J2 sigue como monopolista y el desenlace es $(c_1,c_2)$. En el subjuego que sigue a la entrada, a cada empresa le conviene igualar el precio de la otra ($a_1>h_1$ y $b_1>g_1$ para J1; $a_2>g_2$ y $b_2>h_2$ para J2), de modo que tiene dos equilibrios en estrategias puras: precios altos, $(a_1,a_2)$, y guerra de precios, $(b_1,b_2)$.\footnote{%
  Tiene además un tercero en estrategias mixtas, que da lugar a un tercer ENPS.
} Sustituido el subjuego por cada uno de ellos, J1 resuelve dos juegos reducidos. Si espera precios altos, entra, porque $a_1>c_1$; si espera la guerra, no entra, porque $c_1>b_1$. Hay, por tanto, dos ENPS en estrategias puras: (Entrar, Alto; Alto) y (No entrar, Bajo; Bajo). En el segundo, la amenaza de guerra sí es creíble, porque es un equilibrio del subjuego que sigue a la entrada: si J1 entrara y esperase la guerra, bajar el precio sería su mejor respuesta, y la de J2. La perfección en subjuegos descarta las amenazas que no son equilibrio de lo que queda de juego, pero no elige entre los equilibrios de un subjuego.

\subsubsection{Juegos repetidos}

Un \textbf{juego repetido} consiste en jugar varias veces un mismo juego, el \textit{juego de etapa}, observando cada vez lo que se jugó en las anteriores. Como cada historia da lugar a un subjuego, la perfección en subjuegos es el concepto natural para analizarlos, y la pregunta es si la repetición permite sostener resultados que no son equilibrio del juego de etapa.

\paragraph{Horizonte finito}

\textbf{Proposición} (MAS-COLELL, WHINSTON y GREEN, 1995).- Sea un juego de $I$ jugadores en forma extensiva que consiste en jugar $T$ veces un juego de acciones simultáneas, en el que los jugadores observan las estrategias puras jugadas en cada etapa inmediatamente después de su conclusión, y en el que el pago de cada jugador es la suma de sus pagos en las $T$ etapas. Si el juego de etapa tiene un único equilibrio de NASH, el juego repetido tiene un único ENPS, que consiste en que cada jugador juegue en cada etapa su estrategia de ese equilibrio, independientemente de lo que haya ocurrido antes.

La idea principal subyacente es la aplicación de una consecuencia de la lógica de retroinducción: el movimiento en el último juego debe resultar en el único Equilibrio de Nash de ese juego porque, en ese punto, los jugadores esencialmente se enfrentan únicamente a ese juego. Pero si el movimiento de ese juego está predeterminado, cuando los jugadores juegan el juego inmediatamente antecesor sería equivalente a estar jugando solo ese juego, y así sucesivamente. Una implicación importante es la manera en la que los ENPS eliminan la interdependencia histórica de las estrategias en la clase de juegos considerados en este ejemplo. En general, la estrategia de un jugador podría potencialmente prometer pagos futuros o \textit{castigos} a otros jugadores si toman acciones / movimientos concretos en momentos iniciales del Juego. Pero mientras cada componente del Juego (paso, subjuego, etc.) tenga un único EN, estrategias que conducen a ENPS no pueden ser históricamente dependientes.

La misma lógica produce la \textbf{paradoja de la cadena de tiendas} de SELTEN (1978). Si el juego de entrada secuencial con J2 acomodaticio se juega sucesivamente en $T$ mercados, con un entrante distinto en cada uno, la inducción hacia atrás implica que J2 se acomoda en el último mercado, que nada puede ganar luchando en el penúltimo para disuadir una entrada que en el último no va a disuadir, y así hasta el primero: el único ENPS es que todos entren y J2 se acomode siempre, por más que la intuición sugiera que le convendría luchar al principio para ganarse una reputación. La paradoja se resuelve con información incompleta, como se ve al tratar las creencias. Si el juego de etapa tiene varios equilibrios, la conclusión cambia: amenazar con pasar a un equilibrio peor en las últimas etapas puede sostener resultados que no son equilibrio en las primeras (BENOIT y KRISHNA, 1985).

\paragraph{Horizonte infinito}

Si el juego de etapa se repite indefinidamente, se supone que los jugadores descuentan el futuro con un factor $\delta\in(0,1)$ y se mide el pago de cada uno por la media descontada de sus pagos de etapa:
$$U_i=(1-\delta)\sum_{t=0}^{\infty}\delta^t\,u_i\!\left(a^t\right)$$
(pago medio descontado de un juego repetido infinitamente), donde $a^t$ es el perfil de acciones de la etapa $t$ y el factor $(1-\delta)$ expresa el pago en las mismas unidades que el de una etapa. Una estrategia asigna una acción a cada historia posible, y tras cada historia empieza un subjuego, que es idéntico al juego entero salvo por el pasado. No hay una última etapa desde la que empezar la inducción hacia atrás, pero comprobar la perfección en subjuegos se simplifica con el \textbf{principio de la desviación en un periodo}: en un juego repetido con descuento y pagos de etapa acotados, un perfil es un ENPS si y solo si ningún jugador puede ganar, tras ninguna historia, desviándose en un único periodo y siguiendo después su estrategia (FUDENBERG y TIROLE, 1991).

El \textbf{dilema del prisionero} lo ilustra. Sea $T>R>P>S$, con $2R>T+S$, y sea la estrategia del \textit{gatillo}: cooperar mientras nadie se haya desviado y no cooperar para siempre en cuanto alguien lo haga. Tras una desviación, los jugadores juegan para siempre el equilibrio de la etapa, y nadie gana desviándose. En la fase de cooperación, desviarse un periodo da $T$ hoy y $P$ para siempre después, de modo que no compensa si
$$R\geq(1-\delta)\,T+\delta\,P\qquad\Longleftrightarrow\qquad\delta\geq\frac{T-R}{T-P}$$
(cooperación con estrategias de gatillo en el dilema del prisionero repetido, FRIEDMAN, 1971). Si los jugadores son lo bastante pacientes, la cooperación es un ENPS, aunque en el juego de etapa no lo sea.

El resultado se generaliza en los \textbf{teoremas de tradición oral} (\textit{folk theorems}). Sea $V$ el conjunto de pagos \textit{factibles}, la envolvente convexa de los pagos de etapa, y sea $\underline{v}_i=\min_{\alpha_{-i}}\max_{a_i}u_i(a_i,\alpha_{-i})$ el pago de \textit{minimax} del jugador $i$, el mínimo al que los demás pueden limitarle si él responde de forma óptima. Un pago $v\in V$ es \textit{estrictamente individualmente racional} si $v_i>\underline{v}_i$ para todo $i$. FRIEDMAN (1971) demostró que cualquier pago factible que mejore el de un equilibrio de NASH de la etapa para todos los jugadores puede sostenerse como ENPS si $\delta$ es suficientemente próximo a $1$, con la amenaza de volver para siempre a ese equilibrio. FUDENBERG y MASKIN (1986) extendieron el resultado a todos los pagos factibles y estrictamente individualmente racionales, con una condición de dimensión cuando hay más de dos jugadores. Sin descuento, los resultados equivalentes se deben a AUMANN y SHAPLEY (1976)\textcolor{orange}{\textsuperscript{[Probable]}} y RUBINSTEIN (1979)\textcolor{orange}{\textsuperscript{[Probable]}}.

\begin{center}\begin{minipage}{\linewidth}\centering {\small\textbf{Pagos factibles e individualmente racionales del dilema del prisionero}}\par\vspace{0.3em} \includegraphics{JE_anx_folk.pdf}\par\vspace{0.3em} {\small\textit{Fuente: }Elaboración propia a partir de FUDENBERG y MASKIN (1986).}\end{minipage}\end{center}

En el dilema del prisionero, el minimax de cada jugador es $P$, el pago del equilibrio de la etapa, y los dos teoremas coinciden: cualquier punto de la zona sombreada, desde la cooperación plena $(R,R)$ hasta repartos muy asimétricos, puede sostenerse si los jugadores son pacientes. Es la base de la teoría de la colusión tácita ([\authorfont{Ver Tema 3.A.19}]), y también su límite como predicción: la repetición no elimina la multiplicidad de equilibrios, sino que la multiplica.

\subsection{Creencias y racionalidad secuencial}

A pesar de que el equilibrio de NASH perfecto en subjuegos es muy útil al capturar el \textit{principio de racionalidad secuencial}, a menudo no es suficiente. Cuando un jugador mueve sin saber en qué nodo está, la optimalidad de su acción depende de lo que crea sobre lo ocurrido, y si el único subjuego es el juego entero, la perfección en subjuegos no dice nada. Este apartado introduce las creencias como parte del equilibrio,\footnote{%
    Es muy importante este concepto pues ahora, al hablar de \textit{equilibrio} no estaremos hablando solo de un perfil de estrategias, sino de un \textit{equilibrio} visto como un conjunto de perfil de estrategias y de sistemas de creencias.
} define el equilibrio bayesiano perfecto débil y sus refuerzos, el equilibrio bayesiano perfecto y el equilibrio secuencial, presenta el equilibrio perfecto de mano temblorosa y aplica todo ello a los juegos de señalización y a la reputación.

\subsubsection{Cuando la perfección en subjuegos no basta}

Sea la variante del juego de entrada en la que J1 puede entrar de dos maneras, imitando el producto de J2 o diferenciándose de él, y J2 observa que J1 ha entrado, pero no de qué manera: sus dos nodos, $x_I$ tras la imitación y $x_D$ tras la diferenciación, forman un único conjunto de información, $H_2=\{x_I,x_D\}$. Si J1 no entra, el desenlace es $(c_1,c_2)$; tras imitar, el precio alto y el bajo dan $(a_1,a_2)$ y $(b_1,b_2)$; tras diferenciarse, $(e_1,e_2)$ y $(f_1,f_2)$. Supongamos que a J1 no le compensa entrar si hay guerra de precios ($c_1>b_1$ y $c_1>f_1$), pero sí si J2 se acomoda, al menos con una de las dos maneras de entrar, y que J2 prefiere acomodarse sea cual sea la manera en que J1 haya entrado ($a_2>b_2$ y $e_2>f_2$).

\begin{center}\begin{minipage}{\linewidth}\centering {\small\textbf{Entrada de dos maneras: una amenaza que la perfección en subjuegos no descarta}}\par\vspace{0.3em} \includegraphics{JE_ebp_creencias.pdf}\par\vspace{0.3em} {\small\textit{Fuente: }Elaboración propia a partir de MAS-COLELL, WHINSTON y GREEN (1995).}\end{minipage}\end{center}

El perfil (No entrar, Bajo) del panel (a) es un equilibrio de NASH: si J1 espera la guerra de precios, no entra, y como no entra, a J2 le da igual lo que haga en $H_2$. Como el único subjuego es el juego entero, es también perfecto en subjuegos. Pero la amenaza no es creíble: en $H_2$, bajar el precio no es óptimo para ninguna creencia que J2 pueda tener sobre la manera en que J1 ha entrado. ¿Cómo descartar entonces este equilibrio? Una posibilidad, en el corazón del principio de racionalidad secuencial, es exigir que la acción de J2 tras la entrada sea óptima para alguna creencia, y que esa creencia sea razonable.

\subsubsection{Sistemas de creencias y racionalidad secuencial}

Un \textbf{sistema de creencias} $\mu$ en un juego en forma extensiva, $\Gamma_E$, es una especificación de una probabilidad $\mu(x)\in[0,1]$ para cada nodo de decisión $x$ tal que (MAS-COLELL, WHINSTON y GREEN, 1995):
$$\sum_{x\in H}\mu(x)=1\qquad\forall H\in\mathcal{H}$$
Un sistema de creencias puede entenderse como la especificación, para cada conjunto de información, de las probabilidades que el jugador que mueve en él asigna a cada uno de sus nodos, condicionadas a que la partida haya llegado a ese conjunto de información. En un conjunto de información de un único nodo, la creencia es evidentemente uno. Un equilibrio pasa a ser un par $(\sigma,\mu)$, una \textbf{evaluación}, que dice qué van a jugar los jugadores y qué creen sobre lo ocurrido en cada conjunto de información con varios nodos.

Un perfil estratégico $\sigma=(\sigma_1,\sigma_2,\ldots,\sigma_I)$ en forma extensiva, $\Gamma_E$, es secuencialmente racional en el conjunto de información $H$ dado un sistema de creencias $\mu$ si, denotando $\iota(H)$ el jugador que mueve en $H$:
$$E\left[u_{\iota(H)}\mid H,\mu,\sigma_{\iota(H)},\sigma_{-\iota(H)}\right]\geq E\left[u_{\iota(H)}\mid H,\mu,\tilde{\sigma}_{\iota(H)},\sigma_{-\iota(H)}\right]$$
$\forall\,\tilde{\sigma}_{\iota(H)}\in\Delta(S_{\iota(H)})$. Si el perfil estratégico $\sigma$ satisface la condición para todo conjunto de información $H$, entonces decimos que $\sigma$ es secuencialmente racional dado el sistema de creencias $\mu$.

En palabras, un perfil estratégico es secuencialmente racional si ningún jugador encuentra beneficioso, una vez que se ha alcanzado un conjunto de información, revisar su estrategia dadas sus creencias sobre lo que ya ha ocurrido y las estrategias de sus rivales. La \textit{respuesta óptima} pasa a ser, por tanto, una respuesta óptima dadas las estrategias rivales y las creencias en cada conjunto de información del juego.

\begin{center}\begin{minipage}{\linewidth}\centering {\small\textbf{Pago esperado de J2 tras la entrada según su creencia}}\par\vspace{0.3em} \includegraphics{JE_anx_umbral_mu.pdf}\par\vspace{0.3em} {\small\textit{Fuente: }Elaboración propia.}\end{minipage}\end{center}

En el ejemplo, sea $\mu$ la probabilidad que J2 asigna a $x_I$, la imitación, y $1-\mu$ la que asigna a $x_D$. Mantener el precio le da $\mu\,a_2+(1-\mu)\,e_2$, y bajarlo, $\mu\,b_2+(1-\mu)\,f_2$. En el panel (a) de la figura, con $a_2>b_2$ y $e_2>f_2$, mantener el precio es mejor para cualquier $\mu$, y Bajo no es secuencialmente racional con ninguna creencia. En el panel (b), si J2 prefiere la guerra frente a un imitador ($b_2>a_2$) pero no frente a quien se diferencia ($e_2>f_2$), bajar el precio es óptimo si y solo si
$$\mu\geq\bar{\mu}=\frac{e_2-f_2}{(b_2-a_2)+(e_2-f_2)}$$
(racionalidad secuencial de J2 en la variante del juego de entrada con dos maneras de entrar), y la credibilidad de la amenaza depende de lo que J2 crea.

\subsubsection{Equilibrio Bayesiano Perfecto Débil (EBPD)}

Este tipo de equilibrio extiende el \textit{principio de racionalidad secuencial} introduciendo formalmente la noción de \textit{creencias}\footnote{%
    El concepto de Equilibrio Bayesiano Perfecto fue inicialmente desarrollado para capturar los requisitos de la racionalidad secuencial en juegos dinámicos con información imperfecta (esto es, en juegos dinámicos bayesianos). El EBPD es una variante que introducimos por simplicidad.
}; MYERSON (1991) se refiere a él como equilibrio secuencial débil. Exige la conjunción de dos condiciones:

\begin{lnum}
    \item Las estrategias deben ser secuencialmente racionales\footnote{%
    Otra definición de \textit{secuencialmente racionales} es: En cada conjunto de información, las acciones tomadas por el Jugador con el turno (y sus acciones posteriores) deben ser óptimas dadas las creencias del Jugador en ese conjunto de información y las estrategias rivales subsiguientes.
} dado un sistema de creencias;
    \item Cuando sea posible, las creencias deben ser consistentes con las estrategias\footnote{%
    Esta idea es prácticamente la misma que en el Equilibrio de Nash: en cualquier equilibrio, los jugadores deben tener creencias correctas sobre las elecciones estratégicas de sus oponentes. Para motivar la consistencia específica requerida sobre las creencias, consideremos cómo podríamos definir la noción de creencias consistentes en el caso especial en el cual el equilibrio de cada Jugador asigna una probabilidad estrictamente positiva a cada posible acción en cada conjunto de información. En este caso, cada conjunto de información se alcanza con probabilidad positiva.
}.
\end{lnum}

La noción de que las creencias sean consistentes con el perfil de estrategias es muy sencilla cuando un conjunto de información se alcanza con probabilidad positiva: para cada nodo de decisión del conjunto de información, el jugador calcula la probabilidad de alcanzar dicho nodo dado el perfil de estrategias, y asigna a cada nodo la probabilidad condicionada a que el juego haya llegado a ese conjunto de información, con la regla de Bayes:
$$P(x\mid H,\sigma)=\frac{P(x\mid\sigma)}{\sum_{x'\in H}P(x'\mid\sigma)}$$

\textbf{NOTA.-} Esta definición nos proporciona indirectamente una definición que será de gran utilidad más adelante: dado un equilibrio\footnote{%
    Puede ser EN, EBN, ENPS, cualquier tipo.
}, un conjunto de información \textbf{\textit{estará en la ruta de equilibrio}} si, y solo si, va a ser alcanzado con probabilidad mayor que cero si el Juego se desarrolla de acuerdo con el perfil de estrategias de equilibrio (que hemos hallado); y estará fuera del equilibrio si sabemos con certeza ($P=1$) que no va a ser alcanzado si el Juego se desarrolla de acuerdo con la estrategia de equilibrio. Esto es consistencia de las creencias.

\paragraph{EBPD: Definición}

Por tanto, un perfil de estrategias, $\sigma$, y un sistema de creencias, $\mu$, es un Equilibrio de Bayes Perfecto Débil, en un juego en forma extensiva, si tiene las siguientes propiedades (MAS-COLELL, WHINSTON y GREEN, 1995):

\begin{lnum}
    \item El perfil de estrategias es secuencialmente racional dado un sistema de creencias, $\mu$;
    \item El sistema de creencias, $\mu$, se obtiene de ese perfil de estrategias a través de la Regla de Bayes, siempre y cuando sea posible\footnote{%
    Otra formulación de este requisito es: En los conjuntos de información de la ruta de equilibrio (que estamos analizando), las creencias deben poderse determinar por la Regla de Bayes sobre las estrategias de Equilibrio de los jugadores.
}. Esto es, para cada conjunto de información $H$ tal que la probabilidad dado el perfil de estrategias, $\sigma$, es mayor que cero (probabilidad condicionada), debemos tener que\footnote{%
    A menudo encontraremos que se expresa simplemente $\sigma$ al hablar de un EBPD, sin embargo, siempre debemos entender este como asociado a un conjunto de creencias que satisfacen la definición formal proporcionada.
}:
\end{lnum}
$$\mu(x)=\frac{P(x\mid\sigma)}{P(H\mid\sigma)},\qquad\forall x\in H$$

La definición permite relacionar el EBPD con el equilibrio de NASH. Un perfil de estrategias, $\sigma$, es un EN de un juego en forma extensiva, $\Gamma_E$, si y solo si existe un sistema de creencias $\mu$ tal que:
\begin{enumerate}
    \item[(i)] El perfil de estrategias $\sigma$ es secuencialmente racional dado un sistema de creencias $\mu$ en todo conjunto de información $H$ t.q. $P(H\mid\sigma)>0$;
    \item[(ii)] El sistema de creencias $\mu$ se obtiene del perfil de estrategias $\sigma$ a través de la Regla de Bayes, siempre que sea posible.
\end{enumerate}
La diferencia con el EBPD está, pues, en la primera condición: el equilibrio de NASH solo exige racionalidad secuencial en la senda de equilibrio; el EBPD la exige en todos los conjuntos de información.

\paragraph{El ejemplo}

En la variante del juego de entrada, la evaluación (Imitar, Bajo; $\mu=0$) no cumple la segunda condición: $H_2$ se alcanza con probabilidad $1$, y la regla de Bayes exige
$$\mu=\frac{P(x_I\mid\sigma)}{P(H_2\mid\sigma)}=\frac{1}{1}=1,$$
de modo que la creencia $\mu=0$ contradice lo que J2 sabe de la estrategia de J1: jugaría su acción considerando que J1 ha imitado y, a la vez, pensaría que no está en ese nodo. En la evaluación (No entrar, Bajo; $\mu$), en cambio, $H_2$ no se alcanza, $P(H_2\mid\sigma)=0$, la regla de Bayes da $0/0$ y no dice nada: \textit{a priori}, cualquier $\mu\in[0,1]$ satisface la segunda condición. Pero la primera se exige también fuera de la senda, y con los pagos del panel (a) ninguna creencia hace óptimo bajar el precio: la evaluación no es un EBPD. El único EBPD es que J1 entre de la manera que más le convenga, que J2 se acomode y que $\mu$ sea la que la regla de Bayes deduzca de la manera de entrar. El EBPD descarta así la amenaza que la perfección en subjuegos no descartaba.

Con los pagos del panel (b), en cambio, (No entrar, Bajo; $\mu$) es un EBPD para cualquier $\mu\geq\bar{\mu}$: como la creencia fuera de la senda es libre, basta suponer que J2 cree que J1 ha imitado. Es la debilidad del concepto. Si imitar fuese una mala idea para J1 en cualquier caso, por ejemplo porque $c_1>a_1$, de modo que no ganaría imitando ni aunque J2 se acomodase, esa creencia sería poco razonable: J2 atribuiría la entrada a la jugada que menos le conviene a J1. Descartar creencias así es la idea de la \textit{inducción hacia adelante}, que formaliza el criterio intuitivo. Además, como las creencias fuera de la senda son libres incluso dentro de un subjuego propio, un EBPD puede no ser perfecto en subjuegos (MAS-COLELL, WHINSTON y GREEN, 1995).

\subsubsection{Fortalecimiento del EBPD: Equilibrio Bayesiano Perfecto}

Nos hemos referido al equilibrio anterior como \textit{débil} porque los requisitos de consistencia impuestos sobre las creencias son mínimos: el único requisito, más allá de que sean probabilidades y sumen $1$ dentro de cada conjunto de información, es que sean consistentes con las estrategias de equilibrio en la senda de equilibrio (\textit{equilibrium path}), en el sentido de que se obtengan de ellas con la regla de Bayes. Ninguna restricción se pone sobre las creencias en los conjuntos de información que se alcanzan con probabilidad nula.

El \textbf{equilibrio bayesiano perfecto} (EBP) incorpora una restricción adicional: la regla de Bayes debe aplicarse también, cuando sea posible, fuera de la senda de equilibrio. Si un jugador se desvía y el juego llega a un conjunto de información que el equilibrio no alcanzaba, las creencias de quien mueve allí deben poder recuperarse con la regla de Bayes a partir de las estrategias de equilibrio de lo que queda de juego, como si la desviación ya se hubiera producido. FUDENBERG y TIROLE (1991)\textcolor{orange}{\textsuperscript{[Probable]}} lo formulan para los juegos en varias etapas con acciones observables y tipos independientes, con tres exigencias: que las creencias sobre el tipo de cada jugador se actualicen con la regla de Bayes siempre que sea posible, también tras una desviación; que la desviación de un jugador no altere lo que los demás creen sobre el tipo de un tercero; y que todos compartan la misma creencia sobre el tipo de quien se ha desviado. Con estas restricciones, todo EBP es perfecto en subjuegos.

\subsubsection{Equilibrio Secuencial}

Un concepto estrechamente relacionado, desarrollado por KREPS y WILSON (1982a), es el Equilibrio Secuencial. En contraste con las nociones de EBP, el equilibrio secuencial introduce las restricciones de consistencia indirectamente, a través del límite de sucesiones de estrategias.

En esencia, el equilibrio secuencial requiere que las creencias se justifiquen por venir de estrategias mixtas completas, que asignan probabilidad positiva a todas las acciones, cercanas a las de equilibrio. Puede verse como si exigiésemos a los jugadores que justifiquen sus creencias con una historia en la que, con una probabilidad pequeña, los jugadores se equivocan al elegir sus estrategias\footnote{%
    Nótese que todo Equilibrio Secuencial es un EBPD porque el salto al límite coincide exactamente con las creencias obtenidas de las estrategias de equilibrio a través de la Regla de Bayes. Generalmente, lo contrario no es cierto.
}.

Un perfil de estrategias y un sistema de creencias $(\sigma,\mu)$ es un Equilibrio Secuencial de un Juego, $\Gamma_E$, si tiene las siguientes propiedades (KREPS y WILSON, 1982a, en la notación de MAS-COLELL, WHINSTON y GREEN, 1995):
\begin{enumerate}
    \item[(i)] El perfil de estrategias es secuencialmente racional, dado un sistema de creencias;
    \item[(ii)] Existe una secuencia completa de estrategias mixtas $\{\sigma^k\}_{k=1}^{\infty}$, con $\lim_{k\to\infty}\sigma^k=\sigma$, t.q. $\mu=\lim_{k\to\infty}\mu^k$, donde $\mu^k$ denota las creencias obtenidas desde $\sigma^k$ a través de la Regla de Bayes.
\end{enumerate}
La condición (ii) se llama \textbf{consistencia}. KREPS y WILSON (1982a) demostraron que todo juego finito tiene al menos un equilibrio secuencial, y todo equilibrio secuencial es perfecto en subjuegos y bayesiano perfecto débil.

En la variante del juego de entrada, la consistencia no restringe $\mu$ en la evaluación (No entrar, Bajo; $\mu$): si J1 se equivoca e imita con probabilidad $\varepsilon_I^k$ y se diferencia con probabilidad $\varepsilon_D^k$, la regla de Bayes da $\mu^k=\varepsilon_I^k/(\varepsilon_I^k+\varepsilon_D^k)$, y eligiendo cómo tienden a cero los dos errores puede obtenerse cualquier límite en $[0,1]$. El equilibrio secuencial descarta, por tanto, la amenaza del panel (a), por racionalidad secuencial, pero no la del panel (b), y para eliminar esta hace falta restringir qué errores son razonables, que es lo que hace la inducción hacia adelante.

A pesar de que este concepto restringe las creencias que están fuera de la \textit{ruta de equilibrio} suficientemente como para solventar los problemas del EBPD, en algunos casos resulta demasiado restrictivo. Por ejemplo, implica que cualquier dupla de jugadores que comparten el mismo conjunto de información tengan exactamente el mismo sistema de creencias dada una desviación de otro Jugador que haya causado que el Juego alcance ese nodo de decisión.

\subsubsection{Equilibrio perfecto de mano temblorosa}

SELTEN (1975) propuso refinar el equilibrio de NASH exigiendo que sea robusto a pequeños errores de los demás. Un perfil $\sigma$ es un \textbf{equilibrio perfecto de mano temblorosa} si existe una sucesión de perfiles de estrategias completamente mixtas $\{\sigma^k\}$, con $\sigma^k\to\sigma$, tal que, para todo jugador $i$ y todo $k$, $\sigma_i$ es una mejor respuesta a $\sigma_{-i}^k$:
$$\sigma_i\in MR_i\!\left(\sigma_{-i}^k\right)\qquad\forall i,\ \forall k$$
(equilibrio perfecto de mano temblorosa de SELTEN, 1975, en forma normal). Cada jugador ha de jugar una mejor respuesta no solo a las estrategias de equilibrio de los demás, sino también a sus versiones perturbadas, en las que cualquier estrategia se juega con una probabilidad pequeña pero positiva. Todo juego finito tiene al menos uno, y todo equilibrio perfecto es de NASH (SELTEN, 1975). Como los demás juegan todas sus estrategias con probabilidad positiva, una estrategia débilmente dominada nunca es mejor respuesta a la perturbación: un equilibrio perfecto no usa estrategias débilmente dominadas, y en los juegos de dos jugadores los equilibrios perfectos son exactamente los equilibrios de NASH en estrategias no débilmente dominadas.

En la forma normal $2\times4$ de la versión secuencial del ejemplo, con J2 acomodaticio, dos de los tres equilibrios de NASH usan estrategias débilmente dominadas, (Bajo, Alto) y (Alto, Bajo), y no son perfectos; el único equilibrio perfecto es (Entrar; Alto, Alto), que coincide con el ENPS: si J1 entra con una probabilidad positiva, aunque sea pequeña, J2 ha de responder de forma óptima en $x_1$. Aplicado a la \textit{forma de agentes} del juego en forma extensiva, en la que cada conjunto de información de un jugador es un agente que tiembla de forma independiente, el concepto da el equilibrio perfecto en forma extensiva, y KREPS y WILSON (1982a) demostraron que todo equilibrio perfecto en forma extensiva es secuencial.

\subsubsection{Juegos de señalización}

Un \textbf{juego de señalización} tiene dos jugadores: un emisor informado, que conoce su tipo $\theta\in\Theta$, distribuido según una probabilidad a priori $p(\theta)$ de conocimiento común, y elige un mensaje $m\in M$; y un receptor, que observa el mensaje, pero no el tipo, y elige una acción $a\in A$. Los pagos son $u_E(m,a,\theta)$ para el emisor y $u_R(m,a,\theta)$ para el receptor. Una evaluación $(m^*(\cdot),a^*(\cdot),\mu)$, con $\mu(\theta\mid m)$ la creencia del receptor tras observar $m$, es un equilibrio bayesiano perfecto si (i) cada tipo elige el mensaje que maximiza $u_E\left(m,a^*(m),\theta\right)$; (ii) tras cada mensaje, el receptor elige la acción que maximiza $\sum_\theta\mu(\theta\mid m)\,u_R(m,a,\theta)$; y (iii) para todo mensaje que algún tipo envía en equilibrio, las creencias se obtienen con la regla de Bayes:
$$\mu(\theta\mid m)=\frac{p(\theta)\,\mathbb{1}\{m^*(\theta)=m\}}{\sum_{\theta'}p(\theta')\,\mathbb{1}\{m^*(\theta')=m\}}$$
(creencias de equilibrio de un juego de señalización con estrategias puras). Para los mensajes que nadie envía, la creencia es libre. El equilibrio es \textbf{separador} si tipos distintos envían mensajes distintos, de modo que el mensaje revela el tipo; \textbf{agrupador}, si todos envían el mismo y no se aprende nada; y \textbf{semiseparador} o híbrido, si algún tipo aleatoriza.

\paragraph{El precio límite}

La señalización del cuerpo es un caso del modelo de MILGROM y ROBERTS (1982a). El emisor es J2, acomodaticio con probabilidad $p$ y agresivo con probabilidad $1-p$; el mensaje es el precio que fija antes de la entrada, Alto o Bajo; el receptor es J1, que decide si entra. Tras la entrada, el acomodaticio se acomoda y el agresivo lucha, de modo que entrar le da a J1 $a_1$ frente al primero y $b_1$ frente al segundo, y no entrar, $c_1$. Sean $\pi_\theta^A$ y $\pi_\theta^B$ los beneficios del tipo $\theta$ con el precio previo alto y bajo; $M_\theta$, su beneficio si J1 no entra; y $D_\theta<M_\theta$, el que obtiene si entra. El pago del tipo $\theta$ es su beneficio previo más el posterior.

En el \textbf{equilibrio separador} en el que el agresivo fija el precio bajo y el acomodaticio el alto, las creencias son $\mu(\theta_{ag}\mid\text{Bajo})=1$ y $\mu(\theta_{ag}\mid\text{Alto})=0$: J1 entra tras el precio alto y no entra tras el bajo. El equilibrio exige que ningún tipo quiera imitar al otro:
$$M_{ac}-D_{ac}\leq\pi_{ac}^A-\pi_{ac}^B\qquad\text{y}\qquad\pi_{ag}^B+M_{ag}\geq\pi_{ag}^A+D_{ag}$$
(condiciones de compatibilidad de incentivos del equilibrio separador del ejemplo). La primera es la clave: lo que el acomodaticio ganaría disuadiendo la entrada no compensa lo que le cuesta fijar el precio bajo. En este equilibrio, J1 entra exactamente cuando entraría con información completa, frente al acomodaticio y no frente al agresivo: el precio límite no disuade ninguna entrada que no se hubiera disuadido de todos modos, aunque puede obligar al agresivo a fijar un precio previo más bajo del que fijaría sin necesidad de señalizar.

En el \textbf{equilibrio agrupador} en el que los dos tipos fijan el precio bajo, J1 no aprende nada: $\mu(\theta_{ag}\mid\text{Bajo})=1-p$, y no entra si $1-p>q^*$, el umbral del cuerpo. El precio alto no se observa, y la creencia tras él, $\mu_A$, es libre; para que el acomodaticio no quiera desviarse, J1 ha de entrar tras un precio alto, lo que exige $\mu_A<q^*$, y el acomodaticio ha de preferir el precio bajo sin entrada al alto con entrada, $\pi_{ac}^B+M_{ac}\geq\pi_{ac}^A+D_{ac}$, que es la desigualdad contraria a la del separador. Cuál de los dos equilibrios existe depende, por tanto, de cuánto vale disuadir la entrada frente a lo que cuesta el precio bajo.

\paragraph{El criterio intuitivo}

Como las creencias fuera de la senda son libres, los juegos de señalización suelen tener muchos equilibrios. CHO y KREPS (1987) propusieron descartar las creencias que atribuyen un mensaje no observado a un tipo para el que está \textit{dominado en equilibrio}: el mensaje $m$ está dominado en equilibrio para el tipo $\theta$ si su pago de equilibrio, $u_E^*(\theta)$, supera el máximo que podría obtener con $m$ ante cualquier respuesta óptima del receptor. Un equilibrio supera el \textbf{criterio intuitivo} si, tras eliminar esos tipos de las creencias que siguen a cada mensaje no observado, ningún tipo restante ganaría enviando ese mensaje frente a cualquier respuesta óptima a creencias concentradas en los tipos restantes. Es la formalización más usada de la inducción hacia adelante: el receptor razona sobre qué tipo podría haber querido enviar el mensaje inesperado.

En el ejemplo, sea un equilibrio agrupador en el precio alto, en el que J1 entra ($1-p<q^*$) y la creencia tras un precio bajo, que no se observa, lleva a J1 a entrar también. Si se cumplen, con desigualdad estricta, las dos condiciones del separador, el precio bajo está dominado en equilibrio para el acomodaticio, porque ni disuadiendo la entrada le compensaría su coste; la creencia tras el precio bajo ha de concentrarse en el agresivo, J1 no entraría, y por la segunda condición al agresivo le compensa entonces bajar el precio. El agrupador en el precio alto no supera el criterio, y queda el separador.

\subsubsection{Reputación}

La información incompleta resuelve la paradoja de la cadena de tiendas. KREPS y WILSON (1982b) y MILGROM y ROBERTS (1982b) suponen que, con una probabilidad a priori pequeña, $\varepsilon$, el incumbente es de un tipo \textit{duro} que lucha siempre. Cada vez que lucha, los entrantes siguientes revisan con la regla de Bayes la probabilidad de que sea duro: si $\mu_t$ es la creencia antes del mercado $t$ y el incumbente normal lucha con probabilidad $\beta_t$, la creencia tras observar la lucha es
$$\mu_{t+1}=\frac{\mu_t}{\mu_t+(1-\mu_t)\,\beta_t}$$
(revisión bayesiana de la reputación en el juego de la cadena de tiendas), mientras que acomodarse revela que es normal y lleva la creencia a cero. En el equilibrio, el incumbente normal lucha en los primeros mercados para no revelar su tipo, porque la reputación de duro disuade la entrada en los siguientes, y solo en los últimos aleatoriza; el número de mercados en los que hay entrada no crece con el número total de mercados. Una probabilidad pequeña de un tipo distinto basta para que la amenaza sea creíble durante casi todo el juego. KREPS, MILGROM, ROBERTS y WILSON (1982) obtuvieron el mismo resultado para el dilema del prisionero repetido un número finito de veces: si hay una probabilidad pequeña de que un jugador coopere mientras el otro coopere, la cooperación se sostiene en casi todas las etapas.


\clearpage
\anexo{Juegos cooperativos: Desarrollo formal}\label{anx:cooperativos}

La teoría de juegos cooperativa es la otra gran rama en la que se divide la teoría de juegos. Se diferencia de la \textit{no cooperativa} en que sus unidades fundamentales de análisis son grupos y subgrupos de individuos, las coaliciones, que se supone capaces de alcanzar resultados conjuntamente mediante acuerdos vinculantes, cuyo cumplimiento garantiza una parte externa como un juez. Cabe recordar que el término no cooperativa no significa que esa teoría no sea capaz de explicar la cooperación entre grupos de individuos, sino que se centra en cómo la cooperación puede surgir del comportamiento racional en ausencia de capacidad para realizar acuerdos vinculantes. La cooperativa, en cambio, prescinde de cómo se juega y se pregunta qué repartos de lo que las coaliciones pueden conseguir son estables o justos.

El cuerpo del tema solo presenta esta rama, que se desarrolla al estudiar la colusión y los acuerdos entre empresas ([\authorfont{Ver Tema 3.A.19}]). Este anexo la formaliza de principio a fin: los elementos de un juego cooperativo y sus propiedades; los conceptos de solución basados en la estabilidad, el núcleo y sus variantes, el $\varepsilon$-núcleo y el nucléolo; el basado en la justicia, el valor de SHAPLEY; y, por último, los juegos de utilidad no transferible y la negociación, con la solución de NASH.

\subsection{Conceptos: ¿qué elementos componen un juego cooperativo?}

Los juegos con forma de coalición parten de un supuesto: los jugadores pueden firmar acuerdos vinculantes. Respecto a los pagos, se distinguen dos clases. En los \textit{juegos de utilidad transferible} (TU), las utilidades de los jugadores se pueden sumar, de modo que la coalición obtiene un nivel total que puede redistribuir perfectamente entre sus miembros, por ejemplo con pagos monetarios. En los \textit{juegos de utilidad no transferible} (NTU), los pagos no pueden juntarse en una \textit{cesta conjunta} y redistribuirse después, porque la coalición no obtiene bienes perfectamente divisibles. Los apartados siguientes se centran en los juegos de utilidad transferible, en los que el mejor resultado que una coalición puede garantizarse frente a los demás jugadores recibe el nombre de \textit{valor de la coalición}; el último trata los de utilidad no transferible.

Un juego con forma de coalición se describe con tres elementos. El conjunto de jugadores, $N=\{1,\ldots,n\}$. El conjunto de las partes de $N$, $P(N)=2^N$, que es el conjunto de las coaliciones posibles.\footnote{%
    Denotamos $N$ como la coalición formada por todos los integrantes del Juego: Gran coalición. Además, es importante destacar que una coalición puede integrar cualquier número de jugadores, incluso un único jugador. Es por tanto, la terminología propia de los juegos cooperativos.
} Y la \textbf{función característica}, $v$, que asocia a cada coalición el máximo pago conjunto que puede garantizarse:
$$v:2^N\rightarrow\mathbb{R},\qquad v(\emptyset)=0$$
(función característica de un juego cooperativo con utilidad transferible). El juego se denota $\Gamma=(N,v)$.

\paragraph{El ejemplo}

Usaremos un ejemplo con letras. Tres jugadores pueden repartirse un excedente $V$ si se ponen los tres de acuerdo; dos de ellos solos obtienen $W$, con $0\leq W\leq V$, y uno solo no obtiene nada:
$$v(N)=V,\qquad v(\{i,j\})=W,\qquad v(\{i\})=0\qquad\forall i\neq j$$
(juego simétrico de tres jugadores). Es la generalización del caso clásico en que el excedente es de $100$ euros y una pareja puede conseguir $20$. Lo compararemos con un juego en el que el jugador $1$ es \textit{imprescindible}: cualquier coalición que lo contenga y tenga al menos dos miembros obtiene $V$, y las demás nada, $v(\{1,2\})=v(\{1,3\})=v(N)=V$.

\subsubsection{Propiedades: superaditividad y convexidad}

Un juego $\Gamma=(N,v)$ es \textbf{superaditivo} si
$$v(C\cup D)\geq v(C)+v(D),\qquad\forall C,D\subseteq N\ :\ C\cap D=\emptyset$$
La implicación directa de esta propiedad se puede entender como \textit{más es mejor}: unir dos coaliciones disjuntas nunca les hace perder, de modo que no hay externalidades negativas entre ellas (\textit{estoy mejor sin él} no existe bajo este supuesto). Este supuesto hace que la gran coalición sea, en general, preferible a las coaliciones menores, y por eso suele analizarse cómo se reparte $v(N)$. El juego simétrico del ejemplo es superaditivo si $W\geq0$ y $V\geq W$.

Una propiedad más exigente es la \textbf{convexidad} (SHAPLEY, 1971): el juego es convexo si
$$v(S\cup T)+v(S\cap T)\geq v(S)+v(T)\qquad\forall S,T\subseteq N$$
lo que equivale a que la contribución marginal de cada jugador, $v(S\cup\{i\})-v(S)$, crezca con el tamaño de la coalición a la que se une: hay rendimientos crecientes a la cooperación. En el juego simétrico, las contribuciones marginales son $0$ a la coalición vacía, $W$ a un jugador y $V-W$ a una pareja, de modo que es convexo si y solo si $V-W\geq W$, es decir, $V\geq2W$.

\subsubsection{Conceptos de solución: ¿qué sustituye al equilibrio?}

En un juego cooperativo no se pregunta qué estrategias jugará cada uno, sino qué resultado alcanzarán las coaliciones y cómo se repartirán lo obtenido. Una \textit{solución} de un juego cooperativo es un par $(C,x)$. $C=\{C_1,\ldots,C_k\}$ es la \textit{estructura de coaliciones}, una partición de $N$:
$$\bigcup_{h=1}^{k}C_h=N;\qquad C_h\cap C_l=\emptyset,\ \forall h\neq l$$
y $x=(x_1,x_2,\ldots,x_n)\in\mathbb{R}^n$ es un vector de pagos que cumple dos condiciones: la \textit{racionalidad individual}, porque el pago que cada jugador recibe formando parte de una coalición ha de ser al menos el que obtendría solo; y la \textit{viabilidad} (\textit{feasibility}), porque cada coalición reparte exactamente lo que obtiene:
$$x_i\geq v(\{i\}),\ \forall i\in N\qquad\text{y}\qquad\sum_{i\in C_h}x_i=v(C_h),\ \forall C_h\in C$$
Si se forma la gran coalición, los vectores de pagos que cumplen ambas condiciones se llaman \textbf{imputaciones}:
$$I(v)=\left\{x\in\mathbb{R}^n:\ \sum_{i\in N}x_i=v(N),\ \ x_i\geq v(\{i\})\ \forall i\right\}$$
El lugar del equilibrio lo ocupan dos clases de criterios para elegir entre imputaciones. Los de \textit{estabilidad} buscan los repartos que ninguna coalición quiera romper: es el núcleo y sus variantes. Los de \textit{justicia} buscan un reparto único que responda a lo que cada jugador aporta: es el valor de SHAPLEY.

\subsubsection{Estabilidad}

En los juegos cooperativos, la noción de óptimo se redefine a través del concepto de \textit{estabilidad}.\footnote{%
    Véase: \url{https://youtu.be/mO30o7Q7veE?si=jBBSJP5FnVv3Binm} (Definición de \textit{estabilidad} a través de un ejemplo).
} Un resultado es \textit{estable} si ningún grupo de jugadores tiene incentivos para abandonarlo y formar otra coalición. Formalmente, una coalición $S$ \textit{bloquea} la imputación $x$ si puede garantizarse más de lo que $x$ da a sus miembros:
$$\sum_{i\in S}x_i<v(S)$$
Con superaditividad, el único parámetro con el que se mide la estabilidad de un resultado es el pago que recibe cada jugador comparado con lo que podría obtener con otros: un reparto es estable si ninguna coalición lo bloquea.

\subsection{El núcleo: ¿qué repartos son estables?}

El \textbf{núcleo} es el conjunto de \textit{asignaciones estables}, las que ningún subconjunto de jugadores puede mejorar rompiendo el acuerdo. Formalmente, es el conjunto de imputaciones que ninguna coalición bloquea:\footnote{%
    Se trata de la definición formal de \textit{estabilidad} en juegos cooperativos.
}
$$\text{Núcleo}(\Gamma)=\left\{x\in\mathbb{R}^n\ \middle|\ \sum_{i\in N}x_i=v(N),\ \ \sum_{i\in S}x_i\geq v(S)\ \ \forall S\subseteq N\right\}$$
Es el conjunto de resultados en los que ningún jugador ni ningún grupo tiene incentivos para desviarse, es decir, para dejar la coalición y formar la suya propia. El concepto se parece al de equilibrio de NASH, pero exige más: no basta con que ningún jugador quiera desviarse solo, porque tampoco puede querer hacerlo ningún grupo.

\paragraph{El ejemplo}

En el juego simétrico, el núcleo son los repartos $(x_1,x_2,x_3)$ que cumplen:
$$x_i+x_j\geq W,\ \forall i\neq j;\qquad x_i\geq0,\ \forall i;\qquad x_1+x_2+x_3=V$$
Como $x_i+x_j=V-x_k$, la primera condición equivale a $x_k\leq V-W$: ningún jugador puede llevarse tanto que la pareja de los otros dos prefiera ir por su cuenta. Sumando las tres condiciones de pareja, $2V\geq3W$, de modo que el núcleo no es vacío si y solo si $W\leq\tfrac{2}{3}V$. En el panel (a) de la figura, con $W<V/2$, el núcleo es el hexágono que queda al quitar del símplice de repartos las tres esquinas en las que un jugador se lleva más de $V-W$; con el caso clásico, $V=100$ y $W=20$, ningún jugador puede llevarse más de $80$. En el juego con un jugador imprescindible, en cambio, las condiciones $x_1+x_2\geq V$ y $x_1+x_3\geq V$, con $x_1+x_2+x_3=V$, obligan a $x_2=x_3=0$: el núcleo es un único punto, $(V,0,0)$, el vértice del panel (b), porque cualquier pago a $2$ o a $3$ invitaría a $1$ a formar coalición con el otro.

\begin{center}\begin{minipage}{\linewidth}\centering {\small\textbf{El núcleo y el valor de SHAPLEY en dos juegos de tres jugadores}}\par\vspace{0.3em} \includegraphics{JE_anx_nucleo.pdf}\par\vspace{0.3em} {\small\textit{Fuente: }Elaboración propia.}\end{minipage}\end{center}

\paragraph{Existencia: juegos equilibrados}

El núcleo tiene dos problemas. En algunos casos es el conjunto vacío, de modo que no existe ninguna asignación estable ante las desviaciones de los subgrupos: en el ejemplo, si $W>\tfrac{2}{3}V$. Y solo en escasas ocasiones contiene un único vector. BONDAREVA (1963)\textcolor{orange}{\textsuperscript{[Probable]}} y SHAPLEY (1967) caracterizaron cuándo no es vacío. Una colección de coaliciones con pesos $\lambda_S\geq0$ es \textit{equilibrada} si cada jugador pertenece a coaliciones cuyos pesos suman uno, $\sum_{S\ni i}\lambda_S=1$ para todo $i$; puede entenderse como una manera de repartir el tiempo de cada jugador entre las coaliciones de las que forma parte.

\textbf{Teorema} (BONDAREVA, 1963; SHAPLEY, 1967).- El núcleo de un juego con utilidad transferible no es vacío si y solo si, para toda colección equilibrada de pesos,
$$\sum_{S\subseteq N}\lambda_S\,v(S)\leq v(N)$$
Es decir, si ninguna manera de repartir el tiempo de los jugadores entre coaliciones menores produce más que la gran coalición. En el ejemplo, las tres parejas con peso $\tfrac{1}{2}$ forman una colección equilibrada, y la condición es $\tfrac{3}{2}W\leq V$, la misma que se obtuvo directamente. Los juegos convexos cumplen siempre la condición: su núcleo no es vacío y contiene el valor de SHAPLEY (SHAPLEY, 1971).

\subsubsection{El $\varepsilon$-núcleo y el núcleo mínimo}

Para remediar el primer problema existe una noción \textit{más débil} de estabilidad, el \textbf{$\varepsilon$-núcleo}, que tolera desviaciones pequeñas. Para $\varepsilon>0$:
$$\varepsilon\text{-núcleo}(\Gamma)=\left\{x\ \middle|\ \sum_{i\in N}x_i=v(N),\ \ \sum_{i\in S}x_i\geq v(S)-\varepsilon\ \ \forall S\subsetneq N,\ S\neq\emptyset\right\}$$
Intuitivamente, es el conjunto de resultados en los que ninguna coalición gana más de $\varepsilon$ desviándose, de modo que, si desviarse tiene un coste, por ejemplo de organización, el resultado es estable. Su papel recuerda al de las estrategias mixtas frente al equilibrio de NASH en estrategias puras: amplía el conjunto de candidatos para que haya alguno.\footnote{%
    Recordar, \textit{estrategias puras y/o mixtas}.
} Con $\varepsilon$ suficientemente grande, cualquier imputación forma parte del $\varepsilon$-núcleo, y el objetivo es reducir $\varepsilon$ todo lo posible. Se define
$$\varepsilon^*(\Gamma)=\inf\left\{\varepsilon\in\mathbb{R}\mid\varepsilon\text{-núcleo de }\Gamma\text{ no es vacío}\right\},$$
y el $\varepsilon^*$-núcleo es el \textbf{núcleo mínimo} (MASCHLER, PELEG y SHAPLEY, 1979\textcolor{orange}{\textsuperscript{[Probable]}}).\footnote{%
    No siempre puede formularse como el mínimo del conjunto, porque no sabemos de antemano si se trata de un conjunto abierto o cerrado, y por eso se toma el ínfimo, el número inmediatamente más pequeño a todos los del conjunto, ya sea parte de la frontera o por salto al límite: $(0,1)\rightarrow\{0\}$; $[0,1]\rightarrow\{0\}$.
} Si se admite un $\varepsilon$ negativo, la definición exige que cada coalición obtenga \textit{más} de lo que podría conseguir sola, y el núcleo mínimo selecciona, dentro de un núcleo no vacío, los repartos más alejados de cualquier objeción. En el juego simétrico, el reparto igualitario es el que hace mínima la mayor queja de las parejas, $W-\tfrac{2V}{3}$, y la de los jugadores solos, $-\tfrac{V}{3}$, de modo que $\varepsilon^*=\max\left\{W-\tfrac{2V}{3},-\tfrac{V}{3}\right\}$, que es negativo mientras el núcleo no está vacío ($W<\tfrac{2}{3}V$), y el núcleo mínimo es el punto $(\tfrac{V}{3},\tfrac{V}{3},\tfrac{V}{3})$.

\subsubsection{El nucléolo}

SCHMEIDLER (1969) afinó la idea hasta seleccionar un único reparto. El \textit{exceso} de una coalición $S$ en la imputación $x$ es lo que le falta para alcanzar su valor:
$$e(S,x)=v(S)-\sum_{i\in S}x_i$$
y mide la queja de $S$ frente a $x$; $x$ está en el núcleo si y solo si ningún exceso es positivo. Sea $\vartheta(x)$ el vector de los excesos de todas las coaliciones ordenados de mayor a menor. El \textbf{nucléolo} es la imputación que minimiza $\vartheta(x)$ en el orden lexicográfico: primero hace lo más pequeña posible la mayor queja, después la segunda, y así sucesivamente. SCHMEIDLER (1969) demostró que existe y es único siempre que el conjunto de imputaciones no sea vacío, que pertenece al núcleo cuando este no es vacío y que está dentro del núcleo mínimo. En el juego simétrico, por simetría, es el reparto igualitario; en el del jugador imprescindible, como el núcleo es un único punto, el nucléolo es $(V,0,0)$.

\subsection{El valor de SHAPLEY: ¿qué reparto es justo?}

Un enfoque alternativo al de la estabilidad es el llamado \textbf{valor de SHAPLEY} (SHAPLEY, 1953). En lugar de preguntar qué repartos resisten las objeciones de las coaliciones, parte de unos criterios, llamados axiomas, que debe satisfacer un reparto \textit{justo}, y lo sorprendente es el pequeño número de requisitos que bastan para determinar un reparto único. La pregunta de fondo es qué importancia tiene cada jugador para la cooperación global y qué recompensa puede razonablemente esperar, dado que unos contribuyen más que otros a las coaliciones o tienen más poder de negociación, por ejemplo porque pueden amenazar con destruir el excedente.

\subsubsection{Resultados justos y contribuciones marginales}

La noción de \textit{justicia} en los juegos cooperativos se apoya en la \textbf{contribución marginal} de un jugador $i$ a una coalición $S$ que no lo contiene, el valor que añade al unirse a ella:
$$v(S\cup\{i\})-v(S)$$
Aunque la definición de justicia es muy relativa y varía enormemente en la literatura, cabe distinguir dos maneras de basarla en las contribuciones marginales. La primera es la \textit{contribución marginal en la coalición}: cada jugador debería recibir el valor esperado de lo que aporta a las coaliciones a las que se une. La segunda es la \textit{contribución marginal al juego}: un jugador puede entrar o no en el juego, y cabe valorar el reparto por lo que aporta al juego en su conjunto, forme parte de una coalición o no.\footnote{%
    Este concepto nos permite estudiar los jugadores que \textit{no aportan nada a ninguna coalición}, es decir, aquellos que, independientemente de cómo jueguen, sus contribuciones son nulas y por tanto no deberían recibir nada del resultado.
}

\subsubsection{Axiomas y fórmula}

Un \textit{valor} es una regla $\phi$ que asigna a cada juego $\Gamma=(N,v)$ un vector de pagos $\phi(v)=\left(\phi_1(v),\ldots,\phi_n(v)\right)$. SHAPLEY (1953) le exigió cuatro propiedades:

\begin{lnum}
    \item \textbf{Eficiencia}: se reparte todo lo que obtiene la gran coalición, $\sum_{i\in N}\phi_i(v)=v(N)$.
    \item \textbf{Simetría}: si dos jugadores son intercambiables, porque $v(S\cup\{i\})=v(S\cup\{j\})$ para toda coalición $S$ que no contiene a ninguno de los dos, reciben lo mismo, $\phi_i(v)=\phi_j(v)$.
    \item \textbf{Jugador nulo}: si un jugador no aporta nada a ninguna coalición, $v(S\cup\{i\})=v(S)$ para toda $S$, no recibe nada, $\phi_i(v)=0$. En su versión de \textit{jugador ficticio}, si un jugador aporta siempre exactamente su pago individual, recibe ese pago.
    \item \textbf{Aditividad}: si un mismo grupo de jugadores juega dos juegos, $v$ y $w$, el reparto del juego compuesto, $v+w$, es la suma de los repartos, $\phi_i(v+w)=\phi_i(v)+\phi_i(w)$; la distribución de pagos no depende de cómo se descomponga el juego.
\end{lnum}

\textbf{Teorema} (SHAPLEY, 1953).- Existe un único valor que cumple las cuatro propiedades, y asigna a cada jugador
$$\phi_i(v)=\sum_{S\subseteq N\setminus\{i\}}\frac{|S|!\,(n-|S|-1)!}{n!}\left(v(S\cup\{i\})-v(S)\right)$$
donde $n$ es el número total de jugadores y la suma se extiende sobre todos los subconjuntos de $N$ que no contienen al jugador $i$ (valor de SHAPLEY, 1953). De forma equivalente, sumando sobre las coaliciones $S$ que sí contienen a $i$, con $s=|S|$:
$$\phi_i(v)=\sum_{S\ni i}q(s)\left[v(S)-v(S\setminus\{i\})\right],\qquad q(s)\equiv\frac{(s-1)!\,(n-s)!}{n!}$$
donde $v(S)-v(S\setminus\{i\})$ es la contribución marginal de $i$ a la coalición. Para las coaliciones que no contienen a $i$, $S=S\setminus\{i\}$ y la contribución es nula, de modo que da igual sumar sobre todas las coaliciones o solo sobre las que lo contienen.

La demostración de la unicidad se apoya en los \textit{juegos de unanimidad}: para cada coalición $T$, $u_T(S)=1$ si $T\subseteq S$ y $0$ en otro caso. En $u_T$, los jugadores de fuera de $T$ son nulos y los de dentro, simétricos, de modo que la eficiencia obliga a $\phi_i(u_T)=1/|T|$ para $i\in T$ y $0$ para los demás. Como todo juego puede escribirse de forma única como combinación lineal de juegos de unanimidad, $v=\sum_Tc_T\,u_T$, la aditividad, extendida a múltiplos, determina $\phi(v)$, y el resultado es la fórmula anterior.

\paragraph{Interpretación y ejemplo}

La fórmula admite una lectura muy intuitiva. Imaginemos que la gran coalición se forma incorporando a los jugadores de uno en uno, en un orden elegido al azar entre los $n!$ posibles, y que cada uno recibe su contribución marginal en el momento de unirse. El valor de SHAPLEY de $i$ es la media de esas contribuciones sobre todos los órdenes:
$$\phi_i(v)=\frac{1}{n!}\sum_{\pi}\left[v\left(P_i^\pi\cup\{i\}\right)-v\left(P_i^\pi\right)\right]$$
donde $P_i^\pi$ es el conjunto de jugadores que preceden a $i$ en el orden $\pi$, y los pesos $q(s)$ son las probabilidades de que $i$ se encuentre ya formada la coalición $S\setminus\{i\}$. Es, por tanto, el \textit{valor esperado de la contribución de cada jugador a las coaliciones del juego}.

En el juego simétrico, cada jugador aporta $0$ si llega primero, $W$ si llega segundo y $V-W$ si llega tercero, y como cada posición es igual de probable, $\phi_i=\tfrac{1}{3}(0+W+V-W)=\tfrac{V}{3}$, lo mismo que exigen la simetría y la eficiencia. En el juego del jugador imprescindible, la tabla recoge las contribuciones en los seis órdenes posibles:
$$\begin{array}{c|ccc}
\text{Orden} & \text{J1} & \text{J2} & \text{J3}\\\hline
1,2,3 & 0 & V & 0\\
1,3,2 & 0 & 0 & V\\
2,1,3 & V & 0 & 0\\
2,3,1 & V & 0 & 0\\
3,1,2 & V & 0 & 0\\
3,2,1 & V & 0 & 0\\\hline
\phi_i & \tfrac{2V}{3} & \tfrac{V}{6} & \tfrac{V}{6}
\end{array}$$
El jugador imprescindible recibe $\tfrac{2V}{3}$, y los otros dos, $\tfrac{V}{6}$ cada uno. El ejemplo muestra la diferencia entre los dos criterios: el núcleo da todo a $1$, porque cualquier pago a los demás sería bloqueado, mientras que el valor de SHAPLEY reconoce a $2$ y a $3$ lo que aportan cuando llegan después de $1$. El valor de SHAPLEY no tiene por qué estar en el núcleo, salvo en los juegos convexos.

\subsection{Juegos de utilidad no transferible y negociación}

Cuando la utilidad no es transferible, el valor de una coalición ya no es un número, sino un conjunto de vectores de pagos que puede garantizarse, y los conceptos anteriores se generalizan con dificultad. El caso más importante es el de dos jugadores, el \textbf{problema de negociación}: un par $(U,d)$, donde $U\subset\mathbb{R}^2$ es el conjunto de pagos alcanzables por acuerdo, convexo y compacto, y $d\in U$ es el \textit{punto de desacuerdo}, lo que cada uno obtiene si no hay acuerdo; se supone que algún punto de $U$ mejora $d$ para los dos.

\textbf{Teorema} (NASH, 1950c).- Existe una única regla de negociación $f(U,d)$ que cumple cuatro axiomas: (i) \textit{eficiencia de PARETO}, porque el acuerdo no puede mejorarse para los dos a la vez; (ii) \textit{simetría}, porque si el problema es simétrico, los dos obtienen lo mismo; (iii) \textit{invariancia a las transformaciones afines} de la utilidad de cada jugador, que, como vimos, no alteran sus preferencias; y (iv) \textit{independencia de las alternativas irrelevantes}, porque si se reduce $U$ sin eliminar el acuerdo elegido, el acuerdo no cambia. Esa regla elige el punto que maximiza el \textit{producto de NASH}:
$$u^N=\arg\max_{u\in U,\ u\geq d}\ (u_1-d_1)(u_2-d_2)$$
(solución de negociación de NASH, 1950c).

\begin{center}\begin{minipage}{\linewidth}\centering {\small\textbf{La solución de negociación de NASH}}\par\vspace{0.3em} \includegraphics{JE_anx_negociacion.pdf}\par\vspace{0.3em} {\small\textit{Fuente: }Elaboración propia a partir de NASH (1950c).}\end{minipage}\end{center}

En el panel (a) de la figura, la solución es el punto de la frontera de $U$ en el que la curva de nivel más alta del producto de NASH es tangente a ella. En el panel (b), si dos jugadores reparten un excedente $V$ con utilidad transferible, la frontera es $u_1+u_2=V$ y la solución da a cada uno su pago de desacuerdo más la mitad del excedente neto:
$$u_i^N=d_i+\tfrac{1}{2}\left(V-d_1-d_2\right)$$
Lo que cada uno puede conseguir sin acuerdo, su \textit{opción externa}, determina lo que obtiene con él. La versión asimétrica, que maximiza $(u_1-d_1)^{\alpha}(u_2-d_2)^{1-\alpha}$, recoge distintos poderes de negociación.

La independencia de las alternativas irrelevantes es el axioma más discutido. KALAI y SMORODINSKY (1975) la sustituyeron por una condición de monotonía, que pide que un jugador no salga perdiendo si las posibilidades de acuerdo mejoran a su favor, y obtuvieron otra solución, que reparte las ganancias en proporción a lo máximo que cada uno podría aspirar a obtener. El llamado \textit{programa de NASH} busca, en cambio, fundamentar las soluciones cooperativas en juegos no cooperativos. RUBINSTEIN (1982) demostró que el juego de ofertas alternativas con descuento tiene un único ENPS, y BINMORE, RUBINSTEIN y WOLINSKY (1986) mostraron que, cuando el tiempo entre ofertas tiende a cero, ese equilibrio converge a la solución de NASH: la solución cooperativa es el resultado de una negociación no cooperativa entre jugadores pacientes.


\clearpage
\anexo{Diseño de Mecanismos: Desarrollo formal}\label{anx:mecanismos: teoria}

El diseño de mecanismos invierte la pregunta de la teoría de juegos. En lugar de tomar el juego como dado y predecir cómo se jugará, parte del resultado que se quiere alcanzar, una regla que diga qué decisión tomar según la información privada de los participantes, y busca un juego cuyas reglas lleven a ese resultado en equilibrio, sea cual sea esa información. La dificultad es que el diseñador no la conoce, y los participantes, que sí la conocen, la revelarán solo si les conviene. Este anexo desarrolla formalmente la teoría de principio a fin: el entorno, las funciones de elección social, los mecanismos y la noción de implementación; el principio de revelación, que reduce el problema a buscar reglas compatibles con los incentivos; la implementación en estrategias dominantes, con el teorema de imposibilidad de GIBBARD y SATTERTHWAITE y los mecanismos de VICKREY, CLARKE y GROVES; la implementación bayesiana, con la caracterización de los mecanismos compatibles con los incentivos y el teorema de MYERSON y SATTERTHWAITE; y las extensiones, limitaciones y aplicaciones de la teoría. La última parte aplica todo ello a las subastas.\footnote{%
  Se sigue en lo esencial la notación de MAS-COLELL, WHINSTON y GREEN (1995), capítulo 23, con la de decisiones y transferencias del cuerpo del tema.
}

\subsection{El entorno: decisiones, tipos y preferencias}

La siguiente descripción de elementos, el \textbf{marco general del diseño de mecanismos} (\textit{general mechanism design setting}), puede aplicarse a cualquier problema de diseño de mecanismos:

\begin{lnum}
    \item \textbf{Número finito de jugadores}:
    $$N=\{1,2,\ldots,n\};\qquad i,j,k\in N$$
    \item \textbf{Decisiones}. El conjunto de todas las decisiones sociales posibles, finito o infinito:
    $$D,\qquad d,d'\in D$$
    \item \textbf{Preferencias e información}. Los \textit{inputs} del mecanismo son siempre información privada. La información privada del individuo $i$ se representa por su tipo, como en los juegos bayesianos:
    $$\theta_i\in\Theta_i,$$
    donde $\theta_i$ es un tipo posible del individuo $i$ y $\Theta_i$ el conjunto de todos sus tipos posibles. El perfil de tipos es
    $$\theta=(\theta_1,\theta_2,\ldots,\theta_n)\in\Theta_1\times\Theta_2\times\cdots\times\Theta_n=\Theta,$$
    un vector de $n$ dimensiones, y $\Theta$ es el producto cartesiano de los espacios de tipos de los $n$ individuos. Los tipos se distribuyen según una probabilidad a priori común, $F$, que el diseñador también conoce.
    \item \textbf{Una función de utilidad dependiente del tipo}:
    $$v_i:D\times\Theta_i\rightarrow\mathbb{R},$$
    que depende de la decisión y del tipo del jugador $i$ y asigna un valor real a cada combinación. Así, $v_i(d,\theta_i)>v_i(d',\theta_i)$ significa que el tipo $\theta_i$ prefiere la decisión $d$ a la $d'$. Se supone que $v_i$ no depende del tipo de ningún otro jugador: es el caso de \textit{valores privados}.
    \item \textbf{Reglas de decisión}. Una regla de decisión es una función que asigna una decisión a cada perfil de tipos:
    $$d:\Theta\rightarrow D,\qquad d(\theta)\in D.$$
    Una regla de decisión es \textbf{eficiente} si y solo si maximiza la suma de las utilidades de todos los jugadores frente a cualquier otra decisión posible:
    $$\sum_{i=1}^{n}v_i\left(d(\theta),\theta_i\right)\geq\sum_{i=1}^{n}v_i\left(d',\theta_i\right);\qquad\forall\theta\in\Theta\ \text{y}\ \forall d'\in D$$
\end{lnum}

\textbf{NOTA.-} Con utilidad transferible, una regla eficiente es el equivalente del óptimo de PARETO: si una decisión no maximizase la suma de utilidades, cambiar a la que lo hace y compensar con dinero a los perjudicados mejoraría a todos.

En muchos entornos, como el de las subastas, las utilidades son transferibles, y hay que añadir un sexto componente:

\begin{lnum}\setcounter{numitem}{5}
    \item \textbf{Una función de transferencias}:
    $$t:\Theta\rightarrow\mathbb{R}^n,$$
    donde $t_i(\theta)$ representa el pago que el jugador $i$ recibe, si es positivo, o realiza, si es negativo, en función del perfil de tipos $\theta$.
\end{lnum}

En una subasta, por ejemplo, la decisión es quién se queda el objeto, y la función de transferencias recoge el pago del ganador al vendedor. Con transferencias, la utilidad del jugador $i$ es \textbf{cuasilineal}:
$$u_i(d,t_i,\theta_i)=v_i(d,\theta_i)+t_i$$
(utilidad cuasilineal del diseño de mecanismos, MAS-COLELL, WHINSTON y GREEN, 1995). La función de transferencias es \textbf{viable} (\textit{feasible}) si no necesita financiación externa,
$$\sum_{i=1}^{n}t_i(\theta)\leq0\qquad\forall\theta,$$
y es de \textbf{presupuesto equilibrado} si la condición se cumple con igualdad. En este último caso, el mecanismo, marco normativo o \textit{institución} es autosuficiente: no necesita que se introduzca renta para alcanzar sus objetivos ni la desperdicia. Un par $(d(\cdot),t(\cdot))$ es \textit{eficiente ex post} si la regla de decisión es eficiente y el presupuesto está equilibrado: con utilidad cuasilineal, son las dos condiciones del óptimo de PARETO.

\paragraph{Tres entornos}

Tres entornos sirven de ilustración a lo largo del anexo. En un \textbf{proyecto público}, $D=\{0,1\}$ (no hacerlo o hacerlo), y la utilidad de cada jugador es $v_i(d,\theta_i)=d\,\theta_i$, donde $\theta_i$ es su valoración neta de su parte del coste; es eficiente hacerlo si y solo si $\sum_i\theta_i\geq0$. En una \textbf{subasta} de un objeto indivisible, $D=\{1,\ldots,n\}$ indica quién se lo queda, $v_i(d,\theta_i)=\theta_i$ si $d=i$ y $0$ en otro caso, y la regla eficiente se lo da a quien más lo valora. En el \textbf{comercio bilateral}, un vendedor con coste $\theta_v$ y un comprador con valoración $\theta_c$ deciden si intercambian, $d\in\{0,1\}$, con $v_c=d\,\theta_c$ y $v_v=-d\,\theta_v$, y es eficiente intercambiar si y solo si $\theta_c\geq\theta_v$.

\subsection{Funciones de elección social, mecanismos e implementación}

El objetivo del diseñador se describe con una \textbf{función de elección social}, que asigna un resultado a cada perfil de tipos:
$$f:\Theta\rightarrow X,\qquad f(\theta)=\left(d(\theta),t_1(\theta),\ldots,t_n(\theta)\right)$$
donde $X=D\times\mathbb{R}^n$ es el conjunto de resultados. Si el diseñador conociera $\theta$, le bastaría con aplicar $f$; como no lo conoce, ha de diseñar un juego.

\subsubsection{Definición de Mecanismo}

Un \textbf{mecanismo} es una pareja $(\mathcal{M},g)$, donde
$$\mathcal{M}=\mathcal{M}_1\times\mathcal{M}_2\times\cdots\times\mathcal{M}_n$$
es el espacio de estrategias, o de mensajes, producto cartesiano de los de cada jugador, y
$$g:\mathcal{M}\rightarrow D\times\mathbb{R}^n$$
es una función de resultados, o regla de asignación (definición de mecanismo, MAS-COLELL, WHINSTON y GREEN, 1995). Para cada perfil de mensajes $m=(m_1,m_2,\ldots,m_n)\in\mathcal{M}$,
$$g(m)=\left(g_d(m),\ g_{t,1}(m),\ldots,\ g_{t,n}(m)\right)$$
representa la decisión y las transferencias resultantes. Dado el perfil de mensajes, la función de resultados da la decisión social, su primer componente, y el vector de transferencias de cada individuo, de $1$ a $n$. El recorrido de $g$ es el mismo conjunto de resultados que el de la función de elección social: por cada perfil de mensajes declarado, el mecanismo escoge un resultado.

Por ejemplo, en el entorno de subastas tendríamos:

\begin{lnum}
    \item Un mecanismo cuyo conjunto $\mathcal{M}$ estaría formado por el conjunto de vectores de \textit{pujas} posibles. El mecanismo no recibe como \textit{inputs} las valoraciones reales de los jugadores, sino sus pujas, $m$;
    \item Dadas las \textit{pujas}, el mecanismo determina quién gana y quién pierde, por ejemplo el que puja más alto, y el resto de componentes de $g$ es el vector de transferencias: uno negativo, el pago del ganador, y todos los demás cero, o un vector de balance nulo si se incluye al vendedor.
\end{lnum}

\subsubsection{Implementación}

Un mecanismo, junto con los tipos, su distribución y las preferencias, define un juego bayesiano: cada jugador observa su tipo y elige un mensaje, y su estrategia es una regla $s_i:\Theta_i\rightarrow\mathcal{M}_i$. El mecanismo \textbf{implementa} la función de elección social $f$ si el juego que induce tiene un equilibrio $s^*(\cdot)=\left(s_1^*(\cdot),\ldots,s_n^*(\cdot)\right)$ cuyo resultado coincide con el que $f$ prescribe para cada perfil de tipos:
$$g\left(s_1^*(\theta_1),\ldots,s_n^*(\theta_n)\right)=f(\theta)\qquad\forall\theta\in\Theta$$
(implementación de una función de elección social, MAS-COLELL, WHINSTON y GREEN, 1995). La implementación es \textbf{en estrategias dominantes} si cada $s_i^*(\theta_i)$ es óptima para el tipo $\theta_i$ frente a cualquier mensaje de los demás:
$$u_i\left(g\left(s_i^*(\theta_i),m_{-i}\right),\theta_i\right)\geq u_i\left(g\left(m_i,m_{-i}\right),\theta_i\right)\qquad\forall m_i\in\mathcal{M}_i,\ \forall m_{-i}\in\mathcal{M}_{-i},\ \forall\theta_i\in\Theta_i,\ \forall i$$
y \textbf{bayesiana} si $s^*$ es un equilibrio bayesiano de NASH, de modo que cada tipo solo necesita que su mensaje sea óptimo en esperanza, dadas las estrategias de los demás y su creencia sobre sus tipos. La primera es más exigente, pero más robusta: no requiere que los jugadores conozcan la distribución de los tipos ajenos ni que acierten las estrategias de los demás.\footnote{%
  La definición pide que \textit{exista} un equilibrio que lleve a $f$, pero el juego puede tener otros que no lo hagan. Exigir que todos los equilibrios lleven a $f$ es la \textit{implementación completa}, que se trata al final del anexo.
}

Un mecanismo es \textbf{directo} si cada jugador solo puede anunciar su tipo, $\mathcal{M}_i=\Theta_i$, y el resultado es el que la función de elección social asigna a los tipos anunciados, $g(\hat{\theta})=f(\hat{\theta})$. La función de elección social puede entenderse, por tanto, como un mecanismo en sí misma, que pregunta a cada uno por su tipo, \textit{¿qué es lo que más te gusta?}, y aplica la regla a las respuestas. Un mecanismo más general admite como \textit{inputs} cualquier tipo de mensaje y recupera los tipos a través de las estrategias de equilibrio. $f$ es \textbf{implementable con veracidad} (\textit{truthfully implementable}), o \textbf{compatible con los incentivos}, si en el mecanismo directo decir la verdad, $\hat{\theta}_i=\theta_i$, es un equilibrio. En estrategias dominantes, la condición es
$$u_i\left(f(\theta_i,\hat{\theta}_{-i}),\theta_i\right)\geq u_i\left(f(\hat{\theta}_i,\hat{\theta}_{-i}),\theta_i\right)\qquad\forall\theta_i,\hat{\theta}_i\in\Theta_i,\ \forall\hat{\theta}_{-i}\in\Theta_{-i},\ \forall i$$
y en equilibrio bayesiano,
$$E_{\theta_{-i}}\!\left[u_i\left(f(\theta_i,\theta_{-i}),\theta_i\right)\mid\theta_i\right]\geq E_{\theta_{-i}}\!\left[u_i\left(f(\hat{\theta}_i,\theta_{-i}),\theta_i\right)\mid\theta_i\right]\qquad\forall\theta_i,\hat{\theta}_i\in\Theta_i,\ \forall i$$
(compatibilidad con los incentivos en estrategias dominantes y bayesiana, MAS-COLELL, WHINSTON y GREEN, 1995).

\subsection{El principio de revelación}

\textbf{Teorema} (principio de revelación; GIBBARD, 1973, para estrategias dominantes, y MYERSON, 1979, para el equilibrio bayesiano).- Si una función de elección social $f$ puede implementarse en estrategias dominantes (en equilibrio bayesiano de NASH) mediante algún mecanismo, entonces es implementable con veracidad en estrategias dominantes (en equilibrio bayesiano de NASH) mediante el mecanismo directo.

La demostración es sencilla. Sea $(\mathcal{M},g)$ un mecanismo que implementa $f$ con el equilibrio $s^*$, y sea el mecanismo directo que, ante los anuncios $\hat{\theta}$, aplica el resultado que el original habría producido si cada jugador hubiera jugado su estrategia de equilibrio para el tipo anunciado, $g\left(s^*(\hat{\theta})\right)=f(\hat{\theta})$. Si a un jugador de tipo $\theta_i$ le conviniera anunciar $\hat{\theta}_i\neq\theta_i$, en el mecanismo original le habría convenido enviar $s_i^*(\hat{\theta}_i)$ en lugar de $s_i^*(\theta_i)$, lo que contradice que $s^*$ sea un equilibrio. El mecanismo directo hace por el jugador la cuenta que este haría en el original: miente en su lugar, si conviene, de modo que al jugador no le queda nada que ganar mintiendo.

\begin{center}\begin{minipage}{\linewidth}\centering {\small\textbf{Implementación y principio de revelación}}\par\vspace{0.3em} \includegraphics{JE_anx_mountreiter.pdf}\par\vspace{0.3em} {\small\textit{Fuente: }Elaboración propia a partir de MAS-COLELL, WHINSTON y GREEN (1995).}\end{minipage}\end{center}

El principio tiene una consecuencia metodológica decisiva: para saber qué resultados son alcanzables, basta con estudiar las funciones de elección social compatibles con los incentivos, sin explorar el conjunto infinito de mecanismos posibles. Y una consecuencia económica: la información privada tiene un coste, porque para que cada uno diga la verdad hay que dejarle una \textit{renta informativa}, y no todas las funciones de elección social la permiten.

\subsection{Implementación en estrategias dominantes}

\subsubsection{El teorema de GIBBARD y SATTERTHWAITE}

Sin restricciones sobre las preferencias, la implementación en estrategias dominantes es casi imposible. Sea $X$ finito y supongamos que cada jugador puede tener cualquier ordenación estricta de los resultados. Una función de elección social es \textit{dictatorial} si existe un jugador $i$ cuyo resultado preferido se elige siempre.

\textbf{Teorema} (GIBBARD, 1973; SATTERTHWAITE, 1975).- Si $f$ puede tomar al menos tres valores distintos y el dominio de preferencias es universal, $f$ es implementable con veracidad en estrategias dominantes si y solo si es dictatorial.

Es el equivalente, para los incentivos, del teorema de imposibilidad de ARROW: no hay forma no dictatorial de agregar preferencias arbitrarias sobre tres o más alternativas sin que a alguien le convenga mentir. Las salidas consisten en restringir el dominio de preferencias. Si las preferencias son \textit{unimodales} sobre una dimensión, elegir la alternativa preferida por el votante mediano es compatible con los incentivos en estrategias dominantes. Si la utilidad es cuasilineal y caben transferencias, el dinero permite alinear los incentivos, y es el caso que sigue.

\subsubsection{Entorno cuasilineal: mecanismos de VICKREY, CLARKE y GROVES}

Sea $d^*(\cdot)$ una regla de decisión eficiente. Un \textbf{mecanismo de GROVES} acompaña la decisión eficiente con la transferencia
$$t_i(\hat{\theta})=\sum_{j\neq i}v_j\left(d^*(\hat{\theta}),\hat{\theta}_j\right)+h_i(\hat{\theta}_{-i})$$
(mecanismo de GROVES, 1973), donde $h_i$ es una función cualquiera de los anuncios de los demás. Cada jugador recibe la suma de las utilidades declaradas por los demás en la decisión elegida, más un término que no depende de lo que él anuncie.

\textbf{Teorema} (GROVES, 1973).- En un mecanismo de GROVES, decir la verdad es una estrategia dominante para todos los jugadores.

En efecto, con un anuncio $\hat{\theta}_i$ cualquiera y los anuncios de los demás fijos, el pago del jugador $i$ es
$$v_i\left(d^*(\hat{\theta}_i,\hat{\theta}_{-i}),\theta_i\right)+\sum_{j\neq i}v_j\left(d^*(\hat{\theta}_i,\hat{\theta}_{-i}),\hat{\theta}_j\right)+h_i(\hat{\theta}_{-i})$$
El último término no depende de su anuncio, y la suma de los dos primeros es el bienestar total de la decisión, medido con su tipo verdadero y los anuncios ajenos. Como $d^*$ elige la decisión que maximiza esa suma cuando se le da el tipo verdadero, anunciarlo es lo mejor que puede hacer, digan lo que digan los demás. El mecanismo convierte a cada jugador en propietario residual del bienestar social, de modo que su interés coincide con el colectivo.

El \textbf{mecanismo pivotal} de CLARKE (1971) elige $h_i$ para que cada jugador pague la externalidad que su presencia impone a los demás:
$$t_i(\hat{\theta})=\sum_{j\neq i}v_j\left(d^*(\hat{\theta}),\hat{\theta}_j\right)-\sum_{j\neq i}v_j\left(d^*_{-i}(\hat{\theta}_{-i}),\hat{\theta}_j\right)$$
(mecanismo pivotal de CLARKE, 1971), donde $d^*_{-i}$ es la decisión que sería eficiente si $i$ no participara. La transferencia es siempre no positiva, porque los demás estaban mejor, o igual, con la decisión que maximiza su propio bienestar, y solo es distinta de cero si el jugador es \textit{pivotal}, es decir, si su presencia cambia la decisión. La función de transferencias es viable. En la subasta de un objeto, el mecanismo pivotal es la subasta de segundo precio de VICKREY (1961): el ganador paga la mayor valoración de los demás, que es lo que estos pierden porque él se lleva el objeto, y los demás no pagan nada. En el proyecto público, el jugador pivotal paga lo que su voto cuesta a los demás, $\left|\sum_{j\neq i}\hat{\theta}_j\right|$, y el resto nada.

Los mecanismos de VICKREY, CLARKE y GROVES (VCG) resuelven la eficiencia, pero no el presupuesto. GREEN y LAFFONT (1977) demostraron que, si el dominio de valoraciones es suficientemente rico, los mecanismos de GROVES son los únicos que implementan con veracidad en estrategias dominantes una regla de decisión eficiente, y que ninguno de ellos equilibra el presupuesto para todos los perfiles de tipos: no existe ningún mecanismo eficiente, compatible con los incentivos en estrategias dominantes y de presupuesto equilibrado. Con el mecanismo de CLARKE, los pagos de los jugadores pivotales tienen que destruirse o entregarse a alguien de fuera, porque devolverlos alteraría los incentivos.

\subsection{Implementación bayesiana}

Si se renuncia a las estrategias dominantes y se exige solo que decir la verdad sea un equilibrio bayesiano, el presupuesto puede equilibrarse. Se supone en este apartado que los tipos son independientes.

\subsubsection{El mecanismo de la externalidad esperada}

El mecanismo de d'ASPREMONT y GÉRARD-VARET (1979) y ARROW (1979)\textcolor{orange}{\textsuperscript{[Probable]}} sustituye la utilidad de los demás por su valor esperado:
$$t_i(\hat{\theta})=\xi_i(\hat{\theta}_i)+h_i(\hat{\theta}_{-i}),\qquad\xi_i(\hat{\theta}_i)=E_{\theta_{-i}}\!\left[\sum_{j\neq i}v_j\left(d^*(\hat{\theta}_i,\theta_{-i}),\theta_j\right)\right]$$
(mecanismo de la externalidad esperada, d'ASPREMONT y GÉRARD-VARET, 1979). Si los demás dicen la verdad, el pago esperado del jugador $i$ es el bienestar total esperado más un término que no depende de su anuncio, de modo que decir la verdad es un equilibrio bayesiano por el mismo argumento que en el mecanismo de GROVES. Como $\xi_i$ solo depende del anuncio de $i$, puede elegirse $h_i(\hat{\theta}_{-i})=-\frac{1}{n-1}\sum_{j\neq i}\xi_j(\hat{\theta}_j)$, y entonces la suma de las transferencias es nula para todo perfil de anuncios: cada jugador financia, a partes iguales con los demás, lo que el mecanismo da a los otros. El mecanismo es eficiente, compatible con los incentivos en equilibrio bayesiano y de presupuesto equilibrado.

\subsubsection{Participación voluntaria}

Falta una condición que los mecanismos anteriores no garantizan: que los jugadores quieran participar. Sea $\bar{u}_i(\theta_i)$ lo que el tipo $\theta_i$ obtiene si no participa. La \textbf{restricción de participación} o de racionalidad individual \textit{interina} exige que, conociendo su tipo pero no los ajenos, cada uno espere obtener al menos eso:
$$U_i(\theta_i)=E_{\theta_{-i}}\!\left[u_i\left(f(\theta_i,\theta_{-i}),\theta_i\right)\mid\theta_i\right]\geq\bar{u}_i(\theta_i)\qquad\forall\theta_i,\ \forall i$$
Hay versiones más fuertes, \textit{ex post}, para cada perfil de tipos, y más débiles, \textit{ex ante}, antes de conocer el propio tipo.

\subsubsection{La caracterización de los mecanismos compatibles con los incentivos}

Con tipos unidimensionales, la compatibilidad con los incentivos tiene una caracterización muy útil. Sea $\theta_i\in[\underline{\theta},\bar{\theta}]$ y sea $v_i(d,\theta_i)=\theta_i\,q_i(d)$, donde $q_i(d)$ es, por ejemplo, la probabilidad de que $i$ reciba el objeto. En el mecanismo directo, sean $Q_i(\hat{\theta}_i)=E_{\theta_{-i}}\left[q_i\left(d(\hat{\theta}_i,\theta_{-i})\right)\right]$ la asignación interina que espera recibir quien anuncia $\hat{\theta}_i$, $T_i(\hat{\theta}_i)$ la transferencia interina y $U_i(\theta_i)=\theta_i\,Q_i(\theta_i)+T_i(\theta_i)$ la utilidad interina de decir la verdad.

\textbf{Lema} (MYERSON, 1981).- El mecanismo directo es compatible con los incentivos en equilibrio bayesiano si y solo si, para todo $i$, (i) $Q_i(\cdot)$ es no decreciente y (ii)
$$U_i(\theta_i)=U_i(\underline{\theta})+\int_{\underline{\theta}}^{\theta_i}Q_i(s)\,ds\qquad\forall\theta_i\in[\underline{\theta},\bar{\theta}]$$
(condición de la envolvente de MYERSON, 1981).

La compatibilidad con los incentivos exige que $U_i(\theta_i)=\max_{\hat{\theta}_i}\left[\theta_i\,Q_i(\hat{\theta}_i)+T_i(\hat{\theta}_i)\right]$. Es el máximo de funciones lineales en $\theta_i$, y por tanto convexa, con pendiente $Q_i(\theta_i)$ por el teorema de la envolvente; la convexidad exige que esa pendiente no decrezca, y la integral de la pendiente da (ii). La condición tiene dos consecuencias fundamentales. La primera: la renta informativa de cada tipo, lo que obtiene por encima del tipo más bajo, está determinada por la regla de asignación, y crece con el tipo. La segunda: dos mecanismos con la misma regla de asignación y el mismo $U_i(\underline{\theta})$ dan a cada tipo la misma utilidad esperada y, por tanto, el mismo pago esperado; es el origen del teorema de equivalencia de ingresos de las subastas.

\subsubsection{El teorema de MYERSON y SATTERTHWAITE}

Las rentas informativas tienen un coste que, en el comercio bilateral, no puede cubrirse. Sea un vendedor cuyo coste $\theta_v$ se distribuye en $[\underline{c},\bar{c}]$ y un comprador cuya valoración $\theta_c$ se distribuye en $[\underline{w},\bar{w}]$, con densidades positivas e independientes y soportes que se solapan, de modo que unas veces conviene intercambiar y otras no.

\textbf{Teorema} (MYERSON y SATTERTHWAITE, 1983).- No existe ningún mecanismo eficiente \textit{ex post}, compatible con los incentivos en equilibrio bayesiano, que cumpla la restricción de participación interina y no necesite financiación externa.

La intuición se sigue del lema. Para que cada parte diga la verdad cuando intercambiar es eficiente, hay que dejarle una renta informativa: al vendedor, por no exagerar su coste, y al comprador, por no rebajar su valoración. Cuando los soportes se solapan, la suma de las dos rentas supera el excedente esperado del intercambio, y la diferencia solo puede cubrirla alguien de fuera. Sin un tercero que subvencione, la negociación entre partes con información privada es necesariamente ineficiente: unas veces no se intercambia aunque convendría. Es el reverso del teorema de COASE cuando la información es incompleta.

\subsection{Extensiones, limitaciones y aplicaciones}

\paragraph{Implementación completa}

La implementación definida antes es \textit{débil}: pide que exista un equilibrio que lleve a $f$, pero el mecanismo puede tener otros que no lo hagan. MASKIN (1999) estudió la \textbf{implementación completa en equilibrio de NASH}, en la que todos los equilibrios llevan a $f$ cuando los jugadores conocen los tipos de los demás, aunque el diseñador no. Demostró que una condición de \textit{monotonía} es necesaria y que, con tres o más jugadores y sin que ninguno tenga poder de veto, también es suficiente. MOORE y REPULLO (1988)\textcolor{orange}{\textsuperscript{[Probable]}} mostraron que, con mecanismos secuenciales y perfección en subjuegos, casi cualquier función de elección social es implementable.

\paragraph{Limitaciones}

La teoría descansa en supuestos exigentes. El diseñador ha de conocer la distribución de los tipos y los jugadores han de compartirla, algo que la llamada \textit{doctrina de WILSON} critica por exigir mecanismos que dependen de detalles que nadie conoce; el \textit{diseño robusto de mecanismos} busca reglas que funcionen con poca información sobre las creencias (BERGEMANN y MORRIS, 2005\textcolor{orange}{\textsuperscript{[Probable]}}). Con tipos correlacionados, la teoría predice que el diseñador puede extraer todas las rentas informativas con loterías apropiadas (CRÉMER y MCLEAN, 1988\textcolor{orange}{\textsuperscript{[Probable]}}), un resultado poco verosímil que revela la fragilidad de los mecanismos óptimos ante sus supuestos. Además, los mecanismos pueden ser complejos de entender, vulnerables a la colusión entre participantes y exigir que el diseñador se comprometa a no renegociar las reglas una vez revelada la información.

\paragraph{Aplicaciones}

Las aplicaciones abarcan buena parte de la economía. Las subastas, que se tratan a continuación, son la más desarrollada. En la regulación de un monopolio con costes desconocidos, BARON y MYERSON (1982) obtuvieron el contrato regulador óptimo con la misma técnica de las rentas informativas. La provisión de bienes públicos, la redacción de contratos con información privada ([\authorfont{Ver Tema 3.A.13}]), el diseño de sistemas de votación y el diseño de mercados de emparejamiento, como la asignación de plazas escolares o de órganos para trasplantes a partir del algoritmo de GALE y SHAPLEY (1962), son otras tantas aplicaciones. HURWICZ, MASKIN y MYERSON recibieron el premio Nobel en 2007 por haber fundado la teoría.



\end{document}
